\documentclass[12pt]{article}
\title{Digest of The New York Times}
\author{The New York Times}

% \definecolor{dark-green}{rgb}{0,0.25,0}
% \usepackage[utf8]{inputenc}

\usepackage{fontspec,xunicode,xltxtra}
\setmainfont[Ligatures=Required,Mapping=tex-text]{Georgia}
% \setmainfont[Ligatures=Required,Mapping=tex-text]{Palatino Linotype}
\setsansfont[Ligatures=Required,Mapping=tex-text]{Myriad Pro}
\setmonofont[Ligatures=Required,Mapping=tex-text]{Courier New}
\usepackage{xeCJK}
\setCJKmainfont{宋体}
\setCJKmonofont{黑体}

% \usepackage{indentfirst}


\usepackage[dvips,dvipsnames,svgnames]{xcolor}
\definecolor{light-gray}{gray}{0.95}

%生成PDF的链接
\usepackage[dvipdfmx]{hyperref}
\hypersetup{
    bookmarks=true,         % show bookmarks bar?
    bookmarksopen=true,
    pdfpagemode=UseNone,    % options: UseNode, UseThumbs, UseOutlines, FullScreen
    pdfstartview=FitH,
    pdfborder=1,
    pdfhighlight=/P,
    pdfauthor={wuxch},
    unicode=true,           % non-Latin characters in Acrobat’s bookmarks
    colorlinks,             % false: boxed links; true: colored links
    linkcolor=blue,         % color of internal links
    citecolor=blue,        % color of links to bibliography
    filecolor=magenta,      % color of file links
    urlcolor=cyan           % color of external links
}
% 使用目录
\usepackage{makeidx}
% 使用图形
\usepackage{graphicx}


\usepackage{titletoc}
\dottedcontents{section}[1.5em]{}{2.3em}{1pc}

% % 页眉页脚的设定
% \usepackage{fancyhdr}
% \pagestyle{fancy}
% \lhead{\includegraphics[height=0.2in]{zte_logo}}

% \usepackage{footmisc}
\usepackage{pfnote}
\usepackage{fnpos}

% % 嵌入的代码显示
% \usepackage{listings}
% \lstset{language=C++, breaklines, extendedchars=false}
% \lstset{keywordstyle=\color{blue},
%         commentstyle=\color{SeaGreen},
%         stringstyle=\ttfamily,
%         showstringspaces=false,
%         backgroundcolor=\color{light-gray}}
% \lstset{basicstyle=\small}

\usepackage{booktabs}

\usepackage{tipa}
\usepackage{phonetic}
\usepackage[gen]{eurosym}

% \usepackage{CJK}
\usepackage{makeidx}
% \usepackage{glossary}

\usepackage{fancyhdr}

\usepackage{lastpage}
\cfoot{\textsf{\thepage\ of \pageref*{LastPage}}}

\pagestyle{fancy}
\fancyhead{}
\renewcommand{\headrulewidth}{0pt}
\pagenumbering{arabic}

\usepackage{graphicx}
\usepackage{picinpar}

\setlength{\parindent}{0in}
\setlength{\parskip}{2.5ex plus4mm minus3mm}
% \setlength{\parskip}{.3\baselineskip}

\setcounter{section}{0}

% bibtex
\usepackage[super]{natbib}
% \usepackage[globalcitecopy,labelstoglobalaux]{bibunits}
% \makeatletter
% \renewcommand{\@bibunitname}{\jobname.\the\@bibunitauxcnt}
% \makeatother

% \renewcommand{\refname}{缩略语}
% \renewcommand{\refname}{}

\makeatletter                   % Reference list option change
\renewcommand\@biblabel[1]{#1:  } % from [1] to 1: 
\makeatother                    %

\newcommand{\citenumfont}[1]{\textsf{#1}}

% \usepackage{tipa}
% \usepackage{phonetic}

\newcommand{\phonetic}[1]{{\fontspec[Colour=CC00CC]{Microsoft Sans Serif} #1}}


%%% Local Variables:
%%% TeX-master: "medical_trivia"
%%% End:


\makeindex
\begin{document}
\date{}
\thispagestyle{empty}
\renewcommand\contentsname{\textsf{Digest of The New York Times}}
{\footnotesize\textsf{\tableofcontents}}
\clearpage
\setcounter{page}{1}

\section{For Obama, Nobel Honor Is Both Benefit and Obstacle}

\lettrine{P}{resident}\mycalendar{Oct.'09}{10} Obama is given to big events at big moments, replete
with stirring speeches, lofty backdrops and stadium-size crowds.

But when Mr.~Obama walked into the Rose Garden on Friday morning, having just been awarded the Nobel
Peace Prize -- an honor that would normally be a moment of high celebration, if not the culmination
of a life's work -- he was humble and self-deprecatory, popping a hole in the balloon of his own
accomplishment. He talked about being congratulated by his daughter Malia, who proceeded to remind
him that it was the family dog's birthday, and he suggested that he was undeserving of the award.

``Let me be clear: I do not view it as a recognition of my own accomplishments, but rather as an
affirmation of American leadership on behalf of aspirations held by people in all nations,'' he
said.

Whatever it meant on the world stage, in the United States the award to Mr.~Obama was a decidedly
mixed blessing. It was a reminder of the gap between the ambitious promise of his words and his
accomplishments. It drew attention to the fact that while much of the world was celebrating him as
the anti-Bush, he had not broken as fully as he had once implied he would from the previous
administration's national security policies. And it set off another round of mocking\cite{mock}
criticism from opponents who have chafed at what they see as the charmed and entitled rise of
Mr.~Obama.

So while he accepted the award and said he would travel to Oslo to pick it up, Mr.~Obama also sought
to minimize any impression that he was basking in the glory or forgetting that he was a long way
from achieving the goals -- ridding the world of nuclear weapons, stopping global warming, bringing
peace to the Middle East, among others -- that the judges seemed to expect of him. He finished his
day, somewhat paradoxically under the circumstances, closeted with his national security team,
continuing his consideration of whether to escalate the war in Afghanistan for a second time since
taking office.

There are, without doubt, benefits to Mr.~Obama. Democrats moved quickly to portray the Nobel as an
honor to the United States after years of being an object of some scorn. For the liberal base of the
Democratic Party, the prize is a ratification\cite{ratification} of the belief that Mr.~Obama's
election would carry powerful symbolic meaning. Abroad, it provides Mr.~Obama additional stature to
be lumped with the likes of Nelson Mandela and Lech Walesa.

``I'd like to believe that winning the Nobel Peace Prize is not a political liability,'' said David
Axelrod, a senior adviser to Mr.~Obama. ``But this isn't something I gave a moment of thought to
until today. Hopefully people will receive it with some sense of pride. But I don't know; it's
uncharted waters.''

It was all but impossible to escape the fact that in the politically polarized world where Mr.~Obama
operates -- locked these days in a fierce debate over overhauling the health care plan that is
central to his presidency -- this was another complication. Even before Mr.~Obama made it to the
Rose Garden, the chairman of the Republican National Committee, Michael Steele, had issued a
statement mocking Mr.~Obama, a line of attack that was echoed on conservative blogs and radio talk
shows throughout the day.

``Can you imagine, folks, how big Obama's head is today?'' Rush Limbaugh, the conservative talk-show
host, told his listeners. ``I think it's getting so big that his ears actually fit.''

The timing of the announcement was not ideal. Fairly or not, an emerging criticism on the left and
the right is that after nine months, Mr.~Obama has not gotten much done. That was the theme of a
skit on ``Saturday Night Live'' last week that caught the notice of Mr.~Obama's political team.

Related to the question of whether his record justified the award was the notion of whether
Mr.~Obama, to some degree, remains as much a symbol as a flesh-and-blood political leader. The image
of Europe celebrating him as a global peacemaker recalled the period during the presidential race
when Senator John McCain's campaign portrayed Mr.~Obama as a vapid celebrity playing to huge
European crowds, a line of attack that left the normally sure-footed Obama team
flummoxed\cite{flummox}.

It is hardly clear that the award will have any long-term impact on the public's
perception\cite{perception} of Mr.~Obama or of his agenda. It came at a time when his prospects for
getting the Democratic-controlled Congress to pass some kind of health care bill seems stronger than
ever; as Mr.~Obama's advisers said, a victory like that could go a lot further in shaping public
perceptions than attacks from conservatives.

Still, this White House was leaving little to chance. His muted speech was followed hours later by
an e-mail message to supporters echoing many of those same self-deprecatory themes.

At the same time, Democrats quickly began making the argument that this was an honor less for this
president, and more for this country, and that no one should offer any apologies. ``It's honoring
the country,'' said Bob Kerrey, the former Democratic senator from Nebraska. ``The Nobel committee
couldn't award the peace prize to the voters of the United States, but that's what they are doing.
It's an award Americans should feel good about.''

One recognition of the tricky politics the award presents came from Gov. Tim Pawlenty of Minnesota,
a Republican who is considering challenging Mr.~Obama in 2012.

``I know there's going to be some people who are saying 'Was it based on good intentions and
thoughts, or is it going to be based on results?' '' Mr.~Pawlenty said on his weekly radio show on
WCCO in Minneapolis. ``But I think the appropriate response, or an appropriate response, is when
anybody wins a Nobel Prize, you know that is a very noteworthy development and designation and
award, and I think the proper response is to say congratulations.''

Whatever the case, the White House clearly hopes that this is one celebration where the
congratulations do not go on for too long. A few hours after his Rose Garden appearance, Mr.~Obama
was addressing dignitaries in the East Room about consumer financial regulations, in one of the
orchestrated events that have become so familiar during the first nine months of his presidency.

\section{Group Resists Korean Stigma\cite{stigma} for Unwed Mothers}

\lettrine{S}{EOUL}\mycalendar{Oct.'09}{10}, South Korea -- Four years ago, when she found that she
was pregnant by her former boyfriend, Choi Hyong-sook considered abortion. But after she saw the
little blip of her baby's heartbeat on ultrasound images, she could not go through with it.

As her pregnancy advanced, she confided in her elder brother. His reaction would sound familiar to
unwed mothers in South Korea. She said he tried to drag her to an abortion clinic. Later, she said,
he pressed her to give the child up for adoption.

``My brother said: 'How can you be so selfish? You can't do this to our parents,' '' said Ms.~Choi,
37, a hairdresser in Seoul. ``But when the adoption agency took my baby away, I felt as if I had
thrown him into the trash. It felt as if the earth had stopped turning. I persuaded them to let me
reclaim my baby after five days.''

Now, Ms.~Choi and other women in her situation are trying to set up the country's first unwed
mothers association to defend their right to raise their own children. It is a small but unusual
first step in a society that ostracizes unmarried mothers to such an extent that Koreans often
describe things as outrageous by comparing them to ``an unmarried woman seeking an excuse to give
birth.''

The fledgling\cite{fledgling} group of women -- only 40 are involved so far -- is striking at one of
the great ironies of South Korea. The government and commentators fret over the country's birthrate,
one of the world's lowest, and deplore South Korea's international reputation as a baby exporter for
foreign adoptions.

Yet each year, social pressure drives thousands of unmarried women to choose between abortion, which
is illegal but rampant, and adoption, which is considered socially shameful but is encouraged by the
government. The few women who decide to raise a child alone risk a life of poverty and disgrace.

Nearly 90 percent of the 1,250 South Korean children adopted abroad last year, most of them by
American couples, were born to unmarried women, according to the Ministry for Health, Welfare and
Family Affairs.

In their campaign, Ms.~Choi and the other women have attracted unusual allies. Korean-born adoptees
and their foreign families have been returning here in recent years to speak out for the women, who
face the same difficulties in today's South Korea as the adoptees' birth mothers did decades ago.

One such supporter, Richard Boas, an ophthalmologist from Connecticut who adopted a Korean girl in
1988, said he was helping other Americans adopt foreign children when he visited a social service
agency in South Korea in 2006 and began rethinking his ``rescue and savior mentality.'' There, he
encountered a roomful of pregnant women, all unmarried and around 20 years old.

``I looked around and asked myself why these mothers were all giving up their kids,'' Dr.~Boas said.

He started the Korean Unwed Mothers Support Network, which advocates for better welfare services
from the state.

``What we see in South Korea today is discrimination against natural mothers and favoring of
adoption at the government level,'' said Jane Jeong Trenka, 37, a Korean-born adoptee who grew up in
Minnesota and now leads Truth and Reconciliation for the Adoption Community of Korea, one of two
groups organized by Korean adoptees who have returned to their homeland to advocate for the rights
of adoptees and unwed mothers. ``Culture is not an excuse to abuse human rights.''

In 2007, 7,774 babies were born out of wedlock in South Korea, 1.6 percent of all births. (In the
United States, nearly 40 percent of babies born in 2007 had unmarried mothers, according to the
National Center for Health Statistics.) Nearly 96 percent of unwed pregnant women in South Korea
choose abortion, according to the Ministry for Health, Welfare and Family Affairs.

Of unmarried women who give birth, about 70 percent are believed to give up their babies for
adoption, according to a government-financed survey. In the United States, the figure is 1 percent,
the Health and Human Services Department reports.

For years, the South Korean government has worked to reduce overseas adoptions, which peaked at
8,837 in 1985. To increase adoptions at home, it provides subsidies and extra health care benefits
for families that adopt, and it designated May 11 as Adoption Day.

It also spends billions of dollars a year to try to reverse the declining birthrate, subsidizing
fertility treatments for married couples, for example.

``But we don't see a campaign for unmarried mothers to raise our own children,'' said Lee Mee-kyong,
a 33-year-old unwed mother. ``Once you become an unwed mom, you're branded as immoral and a failure.
People treat you as if you had committed a crime. You fall to the bottom rung of society.''

The government pays a monthly allowance of \$85 per child to those who adopt children. It offers
half that for single mothers of dependent children.

The government is trying to increase payments to help unwed mothers and to add more facilities to
provide care for unmarried pregnant women, said Baek Su-hyun, an official at the Health Ministry.
But the social stigma discourages women from coming forward.

Chang Ji-young, 27, who gave birth to a boy last month, said: ``My former boyfriend's sister
screamed at me over the phone demanding that I get an abortion. His mother and sister said it was up
to them to decide what to do with my baby because it was their family's seed.''

Families whose unmarried daughters become pregnant sometimes move to conceal the pregnancy. Unwed
mothers often lie about their marital status for fear they will be evicted by landlords and their
children ostracized at school. Only about a quarter of South Koreans are willing to have a close
relationship with an unwed mother as a coworker or neighbor, according to a recent survey by the
government-financed Korean Women's Development Institute.

``I was turned down eight times in job applications,'' Ms.~Lee said. ``Each time a company learned
that I was an unwed mom, it accused me of dishonesty.''

Ms.~Choi, the hairdresser, said her family changed its phone number to avoid contact with her. When
her father was hospitalized and she went to see him with her baby, she said, her sister blocked them
from entering his room. When she wrote to him, she said, her father burned the letters. Last year,
about three years after the birth, he finally accepted Ms.~Choi back into his home.

``That day, I saw him in the bathroom, crying over one of my letters,'' she said. ``I realized how
hard it must have been for him as well.''

\section{Afghan War Debate Now Leans to Focus on Al Qaeda}

\lettrine{P}{resident}\mycalendar{Oct.'09}{10} Obama's national security team is moving to reframe
its war strategy by emphasizing the campaign against Al Qaeda in Pakistan while arguing that the
Taliban in Afghanistan do not pose a direct threat to the United States, officials said Wednesday.

As Mr.~Obama met with advisers for three hours to discuss Pakistan, the White House said he had not
decided whether to approve a proposed troop buildup in Afghanistan. But the shift in thinking,
outlined by senior administration officials on Wednesday, suggests that the president has been
presented with an approach that would not require all of the additional troops that his commanding
general in the region has requested.

It remains unclear whether everyone in Mr.~Obama's war cabinet fully accepts this view. While Vice
President Joseph R.~Biden Jr.~has argued for months against increasing troops in Afghanistan because
Pakistan was the greater priority, Secretary of State Hillary Rodham Clinton and Defense Secretary
Robert M.~Gates have both warned that the Taliban remain linked to Al Qaeda and would give their
fighters havens again if the Taliban regained control of all or large parts of Afghanistan, making
it a mistake to think of them as separate problems.

Moreover, Mr.~Obama's commander there, Gen. Stanley A.~McChrystal, has argued that success demands a
substantial expansion of the American presence, up to 40,000 more troops. Any decision that provides
less will expose the president to criticism, especially from Republicans, that his policy is a
prescription for failure.

The White House appears to be trying to prepare the ground to counter that by focusing attention on
recent successes against Qaeda cells in Pakistan. The approach described by administration officials
on Wednesday amounted to an alternative to the analysis presented by General McChrystal. If, as the
White House has asserted in recent weeks, it has improved the ability of the United States to reduce
the threat from Al Qaeda, then the war in Afghanistan is less central to American security.

In reviewing General McChrystal's request, the White House is rethinking what was, just six months
ago, a strategy that viewed Pakistan and Afghanistan as a single integrated problem. Now the
discussions in the White House Situation Room, according to several administration officials and
outsiders who have spoken with them, are focusing on related but separate strategies for fighting Al
Qaeda and the Taliban.

``Clearly, Al Qaeda is a threat not only to the U.S.~homeland and American interests abroad, but it
has a murderous agenda,'' one senior administration official said in an interview initiated by the
White House on Wednesday on the condition of anonymity because the strategy review has not been
finished. ``We want to destroy its leadership, its infrastructure and its capability.''

The official contrasted that with the Afghan Taliban, which the administration has begun to define
as an indigenous group that aspires to reclaim territory and rule the country but does not express
ambitions of attacking the United States. ``When the two are aligned, it's mainly on the tactical
front,'' the official said, noting that Al Qaeda has fewer than 100 fighters in Afghanistan.

Another official, who also was authorized to speak but not to be identified, said the different
views of Al Qaeda and the Taliban were driving the president's review. ``To the extent that Al Qaeda
has been degraded, and it has, and to the extent you believe you need to focus on destroying it
going forward, what is required going forward?'' the official asked. ``And to prevent it from having
a safe haven?''

The officials argued that while Al Qaeda was a foreign body, the Taliban could not be wholly removed
from Afghanistan because they were too ingrained in the country. Moreover, the forces often
described as Taliban are actually an amalgamation of militants that includes local warlords like
Gulbuddin Hekmatyar and the Haqqani network or others driven by local grievances rather than
jihadist ideology.

Mr.~Obama has defined his mission in Afghanistan and Pakistan as trying ``to disrupt, dismantle and
defeat Al Qaeda and other extremist networks around the world.'' But he made it clear during a visit
to the National Counterterrorism Center on Tuesday that the larger goal behind the mission was to
protect the United States. ``That's the principal threat to the American people,'' he said.

Robert Gibbs, the White House press secretary, said Wednesday that Mr.~Obama's ``primary focus is on
groups and their allies that can strike our homeland, strike our allies, or groups who would provide
safe haven for those that wish to do that.''

The discussion about whether the Taliban pose a threat to the United States has been at the heart of
the administration's debate about what to do in Afghanistan. Some in the Biden camp say that the
Taliban can be contained with current troop levels and eventually by Afghan forces trained by the
United States.

Moreover, they suggest that the Taliban have no interest in letting Al Qaeda back into Afghanistan
because that was what cost them power when they were toppled by American-backed Afghan rebels in
2001.

``The policy people and the intelligence people inside are having a big argument over this,'' said
Leslie Gelb, president emeritus of the Council on Foreign Relations who has advised Mr.~Biden. ``Is
the Taliban a loose collection of people we can split up? Can we split the Taliban from Al Qaeda? If
the Taliban comes back to power in parts of Afghanistan, are they going to bring Al Qaeda back with
them?''

Some analysts say that the Taliban and Al Qaeda have actually grown closer since the first American
bombs fell on the Shomali Plain north of Kabul eight years ago Tuesday.

``The kind of separation that existed between the Taliban and Al Qaeda in 2001 really doesn't exist
anymore,'' said Anthony H.~Cordesman, a scholar at the Center for Strategic and International
Studies who has advised General McChrystal. ``You have much more ideological elements in the
Taliban. In the east, they're really mixed in with Al Qaeda.''

Frances Fragos Townsend, who was President George W.~Bush's homeland security adviser, said the two
groups remained linked.

``It's a dangerous argument to assume that the Taliban won't revert to where they were pre-9/11 and
provide Al Qaeda sanctuary\cite{sanctuary},'' she said. Referring to General McChrystal, she added,
``If you don't give him the troops he asked for and continue with the Predator strikes, you can kill
them one at a time, but you're not going to drain the swamp\cite{swamp}.''

Officials said Wednesday that General McChrystal's official request for additional forces was
forwarded to Mr.~Obama last week. Mr.~Gates's spokesman, Geoff Morrell, said Mr.~Gates had given
Mr.~Obama ``an informal copy'' at the president's request.

The meeting on Wednesday was Mr.~Obama's third with his full national security team. Another is
scheduled for Friday to talk about Afghanistan and then a fifth is planned, possibly for next week.
Mr.~Gibbs said the president was still several weeks away from a decision.

\section{China Aims to Steady North Korea}

\lettrine{S}{EOUL}\mycalendar{Oct.'09}{10}, South Korea -- North Korea's leader gave an unusually
exuberant welcome this week to the prime minister of China, whose trip was intensely monitored by
the rest of the world for progress on efforts to halt North Korea's nuclear weapons program.

But the deal struck by the North Korean leader Kim Jong-il and Prime Minister Wen Jiabao of China
had far more to do with the two countries' shared goal -- stability in North Korea.

When Mr.~Kim met Mr.~Wen on Monday, North Korea said it was ready to return to six-nation talks on
ending its nuclear program -- but only if it saw progress in bilateral talks with the United States,
state-run media in China and North Korea reported Tuesday. For months, North Korea had insisted that
it would never return to the talks demanded by the United States, calling them ``dead.''

North Korea's reversal came after China signed a series of agreements that promised aid for the
North and an expansion in economic exchanges, including the construction of another bridge across
their tightly controlled river border.

In somewhat ceremonial language, the two governments also vowed to support each other ``for
generations to come'' -- a pledge\cite{pledge} interpreted as a reference to the ill health of
Mr.~Kim and the possibility that the youngest of his three sons, Kim Jong-un, is being groomed as
his successor.

It was unclear whether Kim Jong-un was formally introduced to Mr.~Wen.

A belligerently anti-American North Korea has often been considered as China's buffer against
American influence in the region. China also wants to ward off an implosion by North Korea, which
would endanger the stability of its own border area.

``The results of Wen Jiabao's trip show that China's foremost concern is to secure stability in
North Korea,'' said Han Suk-hee, an expert on Chinese-North Korean relations at the Graduate School
of International Studies at Yonsei University in Seoul. ``The deals they signed are aimed at
ensuring stability in North Korea even after Kim Jong-il is gone. China effectively announced that
it did not agree with the United States and South Korea on sanctions against North Korea.''

The joint announcement gave China something on the nuclear front to show the world -- if only a
vaguely worded promise to possibly return to talks -- that might ease international pressure on it
to do more on North Korea.

But the limits of the deal were underlined by the South Korean news agency Yonhap, which said North
Korea was in the final stage of restoring its plutonium plant. That plant was mothballed after an
international agreement in 2007 but revived in a fury after the North's missile and nuclear tests
prompted new international sanctions this year.

North Korea has repeatedly used negotiations over ending its nuclear weapons program as a way of
extracting aid and diplomatic concessions from other countries. North Korea has so far received
\$2.2 billion under failed international deals aimed at persuading it to dismantle its nuclear
facilities, according to a report this week by Kwon Young-se, a governing party lawmaker in South
Korea.

While the economic assistance China promised may elevate its leverage over North Korea, it may also
deprive the United States and South Korea of theirs by sapping the strength of sanctions, analysts
said.

``This is a breakthrough for Kim Jong-il,'' said Chang Yong-seok, research director at the nonprofit
Institute for Peace Affairs in Seoul. ``China reconfirmed that it remained the biggest patron of the
North Korean regime. It's a boost for Kim Jong-il's grip on power.

``But one can imagine frustration in Washington.''

Ian C.~Kelly, a State Department spokesman, said the government was waiting to hear the details from
the Chinese.

South Korea said it had already expressed concern about whether the Chinese moves violated United
Nations Security Council sanctions.

``We are expecting China to explain the details of its economic cooperation programs with North
Korea and whether they violated the Security Council resolutions,'' Foreign Minister Yu Myung-hwan
told a group of South Korean news media editors, according to Yonhap.

Mr.~Yu described Beijing's diplomacy as a combination of sanctions and engagement arising out of a
concern that isolation and pressure alone would drive North Korea only to strengthen its nuclear
weapons programs.

China's trade and aid have become more crucial to North Korea's survival, especially with South
Korea refusing to send aid to the North during the past year and a half. Last year, trade between
China and North Korea reached \$2.79 billion, up 41.3 percent from 2007. That has slipped in the
first eight months of this year to \$1.6 billion, a fall of 6.2 percent from the same months last
year.

Some analysts say that China's influence over North Korea may be overstated.

``China has long sought to parlay economic engagement into political influence over North Korea, ''
said Bruce Klingner, senior research fellow at the Heritage Foundation. ``Beijing's efforts,
however, have been for naught. Despite extensive Chinese government largess and business engagement,
the Chinese leadership was unable to persuade North Korea to abandon either nuclear weapons program
nor prevent long-range missile launches and nuclear tests in 2006 and 2009.''

Even in the North's discussion of returning to talks, it held to its longstanding approach of
framing the issue of its nuclear weapons program largely as a matter between itself and the United
States. ``We expressed our readiness to hold multilateral talks, depending on the outcome'' of
bilateral talks with the United States, said the North's state-run Korean Central News Agency in a
report on the Kim-Wen meeting. ``The six-party talks are also included in the multilateral talks.''

``The hostile relations between the D.P.R.K. and the United States should be converted into peaceful
ties through the bilateral talks without fail,'' the agency quoted Mr.~Kim as saying, using the
initials for the Democratic People's Republic of Korea, the North's official name.

Washington said recently that it was open to direct negotiations with the North to coax it back to
the six-nation talks, which also include South Korea, China, Japan and Russia.

But both the United States and South Korea have said that bilateral discussions should not be a
substitute for six-nation negotiations and that they will maintain sanctions until North Korea
relents on its nuclear program.

\section{Climate Agency Sees China's Efforts Paying Dividends}

\lettrine{L}{ittle}\mycalendar{Oct.'09}{10} good can be said about the worst economic slump since the
1930s, but it has produced at least one piece of positive news: the downturn will make it a bit
easier to slow the rise in emissions responsible for climate change.

The International Energy Agency made that prediction in a report Tuesday on global greenhouse gas
emissions. Because of slower economic growth, the agency slashed, by 5 percent, its estimate of how
much greenhouse gas emissions will be produced in 2020.

But the energy agency also cautioned against complacency, stressing that reaching a deal in climate
talks to be held in Copenhagen at the end of the year is crucial to limiting the rise in global
temperatures.

Another reason for cautious optimism, the report said, is that China will be able to slow the growth
of its emissions much faster than commonly assumed because of its rising investment in wind and
nuclear energy and its newfound emphasis on energy efficiency.

But avoiding some of the worst consequences of climate change will still require significant and
rapid investments in clean technology, and more meaningful cuts in carbon emissions, the report
said.

``This gives us a chance to make real progress toward a clean-energy future, but only if the right
policies are put in place promptly,'' said the agency's executive director, Nobuo Tanaka.

As a result of the economic slump, global emissions of carbon dioxide, the main greenhouse gas, are
expected to decline by 3 percent this year, the steepest drop in the 45 years according to figures
compiled by the agency. That compares with an average growth of 3 percent a year over the last
decade.

The report outlines how governments can achieve additional cuts through energy efficiency and
investments in clean technologies. The goal is to keep global temperatures from rising more than 3.6
degrees Fahrenheit. Meeting that target will require reducing emissions by 23 percent in 2030
compared with what they would otherwise be, the agency said.

``The message is simple and stark: if the world continues on the basis of today's energy and climate
policies, the consequences of climate change will be severe,'' Mr.~Tanaka said.

The cost of making the necessary changes would be \$10 trillion from 2010 to 2030, the agency said,
but lower spending on energy bills would offset much or all of that investment. Every year of delay,
however, would add \$500 billion to the necessary investments, the agency said.

One of the report's main findings was that China's recent energy policies could achieve much bigger
cuts than expected.

China overtook the United States as the world's top emitter of global warming gases in 2007 after
more than a decade of rapid industrialization. But China is building more nuclear power plants and
increasing the role of renewable energy, raising efficiency standards in new buildings, and
gradually moving the economy from its manufacturing base to services, the agency said.

If China manages to achieve the predicted savings, that nation will be ``at the forefront of all
global efforts to combat climate change,'' Fatih Birol, the energy agency's chief economist, said in
an interview.

Continued progress in China could increase the pressure on the United States, which has so far
initiated no overall policy to reduce emissions, as well as on other industrialized nations.

While calling for deep cuts in emissions, the report acknowledged that profound transformations in
energy usage would be necessary to achieve the cuts. In the transportation sector, for example,
which is dominated by fossil fuels, 60 percent of worldwide car sales by 2030 would have to be
either hybrids, plug-in hybrids or electric vehicles to meet the agency's goals.

\section{In China, a Headless Mao Is a Game of Cat and Mouse}

\lettrine{I}{t}\mycalendar{Oct.'09}{10}'s not the kind of sculpture of Chairman Mao you typically
see in China. He's on his knees as a supplicant, confessing; his body language and facial expression
indicate deep remorse\cite{remorse}. What's more, the head of this life-size bronze statue, titled
``Mao's Guilt'' and created by the artist brothers Gao Zhen and Gao Qiang, separates from the body
-- by design.

Exhibitions by the Gao brothers, whose work the authorities find politically challenging, have been
shut down in the past, and their studio has been raided. So they keep the head of Mao hidden in a
separate location -- reuniting it with its body only on special occasions to show friends and
colleagues. Normally, the body of the statue remains headless, unidentifiable and nonthreatening.

``It's something I hope all Chinese people will one day be able to accept and understand,'' Gao
Zhen, 53, said of the work. ``We wanted to portray him as a human being, a regular person confessing
for the wrongs he's committed.''

On Sept.~3 the head came out for a Gao brothers ``party'' -- the code name for one of the
invitation-only private exhibitions they hold several times a year. The location of the exhibition
was not disclosed until several hours beforehand and spread via word of mouth and coded text
message. Outside the closed doors of their private home studio, a staff member kept watch for
unwelcome visitors.

Removable heads and underground exhibitions are just two of the guerrilla tactics the Gao brothers
have employed, often with the help of Melanie Ouyang, their broker, to enable fans and friends to
view their work. The Gaos are part of a generation of avant-garde Chinese artists who are pushing
the boundaries of artistic expression. In the increasingly open Chinese art world, nudity is
commonplace where it used to be forbidden, and art parodying the Cultural Revolution has become so
ubiquitous that it is pass\'e. Still, the Gaos are a reminder that, especially as China celebrates
the 60th anniversary of the Communist revolution, limits to expression remain: although artists are
increasingly free to deal with social and political topics, works that explicitly criticize Chinese
leaders or symbols of China are still out of bounds.

``Ash Red,'' a 2006 exhibition the Gao brothers openly advertised and held in their studio in the
798 Art District here, was suppressed by authorities. Several men representing the government walked
into the gallery and presented a list of works that ``needed to be removed,'' said Gao Qiang, 47.
Posters and catalogs for the show were banned, and interviews the brothers had lined up with local
news media were canceled. For several weeks after ``Ash Red'' was shut down, two guards stood
outside the doors of the Gao brothers' home studio, discouraging people from coming inside.

The 798 Arts District has a local management office that, among other things, keeps an eye on art it
deems unacceptable and detrimental to the district. ``They receive pressure from above,'' said Gao
Yuewen, 29, a staff member at the Gao studio, who noted that the Gao brothers were ``classified
differently'' from other artists by the authorities, meaning that they were suspect.

In March another sculpture by the Gaos, ``To Catch a Lady,'' a reproduction of a photograph of the
police hauling away a prostitute during a raid on a brothel, was whisked away during the night. Only
after the brothers filed a complaint to the police did the authorities admit to removing the piece.

The Gao brothers' most extensive work is both explicit and critical, seeking to recast Chairman Mao
-- a figure in China who is simultaneously capable of arousing deep emotions of pain and despair, as
well as admiration, love, and pride -- as a flawed figure. In television movies here Mao's wartime
triumph in resisting Japanese invaders is continually on replay, while the Great Leap Forward, the
dark side of the Cultural Revolution and Mao's failed policies are played down.

For many older Chinese, Mao remains sacrosanct\cite{sacrosanct}. But for a younger generation of
Chinese, Mao, who died in 1976, feels increasingly irrelevant, and there is little shock value in
the Gaos' portrayals of him. In China's larger, more cosmopolitan cities, remnants of Mao's
personality cult are less prevalent\cite{prevalent}. Reciting phrases from Mao's ``Red Book'' has
long since been replaced by shopping for laptop computers, Mini Coopers and other ``ming pai,'' or
famous brand-name consumer items.

For the Gao brothers, Mao holds a more personal meaning. During the Cultural Revolution their father
was labeled a class enemy and dragged off to a place that was ``not a prison, not a police station,
but something else,'' Gao Zhen said. After 25 days had passed, the family members were told he had
committed suicide.

They think otherwise: ``If someone didn't like you at that time, they arbitrarily labeled you a
class enemy,'' Mr.~Gao added. ``We came to Beijing to petition our father's death.'' Eventually the
family was given the equivalent of about \$290 in compensation. ``That was a very painful period of
our life,'' Mr.~Gao continued. ``We were six brothers and a single mother; we didn't have a penny.''

That defining event in their childhood has been both the basis and motivation for much of the
brothers' work, which often seeks to put a spotlight on people, places and events in Chinese society
that are taboo. The mural-size painting ``Forever Unfinished Building'' shows a smorgasbord of
characters representative of many different sectors of society sprawling across a construction site.
Front and center is a Chinese woman refusing to be displaced from her home by rich developers.

Still, many Chinese who are critical of the Gaos' work say it lacks subtlety. ``I understand what
they're trying to say, but I think their pieces are sensationalist -- they're too direct and
gaudy,'' says Feng Ling, 23, an art student who recently came to the Gao brothers' home studio and
saw ``The Execution of Christ,'' in which a firing squad of Chairman Maos take aim at Jesus.

Many artists in China have learned to work political meaning and criticism into their art without
being as obvious as the Gao brothers. Liu Wei, for example, in his recent work ``A Lifestyle,''
placed various pieces of Chinese exercise equipment, found in parks all over the country, in a giant
iron cage that looks like a jail, suggesting the extent to which daily activity and freedom are
circumscribed by state power.

``Most artists nowadays have learned to make political commentary without being overtly political,
so these kinds of 'underground showings' are happening less than they did in the 1990s,'' said Phil
Tinari, 30, founding editor of Artforum's Chinese-language Web site, artforum.com.cn.

The Gaos, in their latest underground show, sought both to disarm their critics and give a sendoff
to the Mao theme; they say ``Mao's Guilt'' will be their last work based on the Chinese leader ``for
a while,'' partly because of the prevalence of art parodying the Cultural Revolution, and also
because his image has been recycled over and over in popular art and film.

``The Gao brothers' work on Mao is provocative\cite{provocative} for many mainland Chinese,'' said
Kai Heinze, 33, director of the Faurschou Gallery in the 798 district. ``Their work sets off a
trigger, challenging people here to understand and tolerate a view of modern Chinese history that
admits shortcoming,''

\section{Untold Stories of China and Taiwan}

\lettrine{W}{hen}\mycalendar{Oct.'09}{10} Ying Meijun bade farewell to her 1-year-old son at the
train station in September 1949, little did she know that it would be 38 years before she saw him
again.

The baby was crying so much that she decided not to take him onto the overcrowded train, so she left
him in the care of his grandmother.

Thinking they were only leaving China temporarily, she promised: ``We'll be back soon.''

By the time she saw her first-born child again in 1987, he was a 40-year-old man wearied by years of
hard labor on a mainland Chinese farm. Fighting back tears, he told his elderly parents how, as a
young child, he used to chase trains that went pass their front door, shouting, ``Mother! Mother!'',
thinking that she would be on them.

Ms.~Ying and her husband, Lung Huaisheng, who was an officer in the military police under Chiang
Kai-shek's Kuomintang government, fled with his family to Taiwan a few months after the Communist
Party declared itself the new ruler of China on Oct.~1, 1949.

Even in his old age, Lung Huaisheng often wept as he took out the shoe soles that his mother knitted
and gave him when they saw each other for the last time at the train station.

These family memories are just some of the heart-wrenching stories told by their daughter, Lung
Yingtai, a Taiwan-born author and University of Hong Kong professor, in her latest book ``Da Jiang
Da Hai 1949'' (``Big River, Big Sea -- Untold Stories of 1949''). The book is published by Taiwan's
CommonWealth Magazine and Hong Kong's Cosmos Books.

Ms.~Lung, who was born two years after the family moved to Taiwan, is a leading cultural critic,
well-known for her sharp and candid writing. Her book of social-political criticism, ``The Wild
Fire,'' published in 1985 when Taiwan was still under Kuomintang's one-party rule, was seen as
influential in the democratization of the island.

Her new book is a tribute to the tens of millions of people ``who were trampled on, humiliated and
hurt by the era.'' It tells the story of the many Chinese families that were broken up by the civil
war that ended in the Kuomintang's defeat in 1949, with some two million escaping to Taiwan. Many,
like her own parents, hastily said goodbye to loved ones in mainland China and would never see them
again.

Apart from the stories of her own family and other Chinese people born in that era -- including
President Ma Ying-jeou of Taiwan -- there are tales of elderly people who as young men fought for
the Kuomintang, the Communist Party, or both, and even Japan (which ruled Taiwan from 1895 until
1945).

Many have not openly talked about their experiences. One 89-year-old man who was held by the
Japanese as a prisoner of war told Ms.~Lung he waited all his life to tell his story.

Ms.~Lung's book has become an instant best seller -- more than 100,000 copies have been sold in
Taiwan and 10,000 in Hong Kong since its publication in early September. Ms.~Lung, who will be
giving a talk at the Frankfurt Book Fair on Oct.~15, said the book did not have an English-language
publisher yet. Although Ms.~Lung had expressed a wish to publish it in mainland China, it seems
almost impossible now, as the government has banned all Internet articles and discussions on the
book.

Ms.~Lung hopes to break down her readers' preconceptions about events around 1949. Under Communist
rule, many mainlanders regard Taiwan as a renegade province that should be taken back by force if
necessary.

``I want to give them a different perspective,'' she said.

As the Taiwan-born offspring of mainland refugees herself, she wants mainland readers, particularly
political leaders like President Hu Jintao of China, to learn about the pain and sufferings of the
people of Taiwan.

``When will there be no war? It's when you can see your enemy's wounds, then you won't be able to
pick up your gun,'' she said.

She hopes the book will make people in China and Taiwan abandon long-held suspicions and prejudices
regarding each other.

``If all that the leaders can think about are political negotiations'' and economic interests ``and
there is no genuine understanding of emotions, then the foundation of peace would not be solid
enough,'' she said.

While researching her book, Ms.~Lung discovered that residents of Changchun in the northeastern
province of Jilin had not heard of the People's Liberation Army's five-month siege of that city in
1948, which resulted in between 150,000 and 650,000 people dying of starvation.

Instead, what they learn about in mainland Chinese history textbooks is the P.L.A.'s ``great
victory'' when it ``liberated'' that city.

Mainland China is not the only side to edit its version of history.

The Kuomintang, which lost 470,000 troops in the northeastern battles and later fled to Taiwan, did
not mention its defeat in the textbooks of Taiwan, either.

Ms.~Lung wanted to tell this history through the tales or ordinary people.

She claims to make no political or moral judgment in her book. There is no ``right side'' or ``wrong
side'' in the stories, she says. The Kuomintang troops, the People's Liberation Army and the
Taiwanese soldiers fighting for their Japanese colonial masters are given an equal hearing. To her,
those individuals were just young people caught up in history.

``In this book I don't care about who is on the right side, the victorious or the defeated side. I
just want to show you that when you dismantle the apparatus of state, what's inside are these
individuals.''

Parts of Ms.~Lung's book also detail the stories of families amid wars and conflicts in the West,
including the loving letters written by her German mother-in-law's first husband before he died in a
Soviet prisoner of war camp during World War II.

Ms.~Lung said she included these because she wanted her Chinese readers to see their own history in
perspective.

``Chinese people on both sides of the straits tend to see history from their own national scope,''
she said. ``But actually who is righteous or unrighteous? It's a very complicated matter.''

``If we continue to be the unthinking cogs in a machine,'' she said, ``then how do you know whether
these tragic misfortunes would not be repeated?''

\section{Critics Protest Gap in Dalai Lama's Schedule}

\lettrine{W}{hen}\mycalendar{Oct.'09}{10} the Dalai Lama landed here on Monday, he set off on a
characteristically hectic, weeklong schedule including lectures, seminars, an awards ceremony and
meetings with a senior State Department official and the House speaker, Nancy Pelosi.

But one appointment not on the calendar of the Dalai Lama, the Tibetan spiritual leader, is a
meeting with the president of the United States -- a gap that has drawn protests from Republican
lawmakers, commentators and some Tibetan leaders, who say the Obama administration is snubbing him
to appease China.

In June, the White House informed the Dalai Lama that President Obama was committed to meeting him,
but not until after he visits Beijing in November, a senior administration official said.

Greeting the Dalai Lama, whom Beijing condemns as a separatist, weeks before Mr.~Obama's first
presidential trip to China could be ``substantially damaging to the relationship,'' said this
official, who spoke on condition of anonymity because of the delicate nature of the issue.

Some White House officials even worried that the Chinese might withdraw the invitation to Mr.~Obama,
the official said, though Beijing had not issued any direct or veiled threats that it would do so.

``We want to have a good U.S.-China relationship, not just for its own sake, but because if we
don't, we won't be able to help Tibet,'' the senior official said. ``If the Tibet relationship is
seen as an irritant to the U.S.-China relationship, then that will cripple our ability to be of
help.''

Every president since George H.~W. Bush in 1991 has met the Dalai Lama when he came to Washington,
usually in private encounters at the White House. In 2007, George W.~Bush became the first president
to welcome him publicly, bestowing the Congressional Gold Medal on him at the Capitol. Mr.~Obama met
the Dalai Lama as a senator.

The Tibetan leader's representative in the United States, Lodi Gyaltsen Gyari, said the Dalai Lama
accepted Mr.~Obama's explanation and looked forward to meeting him before the end of the year.

``We feel this was the right decision, and we know President Obama is very serious about this
issue,'' he said. ``They are also aware of the concerns people have, and they are dealing with the
Chinese, obviously.''

Despite the supportive words, the Dalai Lama's aides were deeply frustrated by Mr.~Obama's decision,
American officials said. At a time of rising Chinese influence, they worry that it could be used as
a pretext by leaders of other countries to refuse to see the Dalai Lama.

The prime minister of Tibet's government in exile, Samdhong Rinpoche, said in September: ``A lot of
nations are adopting a policy of appeasement. Even the U.S.~government is doing some kind of
appeasement. Today, economic interests are much greater than other interests.''

Mr.~Obama sought to reassure the Dalai Lama by sending Valerie Jarrett, a senior adviser and close
friend, to visit him in Dharamsala, in northern India, last month. She was accompanied by Maria
Otero, an under secretary of state for global affairs who is the special coordinator for Tibetan
issues. On Monday, the State Department said the Dalai Lama met with Ms.~Otero at a hotel in
Washington.

Mr.~Obama conveyed his ``respect'' for the Dalai Lama to China's president, Hu Jintao, during their
first meeting at an economic summit in London in April, a senior official said. But the Chinese
leader reiterated Beijing's position that no foreign officials should meet with him.

The Obama administration has worked hard to avoid the bumpy start that many recent presidents have
had with China. Secretary of State Hillary Rodham Clinton played down Tibet and other human rights
issues in her visit to Beijing in February, while the Treasury secretary, Timothy F.~Geithner, has
been careful not to repeat the written testimony during his confirmation hearings, in which he
accused China of manipulating its currency.

Critics say the administration is betraying the Buddhist monks who are jailed in China for fighting
for democracy in Tibet. ``I can almost hear the words of the Chinese guards saying to them that
nobody cares about you in the United States,'' Representative Frank R.~Wolf, a Virginia Republican,
said at a recent hearing.

This article has been revised to reflect the following correction:

Correction: October 8, 2009 A picture of the Dalai Lama, Tibet's spiritual leader-in-exile, with an
article on Tuesday about his visit to Washington carried an erroneous credit in some editions. It
was by Tim Sloan of Agence France-Presse -- Getty Images, not by Kevin Lamarque of Reuters.

\section{China Is Wordless on Traumas\cite{trauma} of Communists' Rise}

\lettrine{C}{HANGCHUN}\mycalendar{Oct.'09}{10}, China -- Unlike in other cities taken by the People's
Liberation Army during China's civil war, there were no crowds to greet the victors as they made
their triumphant march through the streets of this industrial city in the heart of Manchuria.

Even if relieved to learn that hostilities with Chiang Kai-shek's Nationalist Army had come to an
end, most residents -- the ones who had not died during the five-month siege -- were simply too weak
to go outdoors. ``We were just lying in bed starving to death,'' said Zhang Yinghua, now 86, as she
recalled the famine that claimed the lives of her brother, her sister and most of her neighbors.
``We couldn't even crawl.''

In what China's history books hail as one of the war's decisive victories, Mao's troops starved out
the formidable Nationalist garrison that occupied Changchun with nary a shot fired. What the
official story line does not reveal is that at least 160,000 civilians also died during the siege of
the northeastern city, which lasted from June to October of 1948.

The People's Republic of China basked in its 60th anniversary on Thursday with jaw-dropping
pageantry, but there were no solemn pauses for the lives lost during the Communist Party's rise to
power -- not for the estimated tens of millions who died during the civil war, nor the millions of
landlords, Nationalist sympathizers and other perceived enemies who were eradicated during Mao's
drive to consolidate power.

``Changchun was like Hiroshima,'' wrote Zhang Zhenglu, a lieutenant colonel in the People's
Liberation Army who documented the siege in ``White Snow, Red Blood,'' a book that was immediately
banned after publication in 1989. ``The casualties were about the same. Hiroshima took nine seconds;
Changchun took five months.''

The 40,000 who survived did so by eating insects, leather belts and, in some cases, the bodies that
littered the streets. By the time Communist troops took over the city, every leaf and blade of grass
had been consumed during the final desperate months.

There are no monuments or markers recalling the events that decimated Changchun's populace. Most
young people have no knowledge of the darker aspects of the siege, and the survivors, now in their
70s and 80s, are reluctant to give voice to long-buried trauma. ``I've always heard that Changchun
was captured without bloodshed,'' Li Jiaqi, a 17-year-old high school student, said as she sat on
the steps in front of the city's Liberation Memorial.

Chinese scholars have largely steered clear of the subject. Several historians, when asked about the
episode, declined to be interviewed. Zhou Jiewen, a retired nuclear physicist in Changchun who has
become a self-taught expert on the siege, explained that many key details, if widely disseminated,
would tarnish the army's reputation as defenders of the common man. Those include shooting civilians
who tried to escape the city and ignoring the pleas of mothers holding aloft starving children on
the other side of the barbed-wire barricades. ``To cause so many civilians to die was a great
blunder by the P.L.A. and tragedy unparalleled in the civil war,'' Mr.~Zhou said.

While history is often written by the victors, the Communist Party has never been shy about shaping
the past to serve its central narrative. Textbooks portray the revolution as the inevitable outcome
of a popular uprising; the patriotic films that have flooded television in recent months are not
subtle in their glorification of Mao's troops as munificent liberators. The unpleasant aspects of
the revolution, including innocents caught in the cross-fire, are often omitted.

``The party has no use for objective history,'' said Bao Pu, a Hong Kong publisher who infuriated
party leaders last spring by printing the memoir of Zhao Ziyang, the deposed Communist Party leader
who spent 15 years under house arrest after opposing the violent crackdown on democracy protesters
in 1989. ``The basic idea is that history can be rewritten and used as a tool of the state. But this
requires constant censorship. And it has a destructive effect on society.''

Other unintended consequences of suppressing the truth are hard to quantify. Many Chinese,
especially those who grew up during the tumultuous decades of war, famine and political persecution,
carry psychic wounds that are seldom expressed, let alone healed.

Lung Ying-tai, a University of Hong Kong professor who studied the siege of Changchun, said nearly
every elderly army officer she interviewed for her book about the civil war, ``Big River, Big Sea --
Untold Stories of 1949,'' broke down when recounting what he experienced. ``It's an unspeakable
national trauma that has not once been opened up and gently treated for 60 years,'' she said.

The book, which was published last month in Taiwan and promptly banned on the mainland, seeks to
portray the horror of the civil war through the stories of those who survived. ``There are not too
many left who can clearly remember,'' she said.

The elderly survivors who gather in Changchun's Labor Park most days are not eager to tell their
tales. But after some prompting, the details spill out. They describe babies too weak to cry, brides
sold for a morsel of food and the milewide no man's land where thousands perished in full view of
troops under orders from Gen. Lin Biao to turn Changchun into a ``dead city.''

In the first few months of the siege, food could be purchased, albeit at exorbitant prices. By the
end of the summer, people were trading thick gold rings for a biscuit.

``At first we ate rotten sorghum, then corncobs and then the bark off the trees,'' said Meng
Qinghua, 85. ``After a week of not eating you'd get very sleepy. Once that happened, you would start
to die.''

The few airdrops of aid, delivered by American planes, were quickly gobbled up by Nationalist
troops. When those stopped, the soldiers stole food from civilians at gunpoint. In the poorer
quarters of the city, according to ``White Snow, Red Blood,'' 9 of 10 families were wiped out.

Although her family was relatively well off, Zhang Yinghua said there was nothing to be bought by
the end of summer. They opened their pillows and consumed the corn husk filling. Later they boiled
and ate leather.

Then 25, Ms.~Zhang understood that swallowing such unpalatable matter was the only way to survive.
``Every day we would eat a spoonful, just enough to maintain the flicker of life, but the children
would not,'' she said. When her 6-year-old sister and her 9-year-old brother finally died, her
parents, barely able to stand, dragged their bodies to the street.

Some of those charged with enforcing the blockade have come to regret their participation. Wang
Junru said he was 15 when the Communists forced him to join a militia for teenagers. Later, he
joined 170,000 other soldiers ordered to drive back hungry civilians. ``We were told they were the
enemy and they had to die,'' he said.

Whatever zeal he had for the revolution was extinguished by the 23 years he spent in a labor camp --
punishment, he said, for insulting the relative of a party official when he was a college student.
After his release, he spent the rest of his working life hauling logs.

Now 76 and embittered, he said young people should learn about what happened in Changchun -- and
during the rest of the civil war. ``They only know the propaganda,'' he said. ``Maybe if they know
how horrible war is, they can try to avoid it in the future.''

\section{Seeing the Future in 3-D Television}

\lettrine{F}{lat}\mycalendar{Oct.'09}{10}-panel television screens can't get much flatter and
consumers don't want the screens to get much wider, so Japanese television makers are banking on a
whole new dimension to buttress their lineups.

High-definition three-dimensional TV is the future, or so Panasonic and Sony hope, as they seek to
stem a slide in prices and re-energize a market slowed by the global recession.

The biggest problem the companies face, however, is staring them right in the face. Viewers will
need to wear those goofy, ill-fitting glasses, just as they have to when watching 3-D movies in a
theater. Without them, the screen looks nauseatingly blurry.

The expected high price of 3-D systems -- the special television screens, glasses and Blu-ray DVD
players -- could also discourage all but home theater buffs.

Neither Sony nor Panasonic has announced a price for their 3-D offerings. Panasonic said market
research had shown that its 50-inch model might sell for \$2,000. And at \$50 a pair, a family of
four would be paying as much for the glasses as a small high-definition television costs.

A poll from the research firm In-Stat found that while 64 percent of respondents expressed at least
some interest in watching 3-D programming at home, 25 percent of those interested said they would
not pay more than for a regular television.

``From a consumer experience standpoint, I'm skeptical, because I can't see consumers sitting at
home wearing glasses,'' said David Gibson, the head of research at Macquarie Capital Securities.
``It's a good idea, but it's a niche market for now.''

Still, Panasonic is betting that its eye-popping imagery will win over more than a few fans. A
prototype 50-inch Viera plasma 3-D set attracted long lines at the Ceatec electronics show in Tokyo
this week. Beating rivals to the market could also provide a boost for its plasma technology, which
has been losing out to liquid-crystal displays.

In a 3-D clip from ``Toy Story,'' Pixar's computer-animated film, Rex the green tyrannosaurus looked
so ready to leap off the screen that the crowd gasped.

The technology works by rapidly alternating between left and right frames of the video. Viewers wear
glasses that sync with the television over an infrared signal. The right frame is seen only with the
right eye and the left frame with the left eye, creating the illusion of depth.

``We are serious about the future of 3-D,'' said Masayuki Kozuka, a leader on Panasonic's 3-D team.
``We will not have succeeded until half of all TVs we sell are 3-D TVs.''

Sony showed off a similar technology, which it hopes to deploy in some Bravia televisions and Vaio
notebook PCs in 2010. PlayStation 3 video game consoles could also be fitted with 3-D technology,
said Yutaka Nakamura, a Sony spokesman.

Samsung Electronics of South Korea has also shown a 3-D prototype.

The movie industry has been the biggest cheerleader of 3-D technology, and television makers are
hoping that enthusiasm will help persuade consumers to replicate the 3-D experience at home.
Hollywood raised its output of 3-D movies in recent years to lure people into theaters. There will
be around 7,000 3-D movie screens by the end of 2009, according to Sony.

Hollywood, in fact, has long been obsessed with the technology. In the 1950s, a flurry of 3-D films
like ``Man in the Dark'' and ``House of Wax'' were released. Another round of stereoscopic 3-D films
that had viewers wearing red-and-green glasses hit screens in the 1970s. But with poor color quality
and unsteady images that induced nausea, 3-D didn't take off.

Hollywood hopes that the third time will be the charm. Studios are now eager to sell 3-D films on
Blu-ray discs, too. Panasonic and 20th Century Fox worked together to promote the forthcoming James
Cameron sci-fi thriller ``Avatar,'' the first major nonanimation film being released worldwide in
both 2-D and 3-D formats.

Television makers, in turn, hope that by helping to promote 3-D technology on the big screen,
broadcasters will follow Hollywood's lead and start offering 3-D programming. Many stumbling blocks
remain, however, including the added costs of filming and broadcasting 3-D images.

Some analysts are upbeat on the prospects for 3-D television. Alfred Poor, an analyst for the GigaOM
Network of high-tech news Web sites, predicted in a recent report that TV makers could ship up to 46
million 3-D sets by 2013, adding that high sales volume would drive prices down. Analysts at Gartner
played it safer in their ``Hype Cycle'' report, giving a fuzzier timeline of five to 10 years for
3-D to catch on.

Japanese electronics makers have for years been losing market share to competitors from South Korea,
while more recently watching competitors from Taiwan and China take even more. By capitalizing on
Japan's technological prowess, manufacturers here can at least differentiate themselves from
lower-priced rivals, who sell under brand names like Vizio and Haier.

It is a risky bet. High-end spending could be the last to rebound. ``The world has become
increasingly frugal, compounding the plight of Japanese producers, as demand for high-end
manufactured goods dries up,'' said Ryutaro Kono, chief economist for Japan at BNP Paribas.

The theme at the Tokyo electronics show was decidedly high-end. Toshiba's Cell Regza 55-inch
television does not yet offer 3-D, but can display eight high-definition broadcasts at once in eight
windows. The set, which uses the Cell microchip found in PlayStation game machines, can also record
all eight channels nonstop for 26 hours, according to a company spokeswoman, Kaori Hiraki. The price
tag: about a million yen, or \$11,000.

Meanwhile, a Hitachi prototype used a camera sensor to allow users to control the TV by gesturing.

``You're saying to the TV, 'Notice me,' '' said Takashi Matsubara, a Hitachi spokesman.

\section{Back to Basics: Yes, Sergeant!}

\lettrine{I}{F}\mycalendar{Oct.'09}{10} Mark Roozen, a personal trainer in Colorado Springs, set his
group conditioning classes to music, the playlist could start with ``I've Been Working on the
Railroad.'' Mr.~Roozen's routines are as likely to incorporate logs, wheelbarrows and sandbags as
circuit machines, Pilates equipment and other gym staples.

``We'll take a big, old tractor tire and a sledgehammer and hit the tire 25 times'' in a typical
class, said Mr.~Roozen, owner of Performance Edge Training Systems and director of certification for
the National Strength and Conditioning Association. ``Those tires can weigh 250 pounds, and they
have to flip the tire.''

He makes use of other low-tech equipment that has resurfaced in recent years, including kettle bells
(borrowed from the early Soviet military); stability balls (popularized in Switzerland in the
1960s); air-filled balance boards (a physical therapy staple); and medicine balls (``In ancient
times, they'd fill an animal bladder with water or sand,'' Mr.~Roozen said of the balls' origins).

While such rough-hewn techniques and gear may look old-fashioned, they comport with a modern shift
away from developing individual muscle groups and toward so-called functional fitness, which refers
to overall strength and comfort in performing everyday activities, like lifting, walking and
reaching, along with cardiovascular health.

``Rather than just sculpting the body to have great biceps, or for women who don't want the waggle
under their arms, we want you to be able to move functionally and have a healthy body,'' said Shelby
Murphy, editor of Personal Fitness Professional magazine and owner of a personal training studio in
Granbury, Tex. Staples of Ms.~Murphy's group conditioning classes, she said, include push-ups,
pull-ups, mountain climbers and bear crawls, along with squats, walking lunges and other ``big
movements using a combination of muscles.''

Such exercises ``mimic nature,'' she said. ``You're running, you're jumping, you're squatting,
you're hopping, you're skipping -- all of those things our body was designed to do that we're
getting away from.''

Still, these drills have changed since many boot campers last heard a physical education teacher's
whistle, evolving to combine several components at once, Mr.~Roozen said. To increase muscle power
across a range of motion, for instance, he updates weight drills like squats, curls and presses by
adding lateral motions, one-armed or one-legged lifts and other variations.

For organizers of workouts, an advantage to these is that they don't require expensive equipment,
lending themselves to large class sizes.

Kelli Calabrese trains instructors nationally for the Adventure Boot Camp program and also owns
Argyle Adventure Boot Camp in the Dallas suburb of Argyle. Her daily 5:30 a.m. class usually draws
about 60 participants at \$19.95 a class. The group meets at a church-school campus, making use of
fields, blacktop and bleachers (for running drills, of course).

When she began outdoor classes in 2007, Ms.~Calabrese said, ``I didn't want to drag out a truckload
of stuff'' every day. So she improvises, setting up a rotation of workout stations using traffic
cones and baseball bases, each with its own physical challenge. ``Your body has no idea if you're
using a \$4,000 machine or if you're using a rock,'' she said.

Or, as one of her students, Patrick McCarrell, put it: ``The parking lot now is a workout facility.
And there's no air-conditioning and no heater.''

Mr.~McCarrell, 40, said he Googled ``boot camp'' after watching ``The Biggest Loser'' on NBC. In
April, he enrolled in Ms.~Calabrese's program, which runs in four-week sessions with a fifth week
off between camps.

``My goal for the first few camps was just to survive,'' said Mr.~McCarrell, who said he has since
lost 16 pounds.

Julie Crecelius, 38, who also attends the Argyle camp, said she used to work with a personal trainer
but could afford only two hours a week. When she exercised at a fitness center, she said, work and
family demands often derailed her evening workouts. With a morning boot camp, she said, ``I can get
out and get home in time for my husband to get to work.''

Despite the popularity of the camp format, group strength training and conditioning classes are
becoming more common in fitness clubs, too. The International Health, Racquet and Sportsclub
Association reported in its annual survey that nearly a third of member fitness centers offered some
form of so-called boot camp classes in 2008.

Equinox fitness centers offer some 20 varieties of boot camp and, this summer, added four types of
kettle bell workouts to their schedules.

AYC Health and Fitness, a personal training center in Prairie Village, Kan., runs ``30-minute
express'' boot camps every morning, both on-site and at corporate locations. While less expensive
than personal training, the classes often act as a ``feeder'' for higher-end services, said the
owner, Greg Justice.

``It's almost imperative that you have some type of boot camp offering,'' Mr.~Justice said, ``for
pure and simple bottom-line dollars.''

Beyond convenience and economy, the retro fitness kick may also -- like the revival of Pop Rocks and
tie-dye -- appeal to a taste for boomer nostalgia. ``When's the last time most adults jumped rope?''
Ms.~Calabrese asked. ``Second grade? We play tag. We play games. We do baton races, Indian runs,
shuttle runs, all these things that go back to gym class.''

\section{In Test of Water on Moon, Craft Hits Bull's-Eye}

\lettrine{M}{OUNTAIN}\mycalendar{Oct.'09}{10} VIEW, Calif. -- Scientists who slammed a bus-size
projectile into the Moon on Friday, hitting exactly the spot they were aiming for, say it will take
weeks to figure out what they did, and did not, see.

More than 230,000 miles from Earth, the Lunar Crater Observation and Sensing Satellite hit a
bull's-eye -- two, actually -- on the Moon. But the initial images at least, both from the
spacecraft and telescopes on Earth, failed to capture expected plumes of debris rising out of the
impacts.

At 4:31 a.m. Pacific time (7:31 a.m. Eastern time), one piece of Lcross slammed into the bottom of a
crater, excavating hundreds of tons of the Moon. Trailing four minutes behind, a second piece sent
its observations back to Earth before it also slammed into the same crater.

The absence of a visible plume disappointed the hundreds of enthusiasts who braved a chilly evening
outside at the NASA Ames Research Center here, which operated the spacecraft, to watch the live
images transmitted from the spacecraft.

But Anthony Colaprete, the mission's principal investigator, was ecstatic. ``We got the data we need
to address the questions,'' he said at a news conference.

Of greatest interest is whether there is water ice hidden in the crater's perpetual darkness and
frigidness. The data could play into the debate over where NASA's human spaceflight program should
aim next, whether to return to the Moon or head elsewhere in the solar system neighborhood. The
presence of large significant amounts of water could make it easier to set up future settlements
with the ice providing water and oxygen.

Data from NASA's Lunar Reconnaissance Orbiter has already confirmed the presence of hydrogen deep
within permanently shadowed craters near the Moon's poles, and hydrogen is most likely in the form
of water.

Lcross (pronounced L-cross) is a \$79 million companion mission to the Lunar Reconnaissance Orbiter,
sharing the same rocket into space in June. The mission designers took advantage of what would have
otherwise been space junk -- the rocket's two-ton, second stage -- and turned it into a projectile
to hit the Moon, shepherded by a car-size spacecraft.

While the orbiter entered orbit around the Moon, Lcross swung into a wide polar orbit around the
Earth that, by design, would intersect with the Moon's path four months later at 5,600 miles per
hour, or twice the speed of a bullet.

The target of Lcross was Cabeus crater, about 60 miles wide near the south pole.

Dr.~Colaprete said the Lcross spacecraft captured a flash of light as the upper stage hit the bottom
of Cabeus and then captured a thermal image of the resulting crater, about 60 to 65 feet wide, close
to what had been predicted.

What was missing was the plume of debris that was knocked out by the impact. ``We saw a crater. We
saw a flash,'' Dr.~Colaprete said. ``Something had to happen in between.''

Ground-based telescopes that had been pointed to that crater at the bottom of the Moon also failed
to spot the theatrics, at least at first glance of their images.

``As far as I can tell from our quick processing, we did not see any plume,'' said William C.~Keel,
a professor of astronomy at the University of Alabama who was operating a telescope at the Kitt Peak
National Observatory in Arizona.

If the expended rocket stage hit a rocky area or a slope, the debris may not have been tossed high
enough to reach sunlight and thus not have been seen.

But Lcross's spectrometers -- instruments that break down light into wavelengths and detect subtle
changes, perhaps from vapor or fine particles not visible to the eye -- did observe changes before
and after the impact. Dr.~Colaprete said the spectrometer data could identify the water and other
molecules. ``When I saw actually the spectra, I was like, ``We got something,' '' he said.

The analysis will take at least a few days and maybe weeks.

``We're going to take our time'' Dr.~Colaprete said, ``and build up a case for water and the ejecta,
if it's there or the case against it if it's not there.''

While Lcross itself had the best view of the first impact, a host of telescopes in space and on
Earth, including the Hubble Space Telescope and the Keck telescope in Hawaii, were also gazing at
the Moon. The other telescopes also tried to observe the second impact of the Lcross spacecraft.

Several hundred people spent a chilly night on a grassy lawn at the Ames Research Center. They
listened to Charlie Duke, one of the Apollo 16 astronauts, and watched three space-themed films --
``Fly Me to the Moon,'' ``The Dish'' and ``October Sky'' -- projected on a 40-foot screen.

Then they watched the same NASA coverage of the mission, streamed over the Internet, that they could
have watched at home.

``It's adventurous and nerdy at the same time,'' said Karin Atkins of Sunnyvale, Calif., one of
those pulling an outdoors all-nighter.

\section{Death Sentence for Chinese Brawl}

\lettrine{C}{hina}\mycalendar{Oct.'09}{12} sentenced one man to death and another to life in prison
on Saturday for their roles in a deadly toy factory brawl that was blamed for setting off riots in
western China's Xinjiang region this past summer, according to the official Xinhua news agency.

Two courts in southern Guangdong Province, where the toy factory was located, also sentenced nine
other people to prison terms ranging from five to eight years for taking part in the fights in June,
according to Xinhua.

Xiao Jianhua was sentenced to death as the ``principal instigator'' in the factory melee that pitted
a group of Han Chinese against a group of Turkic-speaking Uighurs from the Xinjiang region.

Mr.~Xiao and Xu Qiqi, who was sentenced to life in prison, were both convicted of manslaughter,
Xinhua said.

The fight, at the Early Light Toy Factory in Shaoguan City, was incited by rumors that a group of
Muslim Uighurs had raped two Han Chinese women. It raged at a factory dormitory through the
early-morning hours of June 26, eventually leaving two Uighur men dead and, by some accounts, about
120 other people injured, most of them Uighurs.

The police said no woman was raped. The government said the fight was caused by a misunderstanding
involving a woman who accidentally ventured into a dormitory room of Uighur men. According to the
government's account, a Han Chinese man later posted reports of rape on the Internet, setting off a
bloody showdown.

On July 5, after reports of the brawl spread to Xinjiang, where most of the country's Uighur Muslims
live, demonstrations and riots broke out. There have been longstanding tensions there between Han
Chinese, who have flooded into the area, and native Uighurs.

The government said the Uighurs went on a murderous rampage\cite{rampage} against Han Chinese,
killing more than 200 people and leaving more than 1,000 wounded. Uighurs say the government has not
fully accounted for Uighur deaths in the region; and some Han Chinese have said the death toll among
their group was higher.

The government eventually put down the riots, but the region remains tense.

On Saturday, the Intermediate People's Court in Shaoguan City sentenced the men involved in the
brawl after a short trial, following their Sept.~23 indictments. In contrast to the accounts that
indicated scores were hurt, the court said three other people were severely injured and six slightly
injured.

The court said Saturday that Mr.~Xiao and his accomplices beat the Uighur men with iron bars and
``obstructed'' medical workers from treating the injured.

Of the 11 men sentenced Saturday, three were Uighur men, who were given five to six years in prison.

Xinhua said several other men also faced indictment for their roles in the brawl.

\section{Asia's Biggest Economies Promise Greater Cooperation}

\lettrine{T}{heir}\mycalendar{Oct.'09}{12} economies enmeshed as never before, leaders from China,
Japan and South Korea said over the weekend that they would explore the idea of a free-trade pact,
inching toward deeper regional integration.

Prime Minister Wen Jiabao of China and the leaders of Japan and South Korea declared that they were
``committed to the development of an East Asian community,'' and agreed to expand cooperation across
a wide range of issues, including climate change and sustainable growth.

While officials emphasized that the vision for a regional community was a long-term and inclusive
one, the shift toward a sharper focus on Asia was clear.

``Until now, we have tended to be too reliant on the United States,'' Prime Minister Yukio Hatoyama
of Japan said after a summit meeting Saturday with Mr.~Wen and President Lee Myung-bak of South
Korea.

``The Japan-U.S.~alliance remains important, but as a member of Asia, I would like to develop
policies that focus more on Asia,'' Mr.~Hatoyama said.

Since taking office last month, Mr.~Hatoyama has been floating the idea of creating an East Asian
Community, though he has been vague\cite{vague} about its goals and how it would work.

China, Japan and South Korea -- the region's biggest economies, with roughly 16 percent of global
gross domestic product -- would form the core of the envisioned community, he said.

The Japanese foreign minister, Katsuya Okada, said last week that with the region's differing
political systems and varying stages of growth, a European Union-like integration with a common
currency was unlikely anytime soon. Mr.~Okada has said the idea would be to group Japan, China,
South Korea, India, Australia and New Zealand and the 10 members of the Association of Southeast
Asian Nations, or ASEAN -- though it was not meant to exclude other nations.

Mr.~Wen and Mr.~Hatoyama also said the three leaders had agreed on the need to continue their
economic stimulus efforts, to avoid prematurely derailing the still feeble recovery from the global
financial crisis.

Mr.~Hatoyama appeared to have won favor with Japan's closest neighbors by emphasizing his desire to
assuage\cite{assuage} sensitivities over Tokyo's history of invasion and occupation in the region
before and during World War II.

``We have agreed we will seek common ground and shelve our differences,'' Mr.~Wen said after their
three-way meeting.

But plenty of issues remain, and meanwhile the region is moving toward wider economic integration
through bilateral and other, broader free-trade pacts.

Before considering a free-trade agreement or other economic partnerships, the three countries need
to sign an agreement on investment promoting a ``business friendly environment,'' said Kazuo Kodama,
press secretary to Mr.~Okada, the Japanese foreign minister.

Such efforts are hindered\cite{hinder} by differences over the exact conditions each country would
extend to the others, he noted. Even Japan's talks with South Korea on a two-country pact have
``somewhat stalled for the time being,'' Mr.~Kodama said.

\section{Training to Climb an Everest of Digital Data}

\lettrine{M}{OUNTAIN}\mycalendar{Oct.'09}{12} VIEW, Calif. -- It is a rare criticism of elite
American university students that they do not think big enough. But that is exactly the complaint
from some of the largest technology companies and the federal government.

At the heart of this criticism is data. Researchers and workers in fields as diverse as
bio-technology, astronomy and computer science will soon find themselves overwhelmed with
information. Better telescopes and genome sequencers are as much to blame for this data
glut\cite{glut} as are faster computers and bigger hard drives.

While consumers are just starting to comprehend the idea of buying external hard drives for the home
capable of storing a terabyte of data, computer scientists need to grapple with data sets thousands
of times as large and growing ever larger. (A single terabyte equals 1,000 gigabytes and could store
about 1,000 copies of the Encyclopedia Britannica.)

The next generation of computer scientists has to think in terms of what could be described as
Internet scale. Facebook, for example, uses more than 1 petabyte of storage space to manage its
users' 40 billion photos. (A petabyte is about 1,000 times as large as a terabyte, and could store
about 500 billion pages of text.)

It was not long ago that the notion of one company having anything close to 40 billion photos would
have seemed tough to fathom. Google, meanwhile, churns through 20 times that amount of information
every single day just running data analysis jobs. In short order, DNA sequencing systems too will
generate many petabytes of information a year.

``It sounds like science fiction, but soon enough, you'll hand a machine a strand of hair, and a DNA
sequence will come out the other side,'' said Jimmy Lin, an associate professor at the University of
Maryland, during a technology conference held here last week.

The big question is whether the person on the other side of that machine will have the wherewithal
to do something interesting with an almost limitless supply of genetic information.

At the moment, companies like I.B.M. and Google have their doubts.

For the most part, university students have used rather modest computing systems to support their
studies. They are learning to collect and manipulate information on personal computers or what are
known as clusters, where computer servers are cabled together to form a larger computer. But even
these machines fail to churn through enough data to really challenge and train a young mind meant to
ponder the mega-scale problems of tomorrow.

``If they imprint on these small systems, that becomes their frame of reference and what they're
always thinking about,'' said Jim Spohrer, a director at I.B.M.'s Almaden Research Center.

Two years ago, I.B.M. and Google set out to change the mindset at universities by giving students
broad access to some of the largest computers on the planet. The companies then outfitted the
computers with software that Internet companies use to tackle their toughest data analysis jobs.

And, rather than building a big computer at each university, the companies created a system that let
students and researchers tap into giant computers over the Internet.

This year, the National Science Foundation, a federal government agency, issued a vote of confidence
for the project by splitting \$5 million among 14 universities that want to teach their students how
to grapple with big data questions.

The types of projects the 14 universities have already tackled veer into the mind-bending. For
example, Andrew J.~Connolly, an associate professor at the University of Washington, has turned to
the high-powered computers to aid his work on the evolution of galaxies. Mr.~Connolly works with
data gathered by large telescopes that inch their way across the sky taking pictures of various
objects.

The largest public database of such images available today comes from the Sloan Digital Sky Survey,
which has about 80 terabytes of data, according to Mr.~Connolly. A new system called the Large
Synoptic Survey Telescope is set to take more detailed images of larger chunks of the sky and
produce about 30 terabytes of data each night. Mr.~Connolly's graduate students have been set to
work trying to figure out ways of coping with this much information.

Purdue, meanwhile, looks to carry out techniques used to map the interactions between people in
social networks into the biological realm. Researchers are creating complex diagrams that illuminate
the links between chemical reactions taking place in cells.

A similar effort at the University of California, Santa Barbara, centers on making a simple software
interface -- akin to the Google search bar -- that will let researchers examine huge biological data
sets for answers to specific queries.

Mr.~Lin has encouraged his students to illuminate data with the help of Hadoop, an open-source
software package that companies like Facebook and Yahoo use to split vast amounts of information
into more manageable chunks.

One of these projects included a deep dive into the reams of documents released after the
government's probe into Enron, to create an analysis system that could identify how one employee's
internal communications had been connected to those from other employees and who had originated a
specific decision.

Mr.~Lin shares the opinion of numerous other researchers that learning these types of analysis
techniques will be vital for students in the coming years.

``Science these days has basically turned into a data-management problem,'' Mr.~Lin said.

By donating their computing wares to the universities, Google and I.B.M. hope to train a new breed
of engineers and scientists to think in Internet scale. Of course, it's not all good will backing
these gestures. I.B.M. is looking for big data experts who can complement its consulting in areas
like health care and financial services. It has already started working with customers to put
together analytics systems built on top of Hadoop. Meanwhile, Google promotes just about anything
that creates more information to index and search.

Nonetheless, the universities and the government benefit from I.B.M. and Google providing access to
big data sets for experiments, simpler software and their computing wares.

``Historically, it has been tough to get the type of data these researchers need out of industry,''
said James C.~French, a research director at the National Science Foundation. ``But we're at this
point where a biologist needs to see these types of volumes of information to begin to think about
what is possible in terms of commercial applications.''

\section{Debate Rages About Whether Leno's Show Is Good for NBC}

\lettrine{T}{wo}\mycalendar{Oct.'09}{12} weeks into a new season, the talk of television is the Leno
effect -- and whether it is hazardous to NBC's health.

It is not just a question of how the new ``Jay Leno Show'' itself is faring in the ratings, but also
what the show's occupation of the 10 p.m. hour on NBC means to the network as a whole.

As Shari Anne Brill, the senior vice president and director of program analysis for the advertising
agency Carat, put it, ``It's really looking like dominoes.''

The dominoes in question are the other parts of NBC's schedule affected by the network's decision to
relocate its late-night star, Mr.~Leno, to prime time. Even though, as NBC executives point out, it
is early in this experiment, signs of potential collateral damage have already emerged.

Shows seem to have suffered because they have been displaced to new time periods, like ``Law \&
Order SVU,'' which was the leading drama when it played at 10 p.m. on Tuesdays, but now is finishing
last after moving to 9 on Wednesdays.

Late newscasts on local stations affiliated\cite{affiliate} with NBC are reporting significant
ratings declines, at least partly because of a ratings drop-off in the 10:30 half-hour that precedes
them.

And the late-night programs, led by ``The Tonight Show,'' that have been a perennial source of
strength for NBC are no longer the automatic winners against their CBS competition.

Conan O'Brien, the new ``Tonight'' host, has been swamped recently by the tidal wave of publicity
surrounding his CBS rival, David Letterman; but he also has been inheriting much smaller audiences
than ever before in ``Tonight's'' history.

That means NBC's second act in late night, Jimmy Fallon, is also getting a much weaker lead-in than
Mr.~O'Brien did in the same hour a year ago, and as a result is falling behind his CBS competitor,
Craig Ferguson, in audience totals.

NBC rightly points out that both its late-night stars remain more popular with younger viewers who
are more valuable in selling to advertisers, but in the past NBC's late-night hours were dominant
across the board, not just with narrower audience segments.

And all of this is playing out against a backdrop of reports that General Electric is in talks with
Comcast to sell NBC Universal.

Looked at in isolation, Mr.~Leno has been doing everything NBC expected of him.

His ratings, after a big first week, have leveled off to about five million viewers a night (though
some nights have been much lower) with a 1.5 to 2 rating in the category NBC identifies as
all-important, viewers ages 18 to 49, the group many advertisers want to reach.

Though most 10 p.m. shows with those kinds of numbers get canceled, NBC has said from the beginning
that it could accept much lower ratings because of the enormous cost savings of Mr.~Leno's show
versus expensive hourlong scripted dramas. The network guaranteed advertisers that it would average
only a 1.5 rating.

The network's performance over all has not shown signs of a comeback. While somewhat propped up so
far by professional football on Sunday night, NBC has not added any standout new shows. The new
drama ``Trauma'' has already faltered, and a promising new comedy, ``Community,'' struggled last
week when it was moved to a new 8 p.m. time period.

NBC has only two real points of strength now, two hours worth of the reality show ``The Biggest
Loser'' on Tuesday and the comedy ``The Office'' on Thursday.

NBC has also emphasized that Mr.~Leno needs to be judged over the full year because he will be
offering many more original weeks of shows than his competitors.

But for some, the judgment is already clear-cut. Producers of shows that have in the past, and could
in the future, fill the 10 p.m. hour on several networks are using words like ``complete calamity''
and ``utter disaster'' to describe the current state of NBC -- though they are using the words while
requesting anonymity because of the potential to be in business with NBC in the future.

One producer of several hits lamented the overall absence of a 10 p.m. opportunity for new dramas,
saying NBC was formerly the place where the most innovative dramas on television -- from ``Hill
Street Blues'' to ``E.R.'' -- found a home.

One recent example was the new police show ``Southland,'' which was NBC's best-reviewed drama in
years. In the past the network might have waited for it to build an audience based on its quality;
but last week NBC announced it was canceling the show.

``Southland'' was clearly intended as a 10 p.m. entry. John Wells, the longtime executive producer
of ``E.R.'' who held the same position on ``Southland,'' issued a statement saying, ``I'm
disappointed that NBC no longer has time periods available to support that kind of critically
acclaimed series that was for so many years a hallmark of their success.''

If producers have reason to be dismayed, owners of NBC's affiliated stations may be expected to be
in open revolt. Among the top 15 cities in the country, ratings for the late news -- a prime source
of revenue for local television -- are down 10 to 30 percent

But so far the owners seem to be holding their tongues -- along with their breath. ``You don't make
decisions based on a week or two,'' said Michael Fiorile, the vice chairman of the Dispatch
Broadcasting Group, which owns the NBC affiliate in Indianapolis. ``Six months from now we'll take a
look at the trends.''

Media buyers like Ms.~Brill are saying the early results are really no surprise. ``It's exactly what
I predicted,'' Ms.~Brill said, adding of the decision to move Mr.~Leno to 10, ``it was never a
ratings decision. It was a money decision.''

It still seems to be. Jeff Gaspin, the chairman of NBC Universal Entertainment, said he was certain
of one aspect of the Leno move. ``We'll make money at 10 o'clock this year, I guarantee.''

The rationale for the move of Mr.~Leno was simple: the network could not endure his likely move to
ABC, where he would have created a new late-night program and undermined the strength of ``The
``Tonight Show.''

But NBC is justifying the move by citing both the savings Mr.~Leno's show represents over expensive
10 p.m. dramas and the apparent disintegration of the 10 p.m. hour across the board. Mr.~Gaspin
repeated NBC's conclusion that hits cannot be established at 10 anymore, largely because the hour is
dominated by viewers playing back recorded shows on digital video recorders.

``Look at how ABC is doing at 10 against Jay,'' Mr.~Gaspin said.

Indeed, ABC's performance is certainly providing some cover for NBC's move at 10. Mr.~Leno is
already faring as well or better than two new ABC dramas, ``The Forgotten'' and ``Eastwick,'' and he
is not far behind a third, ``Castle.'' All those shows cost three times as much or more per episode
as Mr.~Leno's show.

Mr.~Gaspin argued that NBC is not abandoning quality drama and cited recent deals for future shows,
including outbidding the other networks for a spy drama from J.~J.~Abrams (``Lost'') and an American
version of the British police classic ``Prime Suspect.''

``Maybe we made some wrong choices with shows this season, but we are still investing in
programming,'' Mr.~Gaspin said.

NBC is not engaging in any speculation about trends because it is arguing it is simply too early to
read them. ``We have to play for the long haul,'' Mr.~Gaspin said. It also might help not to keep
``making proclamations that we're doing this for cost reasons,'' he said.

``Jay is doing fine,'' Mr.~Gaspin said. ``Conan is doing what we expected him to do.'' He added,
``We're going to look at our average over the full year.''

\section{Fox's Volley\cite{volley} With Obama Intensifying}

\lettrine{A}{ttacking}\mycalendar{Oct.'09}{12} the news media is a time-honored White House tactic
but to an unusual degree, the Obama administration has narrowed its sights to one specific
organization, the Fox News Channel, calling it, in essence, part of the political opposition.

``We're going to treat them the way we would treat an opponent,'' said Anita Dunn, the White House
communications director, in a telephone interview on Sunday. ``As they are undertaking a war against
Barack Obama and the White House, we don't need to pretend that this is the way that legitimate news
organizations behave.''

Her comments are only the latest in the volatile exchange between the administration and the
top-rated network, which is owned by the News Corporation, controlled by Rupert Murdoch. Last month,
Roger Ailes, the chairman of Fox News, and David Axelrod, a senior adviser to President Obama, met
for coffee in New York, in what Politico, which last week broke that news, labeled a ``Fox summit.''

While neither party has said what was discussed, some have speculated that a truce, or at least an
adjustment in tone, was at issue. (Mr.~Ailes and Mr.~Obama reportedly reached a temporary accord
after a meeting in mid-2008.) But shots are still being fired, which animates the idea that both
sides see benefits in the feud.

Fox seems to relish\cite{relish} the controversy.

``Instead of governing, the White House continues to be in campaign mode, and Fox News is the target
of their attack mentality,'' Michael Clemente, the channel's senior vice president for news, said in
a statement on Sunday. ``Perhaps the energy would be better spent on the critical issues that voters
are worried about.''

Fox's senior vice president for programming, Bill Shine, says of the criticism from the White House,
``Every time they do it, our ratings go up.'' Mr.~Obama's first year is on track to be the Fox News
Channel's highest rated.

One Fox executive said that the jabs by the White House could solidify the network's audience base
and recalled that Mr.~Ailes had remarked internally: ``Don't pick a fight with people who like to
fight.'' The executive asked not to be named while discussing internal conversations.

Certainly, Fox continues to aggressively bolster its on-air talent, most recently with the hiring of
John Stossel, the libertarian investigative journalist from ABC News, for its spin-off channel, Fox
Business. The business channel is also keen on another administration critic, Lou Dobbs, who met for
dinner with Mr.~Ailes last month, according to two people with direct knowledge of the meeting.

The shift for Fox News -- the favorite network of the Bush administration, now the least favored one
of the Obama administration -- has financial implications for the News Corporation, especially given
the network's status as a growth engine in a perilous time for media companies.

Fox's programs have drawn record numbers of viewers this year. Through last week, Fox averaged 1.2
million viewers at any given time this year, up from one million viewers through the same time last
year. Previously, the channel peaked in 2003, the year the Iraq war started, with nearly 1.1 million
viewers.

But controversial comments by the host Glenn Beck have also prompted an ad boycott. And the
perception of Fox News as an opposition party has also affected its news correspondents, including
Major Garrett, its chief White House correspondent, who Ms.~Dunn says is a fair reporter.
Mr.~Garrett and other Fox correspondents have been directed by their bosses not to appear on the
channel's most opinionated programs.

Still, Paul Rittenberg, who oversees ad sales for Fox, said the channel existed in a climate where
viewers choose cable news channels based on affinity. His channel, he said, stresses in its pitch to
advertisers that ``people who watch Fox News believe it's the home team.''

To many Democrats, of course, the ``home team'' is conservative, a view only compounded by Fox's at
times skeptical coverage of Mr.~Obama this year.

``I've got one television station that is entirely devoted to attacking my administration,'' he said
in June, though he did not mention Fox by name. He added, ``You'd be hard pressed if you watched the
entire day to find a positive story about me on that front.''

The White House has limited administration members' appearances on the network in recent weeks. In
mid-September, when the White House booked Mr.~Obama on a round robin of Sunday morning talk shows,
it skipped Fox and called it an ``ideological outlet,'' leading the ``Fox News Sunday'' anchor Chris
Wallace to appear on Bill O'Reilly's prime-time show and call the administration ``the biggest bunch
of crybabies I have dealt with in my 30 years in Washington.''

Ms.~Dunn called that remark juvenile and stressed that administration officials would still talk to
Fox, and that Mr.~Obama was likely to be interviewed on the network in the future. But, she added,
``we're not going to legitimize them as a news organization.''

In an interview, Mr.~Clemente suggested that there was an element of ``shoot the messenger'' in the
back and forth. ``Sometimes it's actually helpful to have an organization or a person that you can
go up against for whatever reason,'' he said.

Fox argues that its news hours -- 9 a.m. to 4 p.m. and 6 to 8 p.m. on weekdays -- are objective. The
channel has taken pains recently to highlight its news programs, including the two hours led by
Shepard Smith, its chief news anchor. And its daytime newscasts draw more viewers than CNN or
MSNBC's prime-time programs.

``The average consumer certainly knows the difference between the A section of the newspaper and the
editorial page,'' Mr.~Clemente said.

The White House rejects the news and editorial page comparison, and officials there can rattle off
any number of perceived offenses. They date to the month before Mr.~Obama formally started his
presidential campaign, when one of the network's morning hosts falsely claimed that he had attended
a madrassa, an Islamic school. (The incident happened on what Fox calls an entertainment show, ``Fox
and Friends''; the mistake was corrected on the air later.)

More recently, Fox hosts have promoted tea party rallies against big government and steered
attention toward a number of White House czar appointments. Mr.~Beck, in particular, was credited
with forcing Van Jones, a low-level White House adviser for environmental jobs, to resign last
month. Mr.~Beck devoted numerous segments to Mr.~Jones and called him a ``communist-anarchist
radical.''

``If it wasn't for Fox or talk radio, we'd be done as a republic,'' Mr.~Beck said in the wake of the
resignation.

Mr.~Beck, whose 5 p.m. program consistently draws three million viewers, is a ``cultural phenomenon
now,'' Mr.~Shine said. But this success has come at a price: he is the source of considerable
discomfort for Fox's journalists, especially for false statements on his program. In August, for
instance, Mr.~Beck claimed that Mr.~Garrett was ``never called on'' at White House press briefings,
but Mr.~Garrett had asked a question that day.

Weeks earlier, Mr.~Beck labeled Mr.~Obama a racist, leading to an advertising boycott by
ColorOfChange.org, an advocacy group that Mr.~Jones helped found. Dozens of advertisers have
distanced themselves from Mr.~Beck's show, causing headaches for Mr.~Rittenberg's advertising team,
although he said Fox ``hasn't lost a dime'' because the ads were moved to different hours.

Fox has made the channel's tensions with the White House a story. In August, the network's top-rated
host, Mr.~O'Reilly, dispatched one of his opinion program's producers to ask why the administration
seemed ``so thin-skinned'' at a White House briefing. The deputy press secretary disagreed, and said
that Mr.~O'Reilly had interviewed Mr.~Obama during his candidacy last year. The administration's
aggressive stance suggests that it does not view Fox's audience as one that can be persuaded. During
the presidential campaign, Ms.~Dunn said, it booked campaign representatives on Fox to try to reach
undecided voters, but by mid-October, the campaign had mostly withdrawn them from the channel's
programs.

``It was beyond diminishing returns,'' she said. ``It was no returns.''

\section{The Gift of Life, and Its Price}

\lettrine{S}{cary}\mycalendar{Oct.'09}{12}. Like aliens. That is how Kerry Mastera remembers her
twins, Max and Wes, in the traumatic days after they were born nine weeks early. Machines forced air
into the infants' lungs, pushing their tiny chests up and down in artificial heaves. Tubes delivered
nourishment. They were so small her husband's wedding band fit around an entire baby foot.

Having a family had been an elusive goal for Jeff and Kerry Mastera, a blur of more than two years,
dozens of doctor visits and four tries with a procedure called intrauterine insemination, all
failures. In one year, the Masteras spent 23 percent of their income on fertility treatments.

The couple had nearly given up, but last year they decided to try once more, this time through
in-vitro fertilization. Pregnancy quickly followed, as did the Mastera boys, who arrived at the
Swedish Medical Center in Denver on Feb.~16 at 3 pounds, 1 ounce apiece. Kept alive in a neonatal
intensive care unit, Max remained in the hospital 43 days; Wes came home in 51.

By the time it was over, medical bills for the boys exceeded \$1.2 million.

Eight months later, the extraordinary effort seems worth it to the Masteras, who live in Aurora,
Colo. The babies are thriving and developing their own personalities -- Wes, the noisy and
demanding; Max, the quiet and serious. Like many other twins conceived through in-vitro
fertilization, the Mastera boys will go down in the record books as a success -- both for the
fertility clinic that helped create them and the neonatologists who nursed them to health.

But an exploration of the fertility industry reveals that the success comes with a price. While IVF
creates thousands of new families a year, an increasing number of the newborns are twins, and they
carry special risks often overlooked in the desire to produce babies.

While most twins go home without serious complications, government statistics show that 60 percent
of them are born prematurely. That increases their chances of death in the first few days of life,
as well as other problems including mental retardation, eye and ear impairments and learning
disabilities. And women carrying twins are at greater risk of pregnancy complications.

In fact, leaders of the fertility industry and government health officials say that twins are a risk
that should be avoided in fertility treatments. But they also acknowledge that they have had
difficulty curtailing the trend.

Many fertility doctors routinely ignore their industry's own guidelines, which encourage the use of
single embryos during the in-vitro fertilization procedure, according to interviews and industry
data. Some doctors say that powerful financial incentives hold sway in a competitive marketplace.
Placing extra embryos in a woman's womb increases the chances that one will take. The resulting
babies and word of mouth can be the best way of luring new business.

Doctors are also often under pressure from patients eager for children, who have incentives to
gamble as well. Frequently, they have come to IVF as a last resort after years of other treatments,
are paying out of pocket, and are anxious to be successful on the first try. And many do not fully
understand the risks.

Dr.~William E.~Gibbons, incoming president of the American Society for Reproductive Medicine, said
his organization was concerned about the risks of twin pregnancies and would issue new guidelines at
a meeting next week to further discourage multiple births. ``People should be made aware of the
concerns that we think twins are not a good outcome,'' Dr.~Gibbons said.

The industry creates preterm infants with in-vitro and other fertility treatments even as government
and nonprofit groups work to fight the nation's 12.7 percent rate of prematurity, regarded as a
major national health care problem.

While IVF multiples are typically the children of affluent women, much of the effort at reducing
premature birth has been focused on prevention and prenatal care for low-income women. A study
released last week by the March of Dimes cited fertility treatments as one of the main reasons for a
36 percent increase in prematurity in the last 25 years.

The government estimates that caring for premature infants costs \$26 billion a year, including \$1
billion for IVF babies, expenses that eventually get passed through the system and on to businesses
and consumers.

The unusual birth of octuplets in California in January notwithstanding, the American Society for
Reproductive Medicine and its affiliate, the Society for Assisted Reproductive Technology, have
succeeded in reducing the number of larger multiple births from in-vitro fertilization over the last
several years.

The two medical organizations and the federal Centers for Disease Control and Prevention have been
promoting the use of single embryos in many cases to reduce the number of twins. But that has not
translated into major action at the 483 fertility clinics across the country. The latest figures
from the industry show that women under 35, the group most likely to get pregnant from the
treatments, choose to use single embryos in only 4.5 percent of in-vitro rounds.

``You can't convince a couple that having twins is a bad thing,'' said Dr.~Maurizio Macaluso, who
runs the C.D.C.'s women's health and fertility branch. ``That's a major communication problem.''

In 2006, a record 137,085 twins were born in the United States, double the number in 1980. Of that
total, 23,284 were the result of IVF, according to government statistics. The number does not
include twins born as a result of other fertility treatments.

Most fertility doctors acknowledge the potential problems with twins, whether conceived naturally or
through fertility treatments. But many say that the good done by their industry -- creating new
families -- outweighs the bad, and that twins are not such a risky bet because most are healthy.

``At the end of the day, when you dissect the statistics out, our patients are interested in
establishing a family and a pregnancy,'' said Dr.~Michael Swanson, a Colorado fertility doctor who
treated Ms.~Mastera.

It is a tricky cost-benefit analysis, however, and one that potentially involves the worst kind of
collateral damage, the type that figures in the nightmares of expectant parents.

Erin and Scott Hare of Houston lost their twin daughter, conceived through in-vitro fertilization.
Her surviving brother, Carter, who was born at just over 24 weeks, is doing well but needs therapy
for lingering problems.

George and Narine Nazaretyan of Van Nuys, Calif., have twin daughters conceived through IVF. One of
the girls, Natlie, has a severe case of cerebral palsy, which occurs four to six times as frequently
in twins as in single babies.

Cutting down preterm births from IVF would be an easy way to make a small dent in reducing the
nation's prematurity problem. Dr.~Macaluso calls them ``low-hanging fruit'' -- a partial solution
that is within reach.

But there is concern among public health officials that the problem may instead grow as fertility
treatments become available to more people.

``In the past few years, we have felt increasingly uncomfortable because we feel like we are sitting
on the tip of the iceberg,'' Dr.~Macaluso said.

Pressure for Success

In the competitive marketplace for fertility medicine, success rates are the metric by which
in-vitro clinics thrive or fail. The rates -- meaning the chances of producing a baby at each clinic
-- are published by the C.D.C. and are widely used by couples to choose a doctor.

Congress passed a law in 1992 that required the data to be disclosed after some clinics were
exaggerating\cite{exaggerate} their numbers to lure potential clients. But there is evidence that
the law has had the unintended effect of pressuring doctors to transfer multiple embryos to maximize
their success rates.

``If a person does not have as high a pregnancy rate as his neighboring competitor, they're going to
lose those prospective patients,'' said Dr.~David Kreiner, medical director of East Coast Fertility,
a network of fertility centers based in Plainview, N.Y.

A busy fertility clinic can be extraordinarily lucrative, generating millions of dollars a year. And
fertility doctors can take on godlike status in their communities for delivering their priceless
commodities.

Knowing that prospective parents can easily seek IVF elsewhere, doctors give them unusual
autonomy\cite{autonomy} in deciding how many embryos to transfer.

For many in-vitro patients, the high cost of treatment is often a factor in making that decision.
When Ms.~Mastera had the procedure at Conceptions Reproductive Associates of Littleton, Colo., in
September 2008, she was a candidate for single embryo transfer under industry guidelines. At age 32,
and having never before undergone in-vitro fertilization, her chances for pregnancy were excellent.

Conceiving a child had become an obsession for the couple, who had met in 2000 while working as
customer service representatives for Nextel. By 2006 they were married, living in a modest
split-level home in Aurora and ready to start a family.

``I just told Jeff one day, 'I think we should get off birth control; I'm ready,' '' said
Ms.~Mastera, who was 30 at the time. ``He was like, 'O.K., let's do it.' ''

After trying unsuccessfully for more than a year, the couple consulted two fertility specialists in
2007 and spent more than \$15,000 on the four rounds of intrauterine insemination.

``You're an emotional basket case because you're on these hormones,'' Ms.~Mastera said. ``We'd be
constantly worried about money. Like constantly. How are we going to pay our phone bill this month?
Or our mortgage? Because we're having to pay for all these fertility treatments.''

In-vitro fertilization was the next logical step, but the price tag was even more daunting.
Depending on the clinic, its location and the extra services included in the treatment, the
procedure can run \$12,000 a cycle to more than \$25,000.

In 2008, the Masteras consulted their third fertility expert, Dr.~Swanson at Conceptions. They chose
him partly based on his clinic's high success rates as published by the Centers for Disease Control
and Prevention.

Ms.~Mastera's insurance would cover about \$8,000 for the procedure and drugs. Frequently, insurance
does not cover anything. Almost \$18,000 in clinic fees and other costs remained for the couple, who
proceeded to cash out their 401(k) and money market accounts and put the remaining balance on a
low-interest credit card.

Like many families, the Masteras could not afford a second cycle. So when the couple was given a
choice by their doctor of implanting one or two embryos, they decided to increase their chances with
two.

``This was our Hail Mary pass,'' said Ms.~Mastera, now 33. ``We thought, let's just do it. At the
time, it was like, twins, they can be fun. They are fun, but holy cow.''

Some public health experts are frustrated by the disconnect between the medical risk of twins and
society's perception.

Twins are celebrities and celebrities have twins. The pop culture media goes into overdrive when
stars like Angelina Jolie and Jennifer Lopez give birth to twins.

``When they have their twins, it's a very acceptable thing,'' said Dr.~Frank L.~Mannino, medical
director of the neonatal intensive care unit at the University of California, San Diego, Medical
Center. The center expanded the unit last year, citing increased demand for the services, largely
because of the use of fertility treatments.

Dr.~Alan R.~Fleischman of the March of Dimes, which has begun distributing literature discussing
single embryo transfer, said a result of playing down the risks of multiple births is that many
women undergoing in-vitro prefer to have twins.

``It's not just a matter of hoping that if you put two embryos in, you'll get one baby,''
Dr.~Fleischman said. ``There are many women who actually want to have two children and would like to
have their full family with one pregnancy.''

The Risks

The problems began at week 24 of Ms.~Mastera's pregnancy.

The human uterus is designed to carry one fetus. An ultrasound on Jan.~7 showed that the extra
burden of twins was placing pressure on her cervix, causing the amniotic sac surrounding Max to
begin pushing through.

``I was really stunned; I couldn't believe what was going on,'' said Ms.~Mastera, who had not been
having symptoms. ``They came in and said, 'You are going into surgery right now.' '' As doctors
tried to comfort her, Ms.~Mastera said, she began crying uncontrollably.

In a procedure called a cerclage, doctors pushed the sac back into her uterus, then stitched up her
cervix.

``The doctor said that if she were pregnant with one, she would have been just fine,'' said her
husband, Jeff, 31. ``The sheer weight of the babies was just pushing them lower and lower.''

After the procedure, Ms.~Mastera was ordered to remain in bed at home. Doctors hoped she would make
it close to her April due date. But on the night of Feb.~15, Ms.~Mastera's water broke.

Her boys were born the next day, and were not the healthy babies she had dreamed of. ``They wheeled
me in, after I had recovered for an hour,'' she said. ``And it was weird. They looked really frail
and unhealthy.''

Ms.~Mastera said she was aware of the risk of prematurity with twins and had discussed the issue
with Dr.~Swanson before having the in-vitro procedure. Informed-consent documents given to patients
by fertility doctors normally detail the increased risk of twins. But even with those increased
risks, the actual number of serious outcomes like fetal death or brain damage is small.

While the average single pregnancy in the United States now last about 39 weeks, the average twins
are born at just over 35 weeks, , according to the Institute of Medicine. And there is emerging
evidence that babies born at that time can develop learning problems.

``There's increasing evidence that the 34-, 35-, 36-weekers, the bigger prematures -- and there are
large number of those -- have significant problems with learning later on, even though their
mortality isn't high'' said Dr.~Richard E.~Behrman, a pediatrician and former vice president and
dean of medical affairs at Case Western Reserve University who led the institute's committee on
preterm birth in 2006. ``And that's not appreciated in the obstetrical and lay community.''

According to one federal study, about 30 percent of all twins end up in a neonatal intensive care
unit. Twins are eight times as likely as single babies to be born at very low birth weight --
defined as under 3 pounds, 4 ounces. These are the babies who often need the longest care and face
the biggest problems. Dr.~Macaluso calls them ``million-dollar babies.''

Carter Hare, the son of the Texas couple, was one of them. His birth announcement in December 2006
gave a hint of the trouble: ``Carter is born! 24 weeks 4 days.''

His health was touch and go. ``They gave us a very grim outlook,'' his mother, Erin Hare, said
recently.

Doctors were aware early on that Ms.~Hare's pregnancy might be complicated. Three months before
becoming pregnant with in-vitro twins she had miscarried a single in-vitro baby because of a
condition called incompetent cervix, a problem similar to Ms.~Mastera's.

After becoming pregnant with the twins, Ms.~Hare had her cervix stitched closed to keep her babies
in place and remained hospitalized for much of the pregnancy.

At 20 weeks, an ultrasound showed that her baby girl's heart had stopped beating, an apparent result
of a compressed umbilical cord. ``It was hard,'' Ms.~Hare said. ``I actually at that point stopped
accepting visitors. I put a 'Do Not Disturb' sign on the door.''

In her solitude, Ms.~Hare was fighting to continue carrying both her son and the lifeless body of
her daughter, whose sac began bulging through her vagina, threatening the entire pregnancy. Doctors
performed a second cervical surgery, continuing efforts to keep her son from being born before 24
weeks gestation, regarded as the point of viability.

He was born four days beyond that milestone, 12 inches long and weighing one pound, 12 ounces.
Carter spent 102 days in the neonatal intensive care unit at Presbyterian Hospital of Plano, near
Dallas.

This year, Carter entered preschool.

``He's really a little miracle baby,'' said Ms.~Hare, a tax accountant. The family has since moved
to Houston. Despite initial heart and eye problems, Carter did not require surgery. He receives
therapy for sensory problems that sometimes develop in premature babies.

``He wouldn't like being in a pool,'' Ms.~Hare said. ``He would scream. His senses are just off. He
has an extremely high tolerance for pain, then sometimes he's real sensitive. The whole sensory
process, when they develop outside the womb, it just doesn't develop. It takes longer.''

Despite her troubled pregnancy, Ms.~Hare tried another round of in-vitro after Carter was born, this
time undergoing a more invasive surgery that stitches the lower part of the uterus and the upper
cervix together and requires an abdominal incision.

``A lot of people thought I was crazy, but I didn't did feel our family was complete yet,'' Ms.~Hare
said.

On April 29, 2008, a daughter, Lauren, was born by Caesarean section at 36 weeks, 4 days.

The Price Tag

In March, the United States Chamber of Commerce and the March of Dimes held a luncheon in Washington
to discuss preterm babies. ``The human costs are staggering,'' Dr.~Steven K.~Galson, then the acting
surgeon general, told the group. ``The medical costs are staggering. That's why we're here.''

``Today you're going to hear that preterm birth is not just a significant public health issue,''
Dr.~Galson said, ``but that it also impacts businesses and employer health plans.''

The hospitalization and doctor's care for Ms.~Hare and her son exceeded \$1 million. Most of that,
about \$750,000 to \$800,000, was for Carter. The bill was picked up by the self-funded health plan
of the Trammell Crow Company, the Dallas real estate investment company where Ms.~Hare worked.

``The following quarter during the earnings release, somebody asked why there was a sharp increase
in medical costs,'' Ms.~Hare said. No one identified her, but Ms.~Hare knew that her family had
contributed heavily.

In Atlanta, the Centers for Disease Control and Prevention hired an economist to predict what would
happen if single embryo transfer were used in a large number of IVF cases.

Dr.~Macaluso, the C.D.C. reproductive health official, estimates the patients, businesses and
insurance providers would save more than \$500 million annually, even taking into consideration the
cost of extra in-vitro rounds, by lowering neonatal intensive care, special education and other
costs of premature babies.

To reduce the number of twins, some clinics are experimenting with programs that provide IVF with
single embryo transfer with free freezing of extra embryos and free transfer of frozen embryos if
the first try does not work. Others are working to develop ways to identify the specific
characteristics of a single embryo that will turn into a healthy baby.

Dr.~G. David Adamson, a former president of the American Society for Reproductive Medicine, has
advocated reducing the twins rate but says that the value of a baby should also be taken into
account when discussing single embryo transfer.

``How many healthy babies do we want and what's the risk and cost to the individual and society of
having a baby with problems,'' Dr.~Adamson asked. ``It's a complex question, and varies from patient
to patient.''

For the Masteras, the pain and anguish is fading as they watch their boys grow.

Although they are at risk for developmental problems because of their prematurity, the boys are
meeting milestones. Max can now sit up on his own, and Wes crawled for the first time last week.

Ms.~Mastera still feels guilty, worried that she did something to make her boys premature. ``I don't
know if I'll ever forget how horrible I felt about them coming early,'' Ms.~Mastera said. ``And I
don't know if I'll ever forget how we struggled.''

\section{A New Way to Inhale, Not Inject, Insulin}

\lettrine{P}{eople}\mycalendar{Oct.'09}{12} with diabetes often inject themselves with insulin at
mealtime to help control their blood sugar levels. But a new, palm-size device may let them
discretely inhale a dose of insulin instead of using a needle.

A small inhaler and insulin powder created by the MannKind Corporation, a drug developer in
Valencia, Calif., are before the Food and Drug Administration for marketing approval.

The insulin powder, called Afresa, is inhaled into the lungs, dissolves there and then travels into
the bloodstream, says Matthew J.~Pfeffer, chief financial officer at MannKind.

Using insulin or other drugs to control blood sugar helps diabetics avoid serious complications,
including heart disease, kidney failure, blindness and nerve damage.

The use of insulin in an inhaled form is not new. It was introduced by Pfizer with a product called
Exubera in 2006. But the inhaler used with Exubera was large and awkward, some critics contended,
which may have been a reason the product didn't become popular. It was withdrawn less than two years
after federal approval.

But MannKind may have better prospects because of its smaller inhaler and fast-acting insulin.
Mr.~Pfeffer says the MannKind inhaler fits neatly in one hand, and a second-generation versions the
company is using in clinical trials is even smaller, the size of a whistle.

Patients put insulin doses -- pre-packaged in cartridges -- into the inhaler and turn the mouthpiece
to release the insulin. The inhaler uses no electricity or compressed gas. ``The patient's breathing
action does the job,'' Mr.~Pfeffer said. ``The airflow through the cartridge allows the powder to be
inhaled.''

The system now before the F.D.A. is for adults with Type 1 diabetes, which often begins in
childhood, and Type 2 diabetes, which typically occurs when people are older.

Simos Simeonidis, a senior biotechnology analyst at the New York investment bank Rodman \& Renshaw,
who wrote a report on MannKind, said he expected its system to be available next year, if the F.D.A.
approved. (Dr.~Simeonidis has no stock in MannKind, but Rodman \& Renshaw has provided investment
banking services for it.)

Leonid Poretsky, chief of endocrinology at Beth Israel Medical Center in New York and director of
the Gerald J.~Friedman Diabetes Institute there, said that the MannKind system will face many
problems even if it is approved.

``Injections today are essentially painless,'' he said of the short, thin needles that are commonly
used to inject insulin. And you don't necesssarily ``have to draw from a bottle into a syringe.
Injections work so well that the advantages of a new route like this are unclear.''

Dr.~Poretsky was also concerned about using the lungs to transport drugs. ``It's possible for people
to stay on insulin for decades,'' he said. ``The whole issue of exposing the lungs to insulin for a
long period of time has to be examined carefully.''

Dr.~Gerald Bernstein, a New York-based endocrinologist who is a former president of the American
Diabetes Association, agreed that the long-term use of inhalable insulin might carry risks for some
patients. Dr Bernstein is vice president of the Generex Biotechnology Corporation, which is
developing an insulin delivered through the lining of the mouth.

``It's counterintuitive to use the fragile cells of the alveoli,'' the tiny air sacs within the
lungs, ``to get insulin to the bloodstream,'' he said. ``The lungs were developed to transport
gases, not proteins.''

Mr.~Pfeffer of MannKind said that the company's clinical data included no signs of damage to lungs.

DR. DAVID M.~NATHAN, a professor of medicine at Harvard Medical School and director of the
Massachusetts General Hospital Diabetes Center, said that even if safety issues were addressed,
there could be other long-term problems with Afresa. He questioned whether MannKind's inhalable
product could achieve the same level of blood sugar control as that obtained with injections or
insulin pumps.

``Insulin inhalers are tricky to use,'' he said. ``The amount you breathe in can be variable.'' But
he acknowledged that injections aren't perfect and that many people would prefer to get their
insulin in a more convenient way.

Mr.~Pfeffer said trials by MannKind had shown a high consistency in doses by patients using the
inhalers.

If all goes well at the F.D.A., Mr.~Pfeffer says he thinks that MannKind will eventually go to
market with the smaller dispenser now in trials. ``The expectation,'' he said, ``is that we will
launch with the whistle-sized inhaler'' -- which, by the way, was given the internal company name
``Dreamboat'' by its design team.

``Perhaps we'll consider a different name, though,'' he said.

\section{A Cruise, a Terminal Illness and a Technicality}

\lettrine{O}{ne}\mycalendar{Oct.'09}{12} secret to a good long winning streak is picking opponents
you're sure you can beat and avoiding those who look like trouble. The strategy goes a long way
toward explaining why, in the short life of this column, the Haggler has prevailed on behalf of
every consumer whose cause he's undertaken.

Until now. Spoiler alert: the Haggler is about to lose.

Q. Last year, I booked an eight-day trip to Europe for me and my brother, Stephen, who has been
diagnosed with terminal cancer. He doesn't have health insurance, so to raise money for both the
trip and his treatment, I helped organize a rock concert here in Austin. The lineup included Kathy
Valentine, a former member of the Go-Go's, and Bobby Whitlock, a former Eric Clapton sideman.
Ultimately, we raised \$38,000. We spent about \$7,000 on the trip, a price that included \$264
travel insurance, recommended by the travel agent in case Stephen had to cancel.

A few days before the flight, he had a recurrence of a gastrointestinal ailment that is a side
effect of his illness. We canceled our trip, but after we submitted paperwork for the insurance
claim, the company that sold us the policy, CSA Travel Protection, said it would not return any of
our money. The reason: Stephen had not seen a doctor, in person, to confirm that he'd actually been
sick.

The company is correct: Stephen did not see his doctor. We didn't see any reason for a visit because
he'd had this very same ailment before and we knew that nothing could be done about it. We had no
idea that by failing to see a doctor, we were forfeiting \$7,000 without even the opportunity for a
partial refund or credit. That is an unconscionable amount of money to sacrifice on a technicality.
Is there anything you can do? Margaret Moser

Austin, Tex.

A. Ms.~Moser wrote in July, soon after a Haggler column about a similar travel insurance problem. A
call to CSA, which is based in San Diego, was returned by a very pleasant woman named Jennifer
Wilbur, who works for a public relations firm. She underscored the salient paragraphs of CSA's
policy, which includes the company's definition of ``medical treatment'' as ``treatment, advice or
consultation, given in person, by a legally qualified physician.''

Ms.~Wilbur explained that after the Mosers canceled their trip, CSA asked Stephen's doctor, Theresa
Pugh, for her records. The company didn't find an appointment for Mr.~Moser. Case closed.

The Haggler noted that the point of CSA's rules is to separate the truly ill from those who want to
bail out of their trips for nonmedical reasons. Eminently reasonable. But nobody disputes that
Mr.~Moser was, and remains, seriously sick. (As he approaches his 52nd birthday, his health is
declining rapidly and he has moved into his mother's house because he can no longer live
independently.) To focus on the narrow question of whether Mr.~Moser saw a doctor in person ignores
every other salient fact about his condition, not to mention the many visits he'd made to his doctor
before his trip was to begin.

Ms.~Wilbur listened with all the equanimity required of a P.R. pro. Soon, she seemed like one of
those unfortunate bystanders from a ``Road Warrior'' movie, who become human shields strapped to the
front of trucks attacking Mel Gibson. The Haggler needed to speak to one of the truck drivers.

Eventually, there ensued a conversation and an e-mail exchange with Bob Chambers, CSA's director of
operations. He went back over CSA's policy and added that it would be unfair to other customers if
the company were to pay the Mosers' claim. Unfair to those who had seen a doctor, in similar
circumstances, and unfair to those whose claims were denied for the same reason as the Mosers'.

There is a cold, clear logic to this argument that is both totally irrefutable and deeply unjust.
You get the impression that CSA is more interested in minutiae than the basic question of whether
Mr.~Moser was really sick -- which is the point of the minutiae, after all.

In August, the Haggler reached out to the Texas Department of Insurance, which regulates CSA in
Texas. That didn't help. A few weeks after the agency wrote to the company, asking it to review the
Moser case, CSA reported that it had come to the same conclusion for the same reasons.

Time for a Hail Mary pass. Two weeks ago, the Haggler wrote to CSA's president, Les Maine, a kindly
looking gray-bearded gent, judging from a picture on the company's Web site.

``CSA and insurance companies cannot pay a claim simply because they feel sorry for an individual,''
Mr.~Maine wrote back. ``We need proper documentation. To pay claims arbitrarily would put the entire
system, and the integrity of our dealings with all of our other policy holders, in jeopardy.''

It seems to the Haggler that the insurance system would survive if CSA were to return the Mosers'
money. But this is apparently one of those situations where what seems fair in the broadest sense of
the word collides with what is considered fair in the unsentimental realm of contract language.

Two business school ethics professors contacted for this column found nothing to fault in CSA's
conduct, although a third, Thomas Donaldson of the Wharton School of Business of the University of
Pennsylvania, described that conduct as ``legal but blatantly unethical.''

Mr.~Moser, meantime, has already outlived his doctor's prognosis, made in early 2008, that he had a
year to live.

\section{Looking at Life as One Big Subscription}

\lettrine{E}{VERYWHERE}\mycalendar{Oct.'09}{12} you look these days, businesses are selling
subscriptions. Cable television, Internet and cellphone services are sold that way. So are business
software, office printing and car rentals like Zipcars.

For a business model that was used by publishers of periodicals back in the 17th century,
subscriptions seem as functional and popular as ever for a variety of goods and services.
Periodicals themselves are suffering -- Cond\'e Nast announced last week that it was closing Gourmet
and three other magazines -- but not because there's anything wrong with the practice of selling
subscriptions.

Marketers like them for good reason: Convince someone to take a subscription, and the revenue flows
in for months to come. ``It is amazing how inertia takes over,'' says Peter S.~Fader, a marketing
professor at the University of Pennsylvania. Anyone who has signed up for a gym membership that is
paid for but not used understands the genius behind subscriptions.

There is another reason that marketers use them: When a product has built-in obsolescence, like new
versions of software or a magazine with a short newsstand life, subscriptions extend the ownership
period. ``It removes the impedance to upgrading,'' says Erica Mina Okada, an associate professor of
marketing at the University of Hawaii.

Most of us mentally account for nearly all the money we spend, whether or not we realize it. We
assign a book value to a purchase, Professor Okada says, based on what we pay. Usually we won't
replace that product until we think we've gotten our money's worth. That explains why some people
continue to wear ill-fitting shoes rather than chucking them.

The rational way to think about your purchases, says V.~Seenu Srinivasan, a Stanford business school
professor, is to replace a product when the benefits of a new version outweigh the costs --
financial and psychological -- of upgrading. (He has developed a mathematical model for forecasting
when consumers will upgrade, based on this theory.)

A subscription moves consumers over the hurdle of mentally depreciating an existing asset. When you
go on vacation and don't get any movies from Netflix, it is easier to accept having wasted \$30 for
your subscription that month than it would be to have bought a \$30 DVD and never watched it.

Marketers love the subscription concept, but is it any good for consumers? It sometimes seems as if
cellphone deals were designed to alienate customers rather than to lure them into that sweet garden
of inertia.

How many times have you heard someone say, ``I can't wait to get out of my cellphone contract''? In
many cases, the customers think they are paying too much because they don't understand that the
subscription bundles the price of the phone with the price of the service. (That's one purpose of
early termination fees, to cover the cost of the hardware. Another is to keep you from canceling
your subscription.)

Some of the madness of the recent housing bubble can be blamed on an extension of the subscription
mentality. What exactly were homeowners doing when they bought a house for little or no money down
with the intention of holding it for two or three years before upgrading to a better home? They were
treating an investment in real estate as though it were just another consumable product, to be
disposed of with the same emotion one shows in recycling a monthly magazine.

People who thought of their home that way got away with it only as long as the home's value grew
enough to let them sell and trade up.

But thinking of other kinds of purchases as subscriptions may actually be useful, especially when
they are often-replaced products.

My colleague Saul Hansell has pointed out in the Bits blog of The New York Times that you
essentially do that when you buy PCs or digital cameras. Those who love technology might pay \$900
every two years to get a new and improved product, and more utilitarian users might spend \$300
every four years. Those who see technology as a bother might shell out \$150 every eight years to
upgrade.

Frugality is a virtue, as we remind ourselves during this recession. But it's not ridiculous to
spend on upgrades. Certainly some people value them more than others, just as some value the premium
HBO cable package or unlimited texting on their cellphone plan. But a consumer can become more
efficient by upgrading to a more fuel-efficient car or by buying a computer that crashes less often
or boots more quickly. A new smartphone gives a person access to information everywhere, a
liberating experience for anyone searching for a restaurant or a theater where no computer,
bookstore or newsstand is nearby.

PROFESSOR OKADA, who has studied how marketers persuade consumers to replace what they have, says,
``People don't upgrade as frequently as they should if they were acting rationally.''

If more people thought this way, it could fuel consumption. Of course, it would also be likely to
expand our disposal society. The two-year-old DVD player from Best Buy or the lightly used Leksvik
bookcase from Ikea will be piled out by the curb next to old newspapers and magazines.

The publishers of those magazines and newspapers recognize that the curbside stack is dwindling, and
most know that they must offer their products online. Every media company on the Web seems to be
trying to figure out how to return to paid subscriptions. It may be a lost cause, but new electronic
reading devices like the Kindle or the rumored Apple tablet computer could offer a means to wean
readers from the free business model.

\section{Is It a Day to Be Happy? Check the Index}

\lettrine{H}{EY}\mycalendar{Oct.'09}{12}, America, to quote the relentlessly upbeat temp in the
movie ``Office Space,'' it sounds like somebody has a case of the Mondays.

There is a 9.7 percent increase in happiness on Fridays compared with the worst day of the week,
Monday. That is among the discoveries made by Facebook researchers with access to two years of
anonymous ``status updates'' from 100 million users in the United States.

Another conclusion: holidays like the Fourth of July, Thanksgiving and Halloween make Americans
happy, while days when celebrities like Michael Jackson or Heath Ledger die make Americans sad. Such
insights -- obvious though some may be -- leap out from the new Gross National Happiness Index,
unveiled last Monday by Facebook.

The index is an attempt to say something profound based on the reports of daily life that Facebook
users impose on -- um, share with -- their friends.

The idea, one that is generally accepted in social psychology, is that word choice can reveal a
person's mood. This is true in ordinary writing, these experts say, and even more so in writing like
Facebook updates or the tweets of Twitter users, which ostensibly are attempts to describe what you
are doing right now and how you feel. (While tweets limn the psyche in 140-character installments,
Facebook updates are downright Augustinian, with some running as long as 420 characters.)

``When people in their status updates use more positive words -- or fewer negative words -- then
that day as a whole is counted as happier than usual,'' the creator of the index, Adam D.~I. Kramer,
explained in a Facebook blog post unveiling the project, where he also said that he was 72 percent
happier than the average Facebook user.

•

Mr.~Kramer hardly seemed giddy in a telephone interview on Thursday -- he seemed determined,
precise, ambitious, even harried, since he is working to finish his doctoral dissertation in social
psychology at the University of Oregon.

But he did seem optimistic. The Facebook happiness index, he said, could be the first step in
reorienting the nation's sense of self-worth.

``We have tracked the economic health of the nation for a long time,'' he said. ``The reason we
track those things is that the government is full of economists, not psychologists.''

``If we know money doesn't buy happiness,'' he asked, ``why are we optimizing for money?''

The idea that countries pursue the wrong policies because they collect the wrong data has been
gaining traction. The French president, Nicolas Sarkozy, created a commission last year that
included the Nobel Prize winning economists Joseph Stiglitz and Amartya Sen to come up with better
models for judging ``economic performance and social progress.''

Their report, which was issued last month, recommended considering many factors, including
``well-being'' and income inequality.

Though Mr.~Kramer sees the happiness index as superior to the gross domestic product, a quick glance
shows that for now the index -- with its daily leaps up and down and long-term trend lines --
resembles nothing so much as a stock ticker, with Mother's Day replacing a good earnings report in
explaining a spike.

The obvious triggers of peaks and troughs in the happiness index may raise snickers, but to
researchers they are reassuring and hold the promise of less obvious insights.

``It's early,'' said James W.~Pennebaker, a psychology professor at the University of Texas at
Austin who is one of the people responsible for the Linguistic Inquiry and Word Count dictionary,
which Facebook used to assess words for happiness and sadness. (Happy words include, naturally,
``happy,'' ``yay'' and ``awesome''; sad words include ``sad,'' ``doubt'' and ``tragic.'')

Such playful and public-minded research is relatively rare for Facebook when compared with the
much-larger Google, which devotes great resources to experimenting with its data.

But Facebook has one clear research advantage over Google, said Cameron Marlow, the head of
Facebook's 10-member research department: its material is much richer than mere search terms.
Facebook records relationships -- how they grow, how they falter, how they change, how they persist.

``The type of things that people reveal on Facebook are the kind of things that sociologists have
tried to collect through surveys for decades,'' he said. ``Say I want to understand how the divorce
rate is changing -- for a social scientist to produce a report on this, it would take a team of
researchers, a company to collect the data.'' But on Facebook, people routinely provide such
information in great detail as it happens in their lives.

Both companies are trying to answer questions central to their businesses. Facebook's questions are
just more interesting. Google wants to know how to provide the information you want as efficiently
as possible; Facebook wants to know how people get and keep friends.

Mr.~Marlow said his researchers were analyzing the ethnic and racial backgrounds of Facebook users
-- in the aggregate, he stressed, using census data on the distribution of last names -- to see
which groups are overrepresented and which are underrepresented. Such research can naturally evolve,
he said, into understanding how groups interact on Facebook, and how those interactions change over
time, another important indicator of the state of the nation.

``This is an example that was largely sparked by internal discussions,'' he said. ``We are asking
for obvious business reasons to see if there are biases in Facebook.''

•

To Mr.~Pennebaker, the Facebook material has the potential to unlock secrets about how people
interact. ``It is an incredible resource in terms of studying human connections, friendship, love,
hate, stigma,'' he said. ``This kind of work should be required for social psychologists.''

He agreed, though, that it likewise has the potential to turn ``creepy,'' with computers managing to
probe deeper into our minds than any person could.

``I could imagine it would allow us to look at a group of people, get a sense of what their concerns
are, how insecure they feel,'' he said. ``It could be an advertiser's dream. Yes, it is creepy from
a government perspective, but it is even creepier from an advertising perspective.''

At a minimum, we now all know that there is sound basis not to name a restaurant Thank God It's
Monday.

\section{Closing the Deal at the Virtual Checkout Counter}

\lettrine{S}{hoppers}\mycalendar{Oct.'09}{12} rarely drive to the mall, load up their carts and then
abandon them in the middle of the store. On the Web, though, it happens all the time.

In online stores, it is much easier for shoppers to fill their virtual shopping carts -- and much
easier for them to get distracted by an e-mail message or comparison shopping on other sites. Then
there are the design flaws and technical glitches\cite{glitch} that can get in the way of closing a
sale.

These problems have been around since online shopping was invented, but they have taken on more
urgency in the last year as consumer spending has shriveled. So e-commerce companies are trying a
variety of techniques to push shoppers through the virtual checkout line.

There are still plenty of people browsing online, but not so many buyers. In the second quarter, the
number of visitors to e-commerce sites who eventually bought something shrank for most sites from
the year before, by as much as 30 percent for Zappos.com and 26 percent for Gap, according to
comScore.

``It's pretty clear that people are looking at more alternatives, evaluating more options, getting
better prices -- but not buying,'' said Gian M.~Fulgoni, executive chairman of comScore.

Shoppers spent \$130 billion online in the last year, according to comScore. But e-commerce sites
missed out on billions more because customers abandoned their carts once they ran into problems
while checking out, according to Tealeaf, a company that makes software to help e-commerce sites
monitor customers' behavior.

``The small transactions add up,'' said Rebecca Ward, chief executive of Tealeaf, whose customers
include Wal-Mart and Best Buy. ``This is revenue that people really wanted to commit to the company
and were unable to do it, and it often ends up being in the millions of dollars.''

Many shoppers fill their carts just to keep track of things they like or to check shipping rates and
taxes, with little intention to buy. While there is no industrywide data, some e-commerce companies
estimate that only about 3 percent of shoppers who visit an e-commerce site buy something, and when
they do load their shopping carts, as many as two-thirds abandon them.

One of the biggest reasons people procrastinate more when shopping online is the fear of regret,
said Dan Ariely, a professor of behavioral economics at Duke, a visiting professor at M.I.T. and
author of the book ``Predictably Irrational.''

It is much simpler online than offline to discover that an item you bought yesterday is on sale
somewhere else today. In fact, he said, people often spend more time researching a product after
buying it online than before, to prove that they should not regret the purchase.

Online retailers do a few things to fight this inclination. Zappos.com and Overstock.com inject
urgency by alerting customers when an item they have put it in their shopping cart is almost sold
out.

``The Internet gives us this ability to kind of have wish lists -- you can look at 15 products,
throw them in your cart and then sleep on it,'' said Stormy D.~Simon, senior vice president for
branding and customer care at Overstock.com. ``We incentivize them.''

Other sites have developed a new, extreme version of limited-time sales. Gilt offers items for 36
hours or until they run out, and Neiman Marcus runs two-hour, online-only sales.

``All these strategies get people to fear that they will regret not acting'' instead of acting,
Mr.~Ariely said.

Customer reviews can also nudge would-be shoppers who are on the fence about making a purchase. A
company called Bazaarvoice helps e-commerce customers, like Macy's and Dell, publish user reviews
and ratings on their sites. For a customer unsure about new shoes or a new sofa, social validation
can make the difference.

Some sites try to make the check-out process easier by offering alternative ways to pay. Bill Me
Later, which eBay bought last year for \$820 million, lets people click one button to check out.
Later, they get a bill in the mail.

TrialPay, a start-up company, lets people get a product free if they buy another product. A shopper
can buy WinZip software online for \$29.95, for example, or get it free by signing up for Netflix.
Then Netflix pays TrialPay for bringing in the new customer.

Though the Web makes it easier for shoppers to abandon their carts, it also makes it easier for
shops to track would-be customers and encourage them to buy.

``In the real world, the jeweler or optician has no way of knowing who I was or how to get me back
in the store, but online you can do all that, which is why it's such an amazing retail
opportunity,'' said Saul Klein, a partner at Index Ventures, which has invested in e-commerce
companies, including TrialPay.

Index has also backed a start-up called Criteo, which lets e-commerce sites ``follow'' visitors who
leave without making a purchase and show them banner ads when they visit another site. Say a shopper
has been perusing digital cameras on a consumer electronics site, then goes to lunch. Later, he
checks the headlines on a news site, where he is shown an ad for the digital camera site, luring him
back.

Some e-commerce sites encourage shoppers to log in before they fill their carts. Then, if they
leave, the site can send them an e-mail message reminding them that their cart is still there and
perhaps offering a carrot, like free shipping. Tealeaf's software can identify each registered
shopper who got to a certain point in the buying process before giving up.

It also alerts shopping sites about technical problems that might otherwise have been invisible. A
month ago, the clothing retailer Bluefly realized that some international shoppers were unable to
check out. Using Tealeaf's software, Bluefly discovered that the glitch had been there for a year.
Instead of reporting the problem, customers had simply been leaving the site without making the
purchase.

After Bluefly fixed the problem, revenue from international shoppers increased 10 percent in a
month, and Matt Raines, Bluefly's vice president for technology, estimated that the fix would result
in \$1.1 million in additional revenue this year.

Bluefly also runs daily promotions and timed sales and shows Bluefly ads to previous visitors when
they are on other sites. It is starting to offer customers the option to save the items in their
cart to buy later, and is considering running customer reviews.

``When customers are trying to purchase something, we need to do everything in our power to make
sure they can do it,'' Mr.~Raines said.

\section{China to End Required Testing for Hepatitis B}

\lettrine{C}{hinese}\mycalendar{Oct.'09}{13} health officials will abolish mandatory testing for
hepatitis B during physical exams given to prospective college students, factory workers and
government employees, according to Xinhua, the state news agency.

The new rules, announced Saturday by the Health Ministry, mean that people found to be carrying the
hepatitis B virus will not be automatically barred from jobs and classrooms, a form of
discrimination widely decried by health care advocates and the estimated 120 million people in China
thought to be infected. It is not uncommon in China for hepatitis B carriers to be barred from
medical school, teaching positions or jobs in the food production industry.

In announcing their decision, officials suggested that they were yielding to public pressure,
especially from activists who in recent years have begun organizing on the Internet.

``On account of the questions brought up by media and society concerning hepatitis B testing, the
Health Ministry has come to a consensus,'' Deng Haihua, a spokesman said in comments posted on the
ministry's Web site. ``The current social misunderstanding about hepatitis B patients is mainly a
result of a lack of understanding. The experts believe that canceling blood testing for the purpose
of employment and students entering schools will not affect the health of others, nor will it cause
the disease to spread.''

Lu Jun, who runs an organization in Beijing that has been pushing the government to ease its testing
policies, described the new rules as a victory for hepatitis B carriers as well as those trying to
change society through public advocacy and legal action. ``International standards have now been
realized in China,'' he said in an interview.

In July, the offices of Mr.~Lu's group, the Yirenping Center, were raided by the police, who
confiscated literature they said had been illegally published.

Many of those who carry hepatitis B in China were infected by contaminated syringes during childhood
inoculation campaigns in the 1970s and '80s. Others contracted it from their mothers during
childbirth.

\section{Obama Becomes Japan's English Teacher}

\lettrine{W}{hen}\mycalendar{Oct.'09}{13} Utako Sakai was changing the background music in her
beauty parlor recently, she did not opt for the classical piano pieces she usually chose.

Instead, she picked her favorite CD: ``President Obama's Inaugural Address,'' released by Asahi
Press, a Japanese publisher of language books. She says the speech lifts her spirits and helps her
to learn English all at once.

``All our customers love it,'' said Ms.~Sakai, who is based in Ayase City, in Kanagawa Prefecture,
outside Tokyo.

The speech CD and its accompanying book have been a resounding success, selling 200,000 copies since
its release in January. A compilation of President Barack Obama's speeches has done even better,
selling half a million copies since November, solidifying his role as Japan's English teacher.

Publishers have since flooded the market with over a dozen language-learning titles, including
``Speech Training: Learning to Deliver English Speech, Obama Style''; ``Learn English Grammar From
Obama''; and ``Yes, I Can With Obama: 40 Magical English Phrases From Presidential E-mails.''

Asahi Press followed up its inauguration book and CD with a recording of Mr.~Obama's ``World Without
Nuclear Weapons'' speech, also in book and CD form, given in Prague in April.

The publishers are trying to tap into a foreign-language teaching industry that the Yano Search
Institute said was valued at ¥767 billion, or \$8.7 billion, in 2008. The figure includes the cost
of books, CDs, dictionaries, e-learning programs, standardized English tests, and the cost of
private language lessons. The institute, in Tokyo, says the majority of the spending is aimed at
learning English.

Most Japanese people, including those studying English, would have difficulty comprehending a speech
given by a native English speaker. But ``Mr.~Obama's English is easy to understand because he
pronounces words clearly and speaks at a relatively slow clip,'' said Professor Tadaharu Nikaido, a
communication specialist here. ``Movies tend to be the most difficult for Japanese, especially when
actors mumble their words.''

Mr.~Obama sets his range of vocabulary wide enough to accommodate the highly educated and the less
educated, Professor Nikaido added, and at the lower end, it sometimes comes within the range of
non-native speakers' comprehension.

But there are probably a large number of buyers who do not really possess the basic English skills
to understand his speech, said Yuzo Yamamoto, an editor at Asahi Press. Since the sales took off, he
has received postcards from readers saying they had been touched by Mr.~Obama's speeches, but
``those same people have said they were moved even though they didn't understand English well,'' he
said. ``Some even said the only phrase they caught was, 'Yes, we can.' They said they were in tears
nonetheless.''

Mr.~Yamamoto said there was a sincerity about Mr.~Obama's speaking style that listeners could
perceive phonetically, combined with a delivery that was almost musical.

``That seems to result in sensation, the kind of which you get from listening to good music,'' he
said.

Other observers say that Japanese buyers probably feel as though they understand his speeches just
from the nonverbal cues.

``The audience in the background helped too,'' Professor Nikaido said. ``The audience's echo in the
background works the same way as, say, laugh sound effects inserted in TV comedies.''

That may explain why Ms.~Sakai, the beautician, and her customers were so enthralled by Mr.~Obama.
Ms.~Sakai describes her English skills as less than perfect and says she relies on others for
translation when she travels overseas. But when she hears Mr.~Obama talk, she feels perfectly at
ease.

``I feel as though I am not Japanese,'' she says, because she is able to understand the English
speech so well.

The fervor over Obama-speak, some say, reflects other aspects of Japan's changing society. Among the
public, there is now a longing for leaders with a communication style that is more effective,
dynamic and inspiring, something not prevalent in the political culture here, they say.

Since late last year, economic recession has only deepened with no end in sight, and hope has been
difficult to find. ``There is a sense, why isn't there an Obama in Japan who can raise people's
spirits,'' Mr.~Yamamoto said. This year, as the Obama books were beginning to boom, the prime
minister then, Taro Aso, was making a series of political gaffes, Mr.~Yamamoto noted. They included
misreading kanji, the Japanese characters, in public.

The Obama speech phenomenon peaked in the spring, observers say, and has given way to more books and
magazines that focus on how Mr.~Obama's communication skills might be adopted in Japan for business
and political purposes.

Professor Nikaido himself has recently published a book called ``Creating an Audience Frenzy:
Learning From Obama's Strategic Oratory.'' At least half a dozen books and magazine covers of this
kind have appeared in the past several months.

John R.~Harris, a Canadian speechwriter who is based in Chiba, a Tokyo suburb, and who has worked
with Japanese politicians, says he knows a good deal about the dire straits of political
communication in Japan. The art is virtually nonexistent and is only beginning to be discovered, he
said.

``Japan has not been serious about communication,'' he said. ``In a Japanese company or political
party or anyplace where Japanese come together as a group, the process is consensus-forming, and the
outcome has to be consensus, and the consensus is internal. In that, the audience often gets
forgotten.''

\section{Editorial Dispute Threatens a Chinese Magazine}

\lettrine{C}{aijing}\mycalendar{Oct.'09}{13}, a respected Chinese business magazine, has been thrown
into turmoil\cite{turmoil} after a dispute over its editorial policy and business governance led to
the resignation of 11 high-ranking executives and nearly 70 other workers from the business staff,
according to people close to the magazine.

The magazine's general manager, its top advertising executive and the head of its conference unit
are among those who have submitted resignation letters. And Hu Shuli, considered the most powerful
business editor in China, may be forced to resign from the magazine, which is based in Beijing.

The dispute threatens to dismantle one of the country's leading media properties, a thriving
magazine published twice a month that specializes in investigating government corruption and
corporate fraud. Caijing also has partnerships with The Wall Street Journal and Reuters.

For months, Ms.~Hu, who as the magazine's managing editor is its top editorial employee, and the
magazine's business managers have been locked in a struggle with Caijing's owners over the
magazine's editorial policy, according to former employees and current staff members who asked not
to be named because they feared for their jobs.

The owners of the magazine have recently come under pressure from some within the government to tone
down or drastically alter Caijing's aggressive journalism, people at the magazine say. Caijing's
managers have told staff members that they have been fighting to maintain the magazine's editorial
integrity.

Caijing's managers have been seeking to create a more independent publication by changing the
magazine's shareholding structure, seeking outside investors and pressing the owners to allow some
employees to own a stake in the magazine. They also want a larger share of the magazine's profits to
be invested in new operations, including an English-language Web site.

But the magazine's owner, the Stock Exchange Executive Council, which is state-owned and partly
private, has apparently refused to surrender control or restructure its operations.

About 70 percent of employees on the magazine's business staff have resigned, though some say they
will continue to work until replacements are found.

The resignations were first reported in the Monday edition of The South China Morning Post in Hong
Kong. A spokeswoman for Caijing, which has a circulation of about 225,000 copies, declined to
comment. A person who answered the phone at the Beijing headquarters of the magazine's owner said no
one was available to comment.

China's media is tightly censored, particularly on political matters. But business publications have
broader press freedoms, and Caijing is widely respected, even by high-ranking government officials.
The magazine is led by Ms.~Hu, 56, who because of her independent style of journalism is often
labeled ``the most dangerous woman in China'' by other media.

A former propaganda writer for Workers' Daily, the Communist Party publication, she helped found
Caijing in 1998, and has been praised for trying to build a stronger publication.

She was recently the subject of a lengthy profile in The New Yorker magazine, which called her ``an
incurable\cite{incurable} muckraker.''

Several people familiar with the dispute said Ms.~Hu continued to negotiate with the owners, but
that if talks broke down she could resign and form a new media property with Wu Chuanhui, who
recently resigned as Caijing's general manager.

``A lot of people just want to leave and follow Hu Shuli,'' one former staff writer said over the
weekend. ``They are thinking of creating a similar magazine and other divisions and running the
business themselves. Caijing makes a lot of money from advertising, and they think they can manage
it better.''

\section{Disney's Retail Plan Is a Theme Park in Its Stores }

\lettrine{L}{os}\mycalendar{Oct.'09}{13} angeles -- The Walt Disney Company, with the help of Steven
P.~Jobs and his retailing team at Apple, intends to drastically overhaul its approach to the
shopping mall.

At a time when many retailers are still cutting back or approaching strategic shifts with extreme
caution, Disney is going the other way, getting more aggressive and putting into motion an expensive
and ambitious floor-to-ceiling reboot of its 340 stores in the United States and Europe -- as well
as opening new ones, including a potential flagship in Times Square.

Disney Stores, which the media giant is considering rebranding Imagination Park, will become more
akin to cozy entertainment hubs. The chain's traditional approach of displaying row after row of
toys and apparel geared to Disney franchises will be given a high-tech makeover and incorporated
into a new array of recreational activities. The goal is to make children clamor to visit the stores
and stay longer, perhaps bolstering sales as a result. Over the next five years, analysts estimate
that Disney will spend about \$1 million a store to redecorate, reorganize and install interactive
technology.

``The world does not need another place to sell Disney merchandise -- this only works if it's an
experience,'' said Jim Fielding, president of Disney Stores Worldwide. The company plans to unveil
the new look in May in Southern California, Long Island and Madrid, and is close to signing a lease
for that Times Square flagship.

Theaters will allow children to watch film clips of their own selection, participate in karaoke
contests or chat live with Disney Channel stars via satellite. Computer chips embedded in packaging
will activate hidden features. Walk by a ``magic mirror'' while holding a Princess tiara, for
instance, and Cinderella might appear and say something to you.

It's your birthday? With the push of a button, eight 13-foot-tall Lucite trees will crackle with
video-projected fireworks and sound. There will be a scent component; if a clip from Disney's coming
``A Christmas Carol'' is playing in the theater, the whole store might suddenly be made to smell
like a Christmas tree.

The makeover happened only after much internal debate at the company. Indeed, some Disney board
members fretted that the concept was so lavish that parents would try to use the stores as day care
centers. Others worried that people would come for the entertainment but not buy anything.

``It's time to take risks,'' Mr.~Fielding said he told them. ``When consumers are ready to spend
again, we will be ready.''

The involvement of Mr.~Jobs, the Apple chief executive who joined the Disney board with the 2006
acquisition of Pixar, is particularly notable. For the first time, Mr.~Jobs's fingerprints can be
seen on Disney strategy, in the same way that he influenced the look and feel of Apple's own
immensely popular retail chain. While Mr.~Jobs did not personally toil on the Imagination Park
concept, he pushed Disney to move far past a refurbishment.

``Dream bigger -- that was Steve's message,'' said Andy Mooney, chairman of Disney Consumer
Products.

Mr.~Jobs provided access to proprietary information about the development and operation of Apple's
highly successful stores, and Disney executives visited Apple's research operation in Cupertino,
Calif. Mr.~Jobs, who declined to comment, also insisted that Disney build a prototype store to work
out kinks, a costly endeavor that most retailers skip.

The company followed his advice, working for the last year on a full-scale, fully stocked store
inside an unmarked warehouse in Glendale, Calif. The prototype was crucial to shaping an overall
philosophy, Mr.~Fielding said, noting that he discovered the shops needed more ``Pixar-esque winks
and nods.'' To that end, one sales area is now labeled ``WWTD: What Would Tinker Bell Do?''

Disney will adopt Apple touches like mobile checkout (employees will carry miniature receipt
printers in their aprons) and the emphasis on community (Disney's theater idea is an extension of
Apple's lecture spaces). The focus on interactivity -- parents will be able to book a Disney Cruise
on touch-screen kiosks while their children play -- reflects an Apple hallmark. Employees can use
iPhones to control those high-tech trees.

Disney is a merchandising titan whose licensed consumer products generated \$30 billion in global
sales last year, up from \$12 billion when Mr.~Mooney joined the company a decade ago. But Apple is
king of the mall. Its fleet of stores generated sales of about \$4,700 a square foot in 2008, by far
the highest for any retail chain, said Charlie Wolf, an analyst at Needham \& Company. In
comparison, Best Buy's sales are about \$1,000 a square foot.

The Disney board approved the Imagination Park concept on Oct.~1 after touring the prototype and
receiving hand-made books from Mr.~Fielding that pitched the concept as ``the best 30 minutes of a
child's day.'' Now Disney is bringing in landlords, trying to pit them against one another to secure
top-tier locations and favorable leases.

``We will essentially be the only toy retailer left at the mall because everybody else has
evaporated\cite{evaporate},'' Mr.~Mooney said. Mr.~Fielding added, ``Every mall in America is
desperate for newness and freshness.''

The Disney Store chain, introduced in 1987, was initially so successful that the company
overexpanded to more than 600 locations. Buffalo alone had five. But consumers overdosed on the
animated-character merchandise and by 2002 the chain was losing about \$100 million a year.

Judging the upkeep too burdensome and focusing on the safer licensing business, Disney sold the
chain to Children's Place Retail Stores in 2004.

But Children's Place failed to meet contractual renovation obligations and consumers complained of
poor service, Mr.~Fielding said. The vaunted Disney brand was in peril. ``Let's face facts,'' he
said, ``some of those stores looked like a dog's breakfast.''

In March 2008, Disney bought back part of the chain on undisclosed terms and the remaining stores,
about 100 of them, were closed. A spokeswoman for Children's Place declined to comment.

Dressing up a toy store with entertainment is hardly new. F.A.O. Schwarz has its famous floor piano,
Toys ``R'' Us operates a Ferris wheel inside its Times Square outpost and American Girl jazzes up
its retail outlets with theaters and doll hair salons. But the emphasis on programming -- via the
theaters, Disney is essentially creating a mall-based television channel -- and the degree to which
the media giant is pouring on the razzle-dazzle in every store is unusual.

The nation's largest mall operator, Simon Property Group, however, said Disney might be overly
optimistic about how far landlords will bend. ``No one comes in here and dictates terms,'' said
Bruce Tobin, a Simon executive vice president.

That said, however, Mr.~Tobin had pinwheels in his eyes after touring the prototype. ``It's truly
spectacular -- beyond our imagination,'' he said. ``These are going to be true destinations.''

\section{Credit Tightens for Small Businesses }

\lettrine{M}{any}\mycalendar{Oct.'09}{13} small and midsize American businesses are still struggling
to secure bank loans, impeding their expansion plans and constraining overall economic growth, even
as the country tentatively rises from its recessionary depths.

Most banks expect their lending standards to remain tighter than the levels of the last decade until
at least the middle of 2010, according to a survey of senior loan officers conducted by the Federal
Reserve Board. The enduring credit squeeze appears to reflect an aversion\cite{aversion} to risk
among lenders confronting great uncertainty about the economy rather than any lingering effects of
the panic that gripped financial markets last fall, after the collapse of the investment banking
giant Lehman Brothers.

Bankers worry about the extent of losses on credit card businesses as high unemployment sends
cardholders into trouble. They are also reckoning\cite{reckon} with anticipated failures in
commercial real estate. Until the scope of these losses is known, many lenders are inclined to hang
on to their dollars rather than risk them on loans to businesses in a weak economy, say economists
and financial industry executives.

``The banks are just deathly afraid,'' said Sam Thacker, a partner at Business Finance Solutions in
Austin, Tex., which helps small businesses line up financing. ``I don't see commercial banks coming
back to the market anytime soon.'' In the long view, tighter loan standards seem healthy after a
terrible crisis attributed in part to years of recklessly lenient lending.

But some economists worry that bankers have overshot the boundaries of a healthy reaction, as even
strong companies are finding it difficult to borrow.

``The banks are still very risk averse,'' said Robert J.~Barbera, chief economist of the research
and trading firm I.T.G. ``Regional banks are in a particularly tough spot, because they're choking
on commercial and residential real estate.''

Bankers acknowledge that loans are harder to secure than in years past, but they say this attests to
the weakness of many borrowers rather than a reluctance to lend.

``Banks want to lend money,'' said Raymond P.~Davis, chief executive of Umpqua Bank, a regional
lender based in Portland, Ore. ``The problem is the effect that the recession is still having on us.
Some of these businesses are still trying to come out of it. For them to go to a bank, if they are
showing weak performance, it is harder to borrow.''

As the financial crisis has largely eased in recent months, big companies have found credit
increasingly abundant, with bond issues sharply higher.

But for ordinary consumers, the picture is quite different. What was once of a flood of come-ons for
home equity loans and credit cards has been replaced by notices of lowered credit limits.

For many smaller companies, too, borrowing remains tough.

Some 14 percent of small businesses found loans harder to secure in August than in July, according
to the most recent survey by the National Federation of Independent Business. Among companies
borrowing regularly, less than one-third reported that all their credit needs were being met.

``It's quite significant, because small businesses generate significant job growth,'' said Andrew
Tilton, a senior economist at Goldman Sachs. ``And small businesses rely more on bank financing,
whereas large businesses have the alternative of raising money in the capital markets.''

Businesses with fewer than 500 employees hold more than half of the nation's private sector jobs,
according to the Small Business Administration.

In Port Arthur, Tex., Five Star Feeds, which sells gardening supplies, pet food and livestock feed,
is stuck in a holding pattern on a planned expansion as the owner, Tina E.~Bean, tries to persuade
her local banker to provide \$150,000 in credit.

Ms.~Bean began the business in 2001 and has been profitable ever since, she said. She has already
borrowed about \$800,000 for her expansion, more than doubling the size of her storefront as she
adds new product lines, including Western-style clothing like jeans and cowboy boots.

But without the last \$150,000, she is unable to order the new products. Her new shelves remain
bare. She missed the deadline to tap her final \$150,000 in credit, she says, and her banker refused
to offer a new loan, even as she insisted that she could not move into the new space.

``His answer is, `We'll cross that bridge when we get there,' '' Ms.~Bean said. ``Well, I'm there. I
can't open the store if I don't have anything to put in it.''

Until she obtains financing, she has no reason to hire the two or three people she will need to run
the larger space. She has no reason to place orders that would increase work for truck drivers who
ferry blue jeans and pet grooming supplies to retail stores, and no need to bring on accountants,
insurance agents, lawyers, forklift drivers and auto mechanics whose business opportunities would be
incrementally expanded by more commerce.

Among small privately held companies, the amount of debt they carry as a portion of their equity has
slipped by about 5 percent since 2007, according to Sageworks, a financial analysis firm in Raleigh,
N.C. The drop reflects not only how companies have cut their inventories and paid down debt, but
also the tightened credit terms they face when they try to borrow, said a Sageworks spokeswoman,
Melinda Crump.

For companies that work on a contract basis bidding for jobs in advance, tight credit sometimes
precludes their ability to seek new business: They cannot raise capital to hire workers and add
equipment.

In Alexandria, Minn., Bob Novak, president of ACB Construction, has in recent months declined to bid
on large-scale government jobs in Utah because of what he describes as a frustrating lack of credit
for people in his line of work. His work force has shrunk to 1, from 14, in the last two years, yet
every effort to generate fresh business seems to get snuffed out by rejections from banks.

``You might as well not walk through the door of a bank and have construction behind your name,''
Mr.~Novak said. ``You cannot get a loan if you don't look profitable, and you can't look profitable
if you can't bid the work. I've had to walk away from jobs because I cannot get the funding.''

Med-National provides medical and dental services on a contract basis, mostly to the Defense
Department. The San Antonio-based company has been in business for more than 20 years, has never
lost money, and logged \$5.9 million in sales last year, said the executive vice president, Robert
S.~Welborn Jr.

Yet as the company recently sought to secure a pair of contracts with the United States Army, its
plans were delayed for nearly a year as it tried and failed to persuade a bank to lend it \$300,000.
The financing was required to hire the roughly 20 people needed to deliver the services -- doctors
and dentists and support staff.

Mr.~Welborn initially called local banks with whom he had done business in years past, among them
JPMorgan Chase and Wachovia. Their answer surprised him: no.

``It used to be you could practically just jangle your keys and they would give you some money,''
Mr.~Welborn said. ``It isn't anything like that anymore.''

In August, Mr.~Welborn finally lined up financing from Wells Fargo and began hiring. Yet the delay
and the uncertainty have left him reluctant to pursue other contracts.

``If we could borrow the money more easily,'' he said, ``we would be hiring even more people.''

\section{Google and Apple Eliminate Another Tie}

\lettrine{I}{n}\mycalendar{Oct.'09}{13} a move that severs another tie between Google and Apple,
Arthur D.~Levinson, the former chief executive of Genentech, has resigned from Google's board.

In a statement on Monday, Google gave no reason for Mr.~Levinson's departure, but earlier this year,
the Federal Trade Commission began an antitrust investigation into the close ties between the boards
of Google and Apple.

At the time, Mr.~Levinson and Eric E.~Schmidt, Google's chief executive, were directors of both
Google and Apple. The Clayton Antitrust Act of 1914 prohibits such board overlaps when they reduce
competition between two rival companies.

Mr.~Schmidt resigned from the Apple board in August.

In a statement, Jon Leibowitz, the chairman of the commission, suggested that the commission had
been ready to go to court to force the two companies to sever their board ties.

``Google, Apple and Mr.~Levinson should be commended for recognizing that overlapping board members
between competing companies raise serious antitrust issues, and for their willingness to resolve our
concerns without the need for litigation,'' Mr.~Leibowitz said. ``Beyond this matter, we will
continue to monitor companies that share board members and take enforcement actions where
appropriate.''

Two other Apple directors, former Vice President Al Gore and Bill Campbell, chairman of Intuit,
remain advisers to Google. But the commission said it was no longer pursuing its inquiry into the
two companies' connections.

When Apple appointed Mr.~Schmidt to its board in 2006, the move was seen as cementing an alliance
between two Silicon Valley giants with a common enemy: Microsoft. But as Google's business continued
to expand into new areas, it began bumping up against the interests of Apple with growing frequency,
fraying the ties between the two allies.

Apple entered the cellphone business when it introduced the iPhone in January 2007. Google followed
suit 10 months later, when it unveiled its Android software for mobile phones.

The growing tension between the two companies came into public view recently when Google accused
Apple of rejecting a proposed Google Voice application for the iPhone. Apple has said it is still
reviewing the application. The dispute has drawn scrutiny from the Federal Communications
Commission.

Google and Apple also compete in other areas, including Web browsers, photo-editing software and the
distribution of music and video on Apple's iTunes Store and Google's YouTube. And Google is
preparing an operating system for lightweight computers that could compete with the Macintosh
operating system.

When Mr.~Schmidt quit Apple's board, Steven P.~Jobs, that company's chief executive, noted that if
Mr.~Schmidt had stayed, his effectiveness as a director would have been diminished because he would
``have to recuse himself from even larger portions of our meetings due to potential conflicts of
interest.''

Earlier this month, however, Mr.~Schmidt insistedthat Mr.~Levinson should remain on both boards
despite the inquiry, saying that the revenue that Apple and Google received from competing products
was below a legal threshold.

Under the Clayton Act, so-called interlocking directorates are not considered a problem if the
revenue from products in which the companies compete is less than 2 percent of either company's
sales.

\section{In 1918 Pandemic\cite{pandemic}, Another Possible Killer: Aspirin}

\lettrine{T}{he}\mycalendar{Oct.'09}{13} 1918 flu epidemic was probably the deadliest plague in
human history, killing more than 50 million people worldwide. Now it appears that a small number of
the deaths may have been caused not by the virus, but by a drug used to treat it: aspirin.

Dr.~Karen M.~Starko, author of one of the earliest papers connecting aspirin use with Reye's
syndrome, has published an article suggesting that overdoses of the relatively new ``wonder drug''
could have been deadly.

What raised Dr.~Starko's suspicions is that high doses of aspirin, amounts considered unsafe today,
were commonly used to treat the illness, and the symptoms of aspirin overdose may have been
difficult to distinguish from those of the flu, especially among those who died soon after they
became ill.

Some doubts were raised even at the time. At least one contemporary pathologist working for the
Public Health Service thought that the amount of lung damage seen during autopsies in early deaths
was too little to attribute to viral pneumonia, and that the large amounts of bloody, watery liquid
in the lungs must have had some other cause.

Dr.~Starko acknowledged that she did not have autopsy reports or other documents that could prove
that aspirin was the problem. ``There was a lot of chaos in these places,'' she said, ``and I'm not
sure if there are good records anywhere.''

But of the many factors that might have influenced the outcome in any particular case, Dr.~Starko
wrote, aspirin overdose stands out for several reasons, including a confluence of historical events.

In February 1917, Bayer lost its American patent on aspirin, opening a lucrative drug market to many
manufacturers. Bayer fought back with copious advertising, celebrating the brand's purity just as
the epidemic was reaching its peak.

Aspirin packages were produced containing no warnings about toxicity and few instructions about use.
In the fall of 1918, facing a widespread deadly disease with no known cure, the surgeon general and
the United States Navy recommended aspirin as a symptomatic treatment, and the military bought large
quantities of the drug.

The Journal of the American Medical Association suggested a dose of 1,000 milligrams every three
hours, the equivalent of almost 25 standard 325-milligram aspirin tablets in 24 hours. This is about
twice the daily dosage generally considered safe today.

Dr.~Starko's paper, published in the Nov. 1 issue of Clinical Infectious Diseases, has stirred some
interest, if not enthusiastic endorsement, among other experts.

``I think the paper is creative and asking good questions,'' said John M.~Barry, author of a book on
the 1918 flu titled ``The Great Influenza.'' ``But we don't know how many people actually took the
doses of aspirin discussed in the article.''

The pharmacology of aspirin is complex and was not fully understood until the 1960s, but dosage is
crucial. Doubling the dose given at six-hour intervals can cause a 400 percent increase in the
amount of the medicine that remains in the body. Even quite low daily doses -- six to nine standard
aspirin pills a day for several days -- can lead to dangerously high blood levels of the drug in
some people.

Peter A.~Chyka, a professor of pharmacy at the University of Tennessee, said he found Dr.~Starko's
theory ``intriguing.'' Little was known about safe dosages at the time, he said, and doctors often
simply raised the amount until they saw signs of toxicity.

``In the context of what we know today about aspirin and aspirinlike products, Starko has made an
interesting effort to put this together,'' Dr.~Chyka said. ``There are things other than flu that
can complicate a disease like this.''

Although he doubted that more than a small number of deaths could be attributed to aspirin overdose,
Dr.~David M.~Morens, an epidemiologist with the National Institutes of Health, said the paper was
valuable in that ``it makes an attempt to look at environmental or host factors that may be
involved.'' He said, ``We haven't been able to explain all the deaths in young adults with the virus
itself.''

Dr.~Starko was hesitant to estimate how many deaths aspirin overdose could have caused, but
suggested that military archives might be one place to look. ``I'm hoping others will follow up,''
she said, ``by examining available treatment records.''

\section{The Big Surprise of 'Big Bang': The Bigger Audience}

\lettrine{T}{he}\mycalendar{Oct.'09}{14} most vibrant buzz this summer around the Warner Brothers
lot here and CBS Entertainment headquarters in nearby Studio City was not being generated by the
slate\cite{slate} of new shows on the CBS fall schedule. Rather, it focused on the sudden emergence
-- during summer repeats, no less -- of a series that had been on the air for two seasons.

``The Big Bang Theory,'' the CBS comedy about two brilliant physicists and their attempts to relate
to the world around them -- and to the cute blond woman next door -- began drawing surprisingly
strong ratings this summer after it moved to a later time slot on Monday, at 9:30 p.m., immediately
following that network's highest-rated comedy, ``Two and a Half Men.''

In some weeks of the summer ``Big Bang'' repeats drew bigger audiences among certain important
demographic groups than when the same episodes were first broadcast. So far this fall ``Big Bang''
has further expanded its audience, becoming the highest-rated live-action comedy among the
sought-after young-adult demographic group.

If current trends prevail, its total viewership could soon surpass that of ``Two and a Half Men,''
long the most-watched comedy on television. Last Monday's ``Big Bang'' drew 12.96 million viewers,
according to Nielsen, only 5 percent fewer than the 13.63 million for ``Men.''

Already ``Big Bang'' has beaten ``Men'' among viewers age 18 t0 49, the demographic category most
valued by advertisers.

The comedies have more in common than their popularity. They were co-created by Chuck Lorre, they
tape on adjacent stages on the Warner Brothers lot, and they share several writers and much of their
technical crews. And with the upstart closing in on the longtime ratings champion, Mr.~Lorre said,
he sometimes isn't sure how to react when the ratings come in.

``There's a lot of ambivalence\cite{ambivalence},'' he said on Tuesday night, during a break in the
taping of a ``Big Bang'' episode. ``It's 'Yeah!' then 'Awww.' But it's all good. I can't claim to
understand how this works; I'm just thrilled that it's working.''

The cast and crew of ``The Big Bang Theory'' are enjoying their success all the more after surviving
two near-death experiences. The show's first pilot was rejected by CBS, but the network asked
Mr.~Lorre and Bill Prady, his co-creator, to retool their script and try again. The first version
featured the same two male lead characters -- Jim Parsons as Sheldon Cooper, a theoretical
physicist, and Johnny Galecki as Leonard Hofstadter, an experimental physicist -- but also included
a female lead character who was ``very damaged and very tough,'' Mr.~Prady said.

``We had a really hard time casting the role, and in retrospect it was obvious that the problem was
not the actresses but the conception of the character,'' he said. Focus groups that watched the
original pilot were left with protective feelings for the two na\"ive, socially awkward scientists,
and they did not like the prospect of a bitter, manipulative woman taking advantage of them.

``What we all liked was the relationship between these two guys, one who wants his world to be
bigger and the other who wants his world to be smaller,'' Mr.~Prady said. ``I think that's what
everyone looked at and said, 'This is worth trying again.' '' The creators decided to keep the male
characters and to persuade Mr.~Parsons and Mr.~Galecki not to take another series in the year
between the two pilots.

They also called back one of the actresses who auditioned unsuccessfully for the original female
role: Kaley Cuoco, a former child actor who played opposite John Ritter in the comedy ``8 Simple
Rules for Dating My Teenage Daughter.'' Much of the edge was taken off the character of Penny -- so
much that at first she looked to be little more than a jiggly blonde next door with no apparent
motivation for being interested in two science geeks.

It took awhile to find the character's voice, but now Penny ``is one of the guys,'' Ms.~Cuoco said.
``She's not some untouchable creature.''

Over the first two seasons Penny and Leonard edged toward each other and are now in a full-fledged
relationship. But theirs is not the unbelievable type of couple -- a gorgeous female and a paunchy,
slacker male -- that has been so popular in Judd Apatow films recently.

``Penny has been in horrible relationships and picked the wrong guy constantly,'' Ms.~Cuoco said.
``I think she has more baggage\cite{baggage} than the guys.''

As a result, Mr.~Galecki said, ``It went from a show that I think may have made fun of intelligent
people half of the time to a show that defends intelligent people 99 percent of the time.''

The most interesting relationships are those between the two male leads and among their two boon
companions: Rajesh Koothrappali (Kunal Nayyar), an astrophysicist who is shy to the point of
muteness around women, and Howard Wolowitz (Simon Helberg), an engineer who maintains an outsize
confidence in his skill as a ladies' man, despite living with his mother.

Those did not turn the series into an immediate hit, however. ``When it went on the air, it was
disregarded almost immediately,'' Mr.~Lorre said, noting the show's respectful but not great
reviews. Then, a few weeks into its first season, came the second near-death experience -- the
writers' strike shut down production for three months. Once the strike ended, CBS moved the series
from its 8:30 time slot to 8, leading off its Monday-night lineup -- an especially tough position
for a first-year comedy.

The series stayed there in its second season, performing admirably. Then in February, on a night
when a presidential news conference interrupted its regular time slot, CBS scheduled an episode of
``Big Bang'' at 9:30, after ``Two and a Half Men.'' The ratings were so promising, said Kelly Kahl,
a CBS senior executive vice president for prime time, that the arrangement was made permanent.

This season the series has also been enjoying the publicity around the Emmy nomination for
Mr.~Parsons, as best actor in a comedy.

``It's been such a healthy climb the first two seasons,'' Mr.~Parsons said. He also said he thought
that the show had much potential to grow. ``I feel like there's still a strong segment out there
that may not be sold on the concept of four nerds and the pretty girl next door,'' he said. ``I get
that. I think there's a lot more going on that doesn't really fit in that description. It really
doesn't tell you 10 percent of why you would be interested, truly.''

\section{Another Fine Mess: Comics Whack Obama}

\lettrine{I}{s}\mycalendar{Oct.'09}{14} President Obama in trouble with his late-night comedy base?

It's likely he hasn't noticed or doesn't care. He is, after all, in the midst of his oft-invoked
``full plate'' of supposedly ``defining moments'' in his presidency -- a ``defining'' decision on
Afghanistan, ``defining'' legislative battle on health care, among other ``defining'' things.

But there is perhaps another more subtle set of ``defining'' episodes playing out for Mr.~Obama in
the televised comedy salons that had previously, by and large, been relatively gentle spaces for
him. The bits about him are getting harsher. They are no longer just gentle gibes\cite{gibe} about
Bo the dog, big ears, bad bowling and beer summits.

A conspicuous\cite{conspicuous} (if not ``defining'') episode occurred Oct.~3 on Saturday Night Live
in a skit\cite{skit} set in the Oval Office. The president (played by Fred Armisen) was defending
his record against critics who had accused him of turning the United States ``into something that
resembles the Soviet Union or Nazi Germany.'' Not so, protested the faux-Bama.

``When you look at my record, it's very clear what I've done so far,'' he said. ``And that is
nothing.''

The sketch went on to show Mr.~Obama/Armisen running through a checklist of things he had vowed to
do -- closing Guant\'anamo Bay prison, overhauling health care. All were marked, Not Done.

``Looking at this list I am seeing two big accomplishments,'' he said. ``Jack and Squat.'' And ouch.

Mr.~Obama has of course been a puzzle to comedians for some time. They agonized during the campaign
about how his low-key and confident manner did not lend itself to edgy caricature. The challenge was
made greater by the sensitivities inherent to lampooning a black candidate.

And from the outset, Mr.~Obama has been praised as someone who ``gets late night,'' whose ironic and
self-deprecating humor is well-suited to the genre's sensibilities. He was the first sitting
president to appear as a guest of Jay Leno's and David Letterman's. ``You ignore their influence at
your peril,'' said Dan Pfeiffer, the White House's deputy communications director. ``They are often
leading indicators of where the narrative is headed.''

But recent indicators could be proving ominous\cite{ominous}. There has been a
proliferation\cite{proliferation} of jokes that feed on -- or are fed by -- a
resuscitated\cite{resuscitate} old narrative against the president that goes back to last year's
campaign when both John McCain and Hillary Clinton tried to portray Mr.~Obama as an All Talk/No Walk
showboat.

Last Tuesday, Jon Stewart advanced the Saturday Night Live ``do nothing'' theme on ``The Daily
Show.'' It began as a standard Stewart video-clip juxtaposition of Mr.~Obama (and surrogates)
promising to end the military's ``Don't Ask, Don't Tell'' policy. It continued with clips from the
ensuing months of Defense Secretary Robert Gates and National Security Adviser Jim Jones saying they
had not yet gotten around to reversing ``Don't Ask, Don't Tell,'' citing Mr.~Obama's ``full plate''
of business. (``The president needs a metaphor czar,'' Mr.~Stewart said.)

What followed was Mr.~Stewart, exasperated\cite{exasperated} with a man he had supported, throwing
his hands up and essentially imploring the president to, you know, do something.

``All that stuff you've been putting on your plate?'' Mr.~Stewart said. ``It's $[$expletive$]$ chow
time, brother. That's how you get things off your plate.''

After a roar of laughter and applause from the audience, Mr.~Stewart reminded Mr.~Obama that ``You
are president of the United States.'' It sounded like something between a liberal call to action or
cry for help. As Mr.~Stewart grew more animated and the crowd grew louder, the routine took on the
feel of a televised catharsis\cite{catharsis}.

``There have been some clear shots coming across the bow from the comic left,'' observed Ric Keller,
a former Republican congressman from Florida who once wrote jokes for Jeb Bush, the former governor.

Others have noted another creeping caricature\cite{caricature} of Mr.~Obama as a ditherer. On
Thursday, for instance, Mr.~Leno joked that no one should expect the president's decision on sending
more troops to Afghanistan anytime soon. ``Remember, it took five months to decide on a puppy,'' he
said.

Jeff Nussbaum, a Democratic speech and joke writer, disagrees that late-night comedy is a leading
indicator of a cultural zeitgeist. ``To use an economic term, it is more of a lagging indicator,''
he said, something that responds to perceptions that are already entrenched. In practical terms,
President Obama has now been in office almost nine months, Mr.~Nussbaum said, and ``comedians now
have a greater body of work to go after, for better or worse.''

By and large, the bulk of late-night barbs directed at the president remain glancing at best. ``The
jokes are still largely about things like how the media lionizes Obama, or what the opposition is
saying about him,'' said Bob Lichter, of George Mason University's Center for Media and Public
Affairs, who has been tracking themes in late-night humor since 1988.

Mr.~Lichter said it was too soon to tell whether the Saturday Night Live skit is a
``harbinger\cite{harbinger} or outlier'' in how comics will treat Mr.~Obama. At the very least, it
provided a comic articulation\cite{articulation} of a potentially devastating message: ``The danger
is that Mr.~Obama is going to be defined by inaction and not living up to expectations,'' he said.

Mr.~Lichter posits Mr.~Obama's Nobel Prize as a kind of test case for how people are
perceiving\cite{perceive} him, and whether a caricature has taken hold of a man more celebrated than
accomplished. ``It will be telling to see how the comedians treat this,'' he said.

As if trying to strike pre-emptively against inevitable ridiculers, Mr.~Obama seemed eager to
embrace the ``I haven't done anything yet'' conceit in his Rose Garden remarks Friday morning. ``Let
me be clear,'' he said. ``I do not view it as a recognition of my own accomplishments.'' But it was
also striking how so many people seemed to greet the Nobel news with shock followed by laughter, as
if truth and caricature has achieved a newly seamless blend in the Obama imprint.

Really, the words rang distinctively comic: ``Did you hear that Barack Obama won the Nobel Peace
Prize?'' It sounded like the set up to a joke -- one of those jokes where the set-up itself is the
punch-line.

``That's pretty amazing, winning the Nobel Peace Prize,'' Jay Leno said, first out of the box Friday
night. ``Ironically, his biggest accomplishment as president so far \ldots winning the Nobel Peace
Prize.''

\section{Wife's Discomfort Fits Comfortably in CBS's Lineup}

\lettrine{O}{n}\mycalendar{Oct.'09}{14} last Tuesday's episode of ``The Good Wife,'' Alicia
Florrick, the disgraced politician's wife played by Julianna Margulies, put down her fork in the
middle of dinner and told her children that she needed to talk to them about something important
that happened that day, adding, ``I just would rather you hear this from me.''

Too late, she quickly discovered. Her children had already been instant-messaging with friends that
one of their childhood acquaintances was charged with murder, and that their mother was now his
lawyer. When she asked if they wanted to talk about the news she had been expecting to deliver
herself, Alicia's son answered: ``Not really. Been hearing about it all night as it is.''

Thus have the creators of ``The Good Wife,'' the new CBS drama about the spouse of a politician who
has been caught cavorting sexually and straying ethically, started to contemporize one of the
mainstays of broadcast television, the courtroom procedural.

Most of the early attention on the new series focused on its parallels to the sex and political
scandals of Mark Sanford and Eliot Spitzer. But in just three episodes this fall ``The Good Wife''
has moved well beyond that, beginning to probe more deeply into what life is like inside the 24-hour
news cycle.

``It's very difficult for Alicia, with the way culture is consumed now,'' said Robert King, who with
his wife, Michelle King, created the series. ``It's consumed in a way that is a little
cynical\cite{cynical} about people's motives. There is no generosity of spirit toward what somebody
is actually feeling.''

So at Alicia's new entry-level legal job and at her children's new school the Florricks continue to
be bombarded\cite{bombard} by the endless loop of online material: tapes of her husband, an Illinois
state's attorney, sucking a prostitute's toes; a skit on the comedy Web site FunnyorDie.com about
his indiscretions; and an online betting site that places odds on whether Alicia will file for
divorce.

``What is interesting is that everyone feels entitled to an opinion about what these women should
do,'' Michelle King said. ``I think it's just an incredibly uncomfortable place to be in. How do you
go on with your life and try to make a change, and yet still constantly be confronted with that?''

The Kings say they are intent on keeping the series current, referring to real-life situations,
people and cases. Of course that can be a liability as well: a planned story line about insider
real-estate deals surrounding Chicago's hosting of the Olympics had to be revised after Chicago
failed to win the games.

The series, which is broadcast at 10 p.m. Tuesdays on CBS, has emerged as one of the strongest hits
of the new season, attracting an average audience of 13.7 million viewers, according to Nielsen,
good enough to rank 13th among all prime-time broadcast programs (excluding football). Last week CBS
announced that it was ordering 9 additional episodes of ``The Good Wife'' to go along with the 13 it
had already ordered.

Nina Tassler, the president of CBS Entertainment, said the show benefits by attracting not only fans
of the legal-show genre but also those viewers who have found Alicia to be ``a lead character who is
identifiable and relatable.'' As a bonus, she added, ``it's a subject that's certainly in the
zeitgeist. It's art imitating life.''

To be sure ``The Good Wife'' has benefited from the strength of the programs that precede it on
Tuesday night: ``NCIS,'' the No.~1 series on television this season, and its new spinoff, ``NCIS:
Los Angeles,'' the top-rated first-year drama. But it also appears to have generated interest
because of its willingness to go beyond the normal bounds of the procedural drama that is the
backbone of the CBS prime-time schedule, adding a story arc that contains an element of
conspiracy\cite{conspiracy} and which has already expanded far beyond the simple sex scandal.

``We knew that the show had to explore the scandal or it would just turn into a dry procedural,''
Mr.~King said. ``So the scandal starts opening up other doors into Chicago politics. And Alicia
finds it's very difficult for her work to stay completely clean because of what happens there. Every
time she opens one door, it seems to open another door of scandal.''

Putting the series together presents some challenges. The writing staff is based in Los Angeles,
while the series is filmed in New York (which is made to look like Chicago). The Kings said they
meet by teleconference with the director of each episode for several hours before shooting starts in
order to review what tone they are seeking in each scene.

This is the second broadcast series for the Kings, who have worked together for 10 years. Their
first series, ``In Justice,'' about a team of lawyers who reopened the cases of people they thought
might have been wrongly convicted, lasted 13 episodes on ABC as a midseason replacement beginning in
January 2006.

And though they often sat side by side while writing the pilot episode of ``The Good Wife,'' now one
tends to make a first draft before giving the episode to the other for revisions. The series also
employs four additional writers, and Dee Johnson, an executive producer who works as a writer.

Certainly some of the success of ``The Good Wife,'' the Kings said, can be traced to Ms.~Margulies,
who expressed interest after reading the script of the pilot episode. Speaking last summer at a
meeting of television writers, Ms.~Margulies said she was drawn both by the opening scene -- the
now-familiar tableau of the humiliated wife standing at the side of the husband who is confessing
his infidelity -- and by the question of why she would stay with him.

``When I met with Michelle and Robert, I kept saying, 'I don't want to do a legal show,' '' said
Ms.~Margulies, whose previous series, ``Canterbury's Law,'' lasted just six episodes on Fox in early
2008. ``What interests me is finding out who these characters are.''

\section{Immigration Rally Draws Thousands}

\lettrine{T}{housands}\mycalendar{Oct.'09}{14} of immigrants came to Capitol Hill on Tuesday for a
day of lobbying and an afternoon rally calling for comprehensive immigration reform.

The event was timed to the unveiling of an immigration bill by Representative Luis V.~Gutierrez,
Democrat of Illinois and chairman of the Immigration Task Force of the Congressional Hispanic
Caucus.

With President Obama's stated commitment to immigration reform, advocates for immigrants said they
hoped to revive a debate that has been overshadowed by other priorities, like the economy and the
wars in Afghanistan and Iraq. As deportations continue to rise, immigration reform is needed now,
they said, to allow illegal immigrants to obtain legal status and to stop families from being torn
apart.

``We need a bill that says if you come here to hurt our communities, we will not support you; but if
you are here to work hard and to make a better life for your family, you will have the opportunity
to earn your citizenship,'' Mr.~Gutierrez said in a prepared statement. ``We need a law that says it
is un-American for a mother to be torn from her child, and it is unacceptable to undermine our work
force by driving the most vulnerable among us further into the shadows.''

Immigration overhaul faces a difficult road. President George W.~Bush twice failed to get Congress
to pass similar legislation. Mr.~Obama recently said his administration would pursue reform this
year but expected no action on legislation before 2010.

Tuesday's event was sponsored by various immigrant advocacy groups, including the Reform Immigration
for America campaign, the National Capital Immigration Coalition and Families United/Familias
Unidas. It attracted convoys of buses, vans and cars carrying more than 3,000 demonstrators from at
least 17 states.

Immigrants, religious leaders, members of Congress and immigrant advocates planned to gather on the
West Lawn for speeches and a prayer vigil at 3 p.m. Similar rallies were being held in at least 20
cities around the nation, including Los Angeles, Milwaukee, Denver and Albany.

``I'm here representing the undocumented workers who cleaned ground zero and its surrounding area
after the 9/11 terrorist attacks,'' said Rubiela Arias, 43, an illegal immigrant from Colombia who
came to Washington with an immigrant advocacy group called Make the Road New York.

Ms.~Arias described how she came from Medellín to New York in 1998 with her 5-year-old son, seeking
a safer place for her family.

``I worked for eight months cleaning the dust and debris surrounding the World Trade Center,'' said
Ms.~Arias, who cleans offices in Manhattan and was dressed in a light-blue T-shirt with a sticker
reading, ``Reform Immigration for America.'' ``There was no question about immigration status. We
were all New Yorkers; we were all Americans.''

In June, Senator Charles E.~Schumer, Democrat of New York, announced what he called seven principles
that would give form to his own reform proposal. Among them were the need to ``curtail future
illegal immigration,'' to have ``operational control of our borders'' and a ``biometric-based
employer verification system.'' Mr.~Schumer, who has been working with Senator Lindsey Graham,
Republican of South Carolina, said he would introduce a bill by Labor Day but missed that deadline.

Mr.~Gutierrez's bill, which is likely to propose less restrictive terms than Mr.~Schumer's plan for
allowing illegal immigrants to become citizens, is partly meant to pressure his Congressional
colleagues.

A main purpose of the rally was to highlight the way current immigration law splits families.

``Families deserve better than this from our government,'' said Peter Derezinski, a 17-year-old high
school senior and a United States citizen whose father was deported to Poland in April 2008 after 18
years as a truck driver and an air-conditioning repairman in Chicago. ``We need to fix our broken
immigration system so our parents who have contributed to this nation's economy in a positive way
have a chance of reuniting with their children.''

Robin Ferschke, who was traveling from Maryville, Tenn., said she planned to talk to lawmakers about
changing the law so that her daughter-in-law and grandson could live legally in the United States.
Ms.~Ferschke's son, Sgt. Michael Ferschke, a 22-year-old Marine radio operator, was killed in Iraq
in 2008, leaving his Japanese widow and their infant son in immigration limbo.

While Sergeant Ferschke was deployed to Iraq, he learned that his girlfriend was pregnant. They
decided to get married by proxy, a method that has a long history in the military when the bride and
groom cannot be in the same place for a ceremony. The boy was born in Japan and holds dual
citizenship.

But under a 1950s legal standard meant to curb marriage fraud, the wedding is not recognized for
immigration purposes even though the military recognizes the union.

``The laws we have now are inhumane and need to be changed,'' Ms.~Ferschke said. ``So I came to beg
lawmakers to change that and not force my daughter-in-law and my grandson to leave the country.''

\section{Justices Seem Sympathetic to Defendant Given Bad Advice}

\lettrine{S}{everal}\mycalendar{Oct.'09}{14} Supreme Court justices on Tuesday appeared sympathetic
to a criminal defendant who unwittingly agreed to be deported by pleading guilty to a drug crime.
But the justices seemed uncertain about whether they could fashion a legal rule that would address
extreme cases without causing turmoil\cite{turmoil} in the criminal justice system.

``Your argument has an appeal,'' Justice Samuel A.~Alito Jr.~told Stephen B.~Kinnaird, a lawyer for
the defendant, ``because removal is such a harsh consequence, particularly for someone like your
client, who had been in the United States for a long time.''

The question in the case, Padilla v. Kentucky, No.~08-651, was whether bad legal advice about a
collateral consequence of a guilty plea amounted to ineffective assistance of counsel under the
Sixth Amendment.

The defendant, Jose Padilla, has lived in the United States for 40 years, served in the Vietnam War
and is a legal permanent resident.

Mr.~Padilla, a commercial truck driver, was arrested in 2001 after the authorities in Kentucky found
more than 1,000 pounds of marijuana in his truck.

He pleaded guilty to marijuana trafficking, a felony, and received a five-year sentence.

Mr.~Padilla, a native of Honduras, later said he had agreed to the plea based on his lawyer's
incorrect advice that it would not affect his immigration status. In fact, the plea made it all but
certain that Mr.~Padilla would be deported once he served his time.

Several justices said the case presented difficult line-drawing problems, including what to do about
bad legal advice concerning the loss of a driver's license or the right to vote.

``What about advice on whether pleading guilty would cause him to lose custody of his children?''
Justice Antonin Scalia asked. ``What if pleading guilty will affect whether he can keep his truck,
which is his main source of livelihood?''

Mr.~Kinnaird said the court could limit its ruling to affirmative advice, as opposed to failures to
alert clients to potential problems, and to severe consequences.

Some justices indicated that it would be preferable to have trial judges make sure that defendants
understand the consequences of their pleas rather than requiring criminal defense lawyers to give
advice about matters outside their expertise.

Mr.~Kinnaird said that his client, had he been properly informed, would have taken his chances at
trial. He said it was not clear that prosecutors could prove Mr.~Padilla knew he was transporting
marijuana.

Wm. Robert Long Jr., an assistant attorney general in Kentucky, said Mr.~Padilla should not be
allowed to take back his guilty plea based on the bad advice he received.

``Sometimes criminal defendants risk ordinary errors with their representation,'' Mr.~Long said.

Justice Ruth Bader Ginsburg noted that an increasing number of crimes required deportation. ``There
is nothing mysterious about that,'' she said. ``There is nothing intricate about making that
determination. So why wouldn't a lawyer whose client is an alien have an obligation, when there is
an aggravated felony as the charge, to say, 'This will be the consequence'?''

The criminal defense lawyer's job, Mr.~Long responded, ``is to ensure the fair and just
determination of guilt, not to advise on collateral matters such as deportation, child custody and
the like.''

\section{Republican's Vote Lifts a Health Bill, but Hurdles Remain}

\lettrine{A}{fter}\mycalendar{Oct.'09}{14} months of relentless courting and suspense, Senator
Olympia J.~Snowe, Republican of Maine, cast her vote with Democrats on Tuesday as the Senate Finance
Committee approved legislation to remake the health care system and provide coverage to millions of
the uninsured.

With Ms.~Snowe's support, the committee backed the \$829 billion measure on a vote of 14 to 9, with
all the other Republicans opposed.

``Is this bill all that I would want?'' Ms.~Snowe said. ``Far from it. Is it all that it can be?
No.~But when history calls, history calls. And I happen to think that the consequences of inaction
dictate the urgency of Congress to take every opportunity to demonstrate its capacity to solve the
monumental issues of our time.''

Ms.~Snowe's remarks silenced the packed committee room, riveted colleagues and thrilled the White
House. President Obama had sought her vote, hoping that she would break with Republican leaders and
provide at least a veneer of bipartisanship to the bill, which he has declared his top domestic
priority.

Mr.~Obama, speaking in the Rose Garden, described the committee's action as ``a critical milestone''
and declared, ``We are now closer than ever before to passing health reform.'' But he added: ``Now
is not the time to pat ourselves on the back. Now is not the time to offer ourselves
congratulations. Now is the time to dig in and work even harder to get this done.''

With its vote Tuesday, the Finance Committee became the fifth -- and final -- Congressional panel to
approve a sweeping health care bill. The action will now move to the floors of the House and the
Senate, where the health care measures still face significant hurdles.

Aside from Ms.~Snowe, no Republicans in Congress have publicly endorsed the bills in their current
form. And Republican leaders are strongly opposed, saying the bills cost too much, raise taxes, cut
Medicare and dangerously expand federal power.

Pressure from lobbyists is sure to grow in the coming weeks. And many more lawmakers will get
involved in what promise to be impassioned and highly politicized debates in the Senate and the
House.

After the Finance Committee vote, the chief architect of the bill, Senator Max Baucus, Democrat of
Montana and chairman of the committee, declared: ``It's clear that health care reform will pass this
year. Our action today provides terrific momentum.''

Senator Charles E.~Grassley of Iowa, the senior Republican on the Finance Committee, said the bill
put the nation on ``a slippery slope toward more and more government control of health care.''

Ms.~Snowe helped write the Finance Committee bill, in months of bipartisan negotiations, but had not
committed to vote for it. She said Tuesday that she shared many of her Republican colleagues'
reservations about the legislation, and pointedly warned Democrats that they could lose her support
later in the legislative process.

``My vote today is my vote today,'' she said. ``It doesn't forecast what my vote will be tomorrow.''
And she observed, ``There are many, many miles to go in this legislative journey.''

Ms.~Snowe gave no clue how she would vote in the first few hours of committee deliberations Tuesday
and she did not alert the White House to her plans.

While colleagues spoke, she kept her head buried in papers, fidgeted and spoke occasionally with
aides. When Mr.~Baucus stepped over to speak to her, a small army of photographers snapped pictures,
with cameras clicking like a chorus of chirping crickets.

The Congressional Budget Office said the bill would cost \$829 billion over 10 years. The costs
include \$345 billion for the expansion of Medicaid and \$461 billion for subsidies to help
lower-income people buy insurance.

The budget office said the costs would be completely offset by new fees and taxes and by cutbacks in
Medicare, so federal budget deficits in the next 10 years would be \$81 billion lower than now
projected.

But Douglas W.~Elmendorf, director of the Congressional Budget Office, said his agency had not
estimated the impact of the bill on overall national health spending, public and private, and could
not say whether it would ``bend the cost curve,'' as Mr.~Obama and lawmakers want.

Likewise, Mr.~Elmendorf said he did not know for sure how the bill would affect premiums.

Several senators said they would fight for changes on the Senate floor.

Liberal Democrats, like Senator John D.~Rockefeller IV of West Virginia, said they would push for a
public insurance plan. Senators Ron Wyden of Oregon and Robert Menendez of New Jersey, both
Democrats, said they would seek changes to make insurance more affordable to middle-income families.
And Senator John Kerry of Massachusetts said he wanted to require employers to provide insurance to
their employees.

The bill does not include such an employer mandate. But employers with more than 50 workers would
have to reimburse the government for some or all of the cost of federal subsidies provided to
employees who buy insurance on their own.

Ms.~Snowe said she liked the Finance Committee bill because it would prohibit insurance companies
from discriminating against people on account of health status or sex and would create a network of
insurance exchanges where individuals, families and small businesses could shop for coverage, with
subsidies from the federal government.

At the same time, Ms.~Snowe said she shared Republican ``concerns about vast governmental
bureaucracies and governmental intrusions.'' That, she said, is why she had opposed amendments to
create a government insurance plan and would continue to do so.

Ms.~Snowe said she was open to a compromise under which a public plan could be ``triggered'' in
states where people could not otherwise find affordable insurance. She said her ``paramount
concern'' was that insurance might be too expensive for some people, even with government subsidies.

The Congressional Budget Office said the Finance Committee bill would provide coverage to 29 million
people, but still leave 25 million uninsured in 2019. Of those left uncovered, about a third would
be illegal immigrants.

\section{Study Finds Pro and Cons to Prostate Surgeries}

\lettrine{P}{rostate}\mycalendar{Oct.'09}{14} cancer patients who chose minimally invasive surgery
rather than more extensive operations to remove the prostate were less likely to experience
complications like pneumonia, but reported higher rates of long-term problems, including impotence
and incontinence, according to one of the largest studies to compare outcomes to date.

Patients achieved similar rates of cancer control regardless of which surgery they had, the analysis
found.

The study, in Wednesday's Journal of the American Medical Association, is not considered the last
word on the subject, several experts agreed. But it raises questions about claims of superiority of
minimally invasive laparoscopic and robotic-assisted surgeries, which have led to a surge in their
popularity.

``People intuitively think that a minimally invasive approach has fewer complications, even in the
absence of data,'' said Dr.~Jim C.~Hu, the study's lead author, who is director of urologic robotic
and minimally invasive surgery at Brigham and Women's Hospital in Boston. ``Men who were well
educated and had higher incomes were actually more likely to embrace this approach, often due to
aggressive marketing by hospitals that had spent \$1.5 million to acquire the robots. I think the
technology has been oversold.''

In one version of prostate removal, called open surgery, a surgeon makes an incision that is several
inches long. With minimally invasive surgery, also called laparoscopic surgery, the surgeon operates
through a series of small incisions using tools and a camera for the operation. With robotic surgery
the surgeon sits at a computer and manipulates a robot to do the operation through the small
openings.

In 2003, minimally invasive radical prostatectomies, which include robotic surgeries, made up fewer
than 10 percent of prostate removal surgeries. By 2006-7, they constituted 43 percent of procedures.

The Harvard researchers who did the study assessed the outcomes of 1,938 men who had minimally
invasive prostate surgery from 2003 to 2007 and 6,899 men who had open surgery. They used
Surveillance, Epidemiology and End Results, or SEER, data from the National Cancer Institute
representing 26 percent of the American population, linking it with Medicare data.

The men in the study -- all of them 65 or older -- who underwent minimally invasive surgery had
shorter hospital stays, fewer respiratory complications and other surgical complications, and were
far less likely to receive a blood transfusion. But they had more complications involving the
genital and urinary organs immediately after surgery, with 4.7 percent having those complications,
compared with 2.1 percent of open surgery patients.

When the researchers looked at lasting complications more than 18 months later, they found that men
who had minimally invasive surgery were at greater risk of suffering from incontinence and erectile
dysfunction than those who had open surgery.

For each 100 men who had minimally invasive surgery, some 15.9 percent were at risk of being
incontinent each year, while 26.8 percent experienced erectile dysfunction, compared with 12.2
percent and 19.2 percent, respectively, each year for every 100 men who had open surgery, the study
calculated.

Several surgeons who specialize in robot-assisted procedures said the study was limited because it
was unable to distinguish between those using robot technologies and older minimally invasive
techniques.

Many experts said the outcomes of experienced surgeons were better than those reported in the study.

``I almost exclusively do robotic prostatectomy now because I think that, despite this manuscript,
there is clear evidence that it is comparable, in terms of continence, potency and tumor control,''
said Dr.~Joseph Smith, the chairman of urologic surgery at Vanderbilt University School of Medicine.

But Dr.~Smith added, ``I don't think there's anything demonstrating it to be superior.''

Dr.~Peter Scardino, chief of surgery at Memorial Sloan-Kettering, said the study was important
because it reported on data that did not come just from one medical center or one region.

``At the end of the day,'' Dr.~Scardino said, ``what all the studies will show is that it's not the
tools the doctor uses, but the experience and skill of the surgeon. There's nothing magical about
the laparoscopic or robotic.''

\section{Hybrid Cars May Include Fake Vroom for Safety}

\lettrine{F}{or}\mycalendar{Oct.'09}{14} decades, automakers have been on a quest to make cars
quieter: an auto that purrs, and glides almost silently in traffic.

They have finally succeeded. Plug-in hybrid and electric cars, it turns out, not only reduce air
pollution, they cut noise pollution as well with their whisper-quiet motors. But that has created a
different problem. They aren't noisy enough.

So safety experts, worried that hybrids pose a threat if pedestrians, children and others can't hear
them approaching, want automakers to supply some digitally enhanced vroom. Indeed, just as
cellphones have ring tones, ``car tones'' may not be far behind -- an option for owners of electric
vehicles to choose thesound their cars emit.

Working with Hollywood special-effects wizards, some hybrid auto companies have started tinkering in
sound studios, rather than machine shops, to customize engine noises. The Fisker Karma, an \$87,900
plug-in hybrid expected to go on sale next year, will emit a sound -- pumped out of speakers in the
bumpers -- that the company founder, Henrik Fisker, describes as ``a cross between a starship and a
Formula One car.''

Nissan is also consulting with the film industry on sounds that could be emitted by its forthcoming
Leaf battery-electric vehicle, while Toyota has been working with the National Highway Traffic
Safety Administration, the National Federation of the Blind and the Society of Automotive Engineers
on sounds for electric vehicles.

``One possibility is choosing your own noise,'' said Nathalie Bauters, a spokeswoman for BMW's Mini
division, who added that such technology could be added to one of BMW's electric vehicles in the
future.

The notion that battery E.V.'s and plug-in hybrids might be too quiet has gained backing in
Congress, among federal regulators and on the Internet. The Pedestrian Safety Enhancement Act of
2009, introduced early this year, would require a federal safety standard to protect pedestrians
from ultra-quiet cars.

Karen Aldana, a spokeswoman for traffic safety agency, which is also working on the issue, said,
``We're looking at data on noise and E.V. safety, but manufacturers are starting to address it
voluntarily.''

A Toyota spokesman, John Hanson, said: ``I don't know of any injuries related to this, but it is a
concern. We are moving rapidly toward broader use of electrification in vehicles, and it's a fact
that these cars are very quiet and could pose a risk to unsighted people.''

A study published last year by the University of California, Riverside and financed by the National
Federation of the Blind evaluated the effect of sounds emitted by hybrid and internal-combustion
cars traveling at 5 miles per hour.

People listening in a lab could correctly detect a gas-powered car's approach when it was 28 feet
away, but could not hear the arrival of a hybrid operating in silent battery mode until it was only
seven feet away.

Some electric-vehicle drivers have taken a low-tech approach to alerting pedestrians. When Paul
Scott of Santa Monica, Calif., drives his 2002 Toyota RAV4 electric car, he often rolls down the
windows along busy streets and turns up his radio so people know his virtually silent vehicle is
there.

Mr.~Scott, vice president of the advocacy group Plug In America, said he would prefer giving drivers
control over whether the motor makes noise, unlike, say, the Fisker Karma, which will make its
warning noise automatically.

``Quiet cars need to stay quiet -- we worked so hard to make them that way,'' he said. ``It's the
driver's responsibility not to hit somebody.''

Mr.~Scott has already warmed up to the idea of a car ring tone.

``It should be a manually operated noisemaker, a button on the steering wheel triggering a recording
of your choice,'' he said. ``It could play 'In-a-Gadda-Da-Vida,' or anything you like.''

\section{Still Hoping to Sell Music by the Month}

\lettrine{T}{he}\mycalendar{Oct.'09}{14} idea of selling monthly subscriptions to a vast catalog of
online music has met with only limited success. That isn't stopping a new batch of entrepreneurs
from trying to make it work.

The latest and perhaps most surprising entrants to the field are the European entrepreneurs Niklas
Zennstrom and Janus Friis. In 2001, they created and financed Kazaa, one of the original
peer-to-peer file-sharing services that hurt the music industry. The two have created and financed a
secretive start-up called Rdio, with offices in Los Angeles and San Francisco.

Rdio and similar start-ups are reinventing a concept pioneered earlier this decade by Rhapsody, a
service majority-owned by RealNetworks, and the tamed version of Napster, now owned by Best Buy. A
few hundred thousand Rhapsody and Napster subscribers pay monthly fees of around \$15 for the right
to stream an unlimited number of songs, at any time, from their PCs and mobile devices.

But with modest membership growth at best, neither service has managed to challenge iTunes, with its
many millions of users -- or enticed music lovers from pirating music. Moreover, Yahoo, AOL and MTV
Networks have abandoned their own music subscription efforts.

But as CD sales continue to plummet, and the music industry searches for a profitable future,
entrepreneurs with various approaches say they believe they can finally make music subscriptions
work. Rdio is hoping to introduce a music subscription service by early next year that offers
seamless access to music from both PCs and cellphones. The big challenge will be to get licenses
from the major music labels, which have not viewed past digital music efforts by Mr.~Zennstrom and
Mr.~Friis favorably.

``The ironies are very interesting,'' said Drew Larner, Rdio's chief executive, who says talks with
music labels are continuing and confidential.

Since they started and sold Kazaa years ago, the founders ``have shown they understand content and
they have always been up front with the labels about what they are trying to do.''

Mark Piibe, the head of digital business development at the EMI music label, confirmed that talks
were under way with Rdio and said there was no reluctance to deal with the pair. ``They're
businessmen with a real track record of innovation,'' he said. ``They are bringing a lot of new
ideas to music distribution and there is no reason why we wouldn't talk to them seriously.''

Two other new music subscription companies will most likely beat Rdio to the United States market.
Spotify, a year-old company that offers a free, ad-supported service but tries to get users to sign
up for about \$16-a-month ad-free version, already has several million users in Europe.

Spotify says it hopes to bring the service to the United States early next year, but it could look
somewhat different here. American music labels are increasingly resistant to the idea of licensing
their catalogs to any new service offering free music with ads, because they have already backed
free music downloading sites like MySpace Music and Imeem.

``We like Spotify as our partner in Europe, but we would like them to move more toward a paid
subscription environment,'' said Thomas Hesse, president of global digital business at Sony Music.

That raises the possibility that a lesser known company will lead the wave of new music subscription
services -- Mog, a three-year-old blogging network that recently raised \$5 million, led by the
venture capital firm Menlo Ventures. Mog has licensing deals with all four major American music
labels: the Universal Music Group, Sony Music Entertainment, the Warner Music Group and EMI Music,
plus thousands of independent labels through the digital distributors Orchard and Ioda.

Based in Berkeley, Calif., Mog is getting specific about its plans for its ``All Access'' service,
which it says it will introduce by Thanksgiving. For \$5 a month, members can listen to as much
music as they want from their computer; for \$12 to \$15, users can access music on their mobile
devices as well.

The service is a sort of cross between Pandora, the free online radio site, and music subscription
services like Rhapsody. Users can listen to the songs of any artist and build a playlist with only
that music. By manipulating a slide bar, users can also gradually add a smattering of tunes by
similar artists, in effect customizing their own online radio station. People can then share their
playlists from their site on the Mog service.

David Hyman, Mog's founder and the former chief executive of Gracenote, calls it ``radio without
restriction.''

\section{China's Growing Economy Mints Billionaires}

\lettrine{T}{he}\mycalendar{Oct.'09}{14} superrich in China have bounced back from the financial
crisis with a vengeance, and China now has more known U.S.-dollar billionaires than any other
country except the United States, according to a report released Tuesday.

The annual Hurun Report said that China has 130 known dollar billionaires, up from 101 last year.
The number in the United States is 359, while Russia has 32 and India 24, according to Forbes
magazine.

China's rich are getting richer, with the average wealth on the list \$571 million, up almost
one-third from last year, said Rupert Hoogewerf, the report's compiler. The Hurun Report is a luxury
publishing and events group, its Web site says.

``With the greatest wealth destruction in the West of the last 70 years, we've seen China buck the
trend and the wealth seems to be still growing,'' Mr.~Hoogewerf said in an interview on the
sidelines of an event to unveil the 2009 rich list.

``They've put the credit crunch behind them,'' he said. ``The key driver has been urbanization.
You've got all these cities being built, and that requires property developers, iron and steel
manufacturers. The latest thing is cars.''

Topping the list was Wang Chuanfu, chairman of the electric car and battery maker BYD, in which the
American billionaire Warren Buffett holds a stake. Mr.~Wang's personal wealth is estimated at \$5.1
billion. He was also the fastest riser from last year, moving up 102 places.

Second place went to Zhang Yin and family, owners of the paper recycler Nine Dragons Paper, while in
third place was Xu Rongmao and family, the owners of the Shimao property group.

Huang Guangyu, who founded Gome Electrical Appliances and owns unlisted property businesses, sank to
17th place from the top position he held last year. He is being investigated for alleged financial
irregularities.

One famous name fell off the list this year -- the basketball player Yao Ming, who has struggled
with a foot injury for the past few months.

China's ruling Communist Party once condemned entrepreneurs and private business people as
capitalist exploiters but has welcomed them since the late reformist leader, Deng Xiaoping, began
landmark economic reforms in the 1970s.

One-third of the people on the 1,000-name Hurun list are estimated to be Party members, according to
the report.

Mr.~Hoogewerf said the actual number of dollar billionaires could be higher than estimated.

``Either they are superdiscreet, or perhaps they haven't come to the surface,'' he said. He added
that the transparency of wealth was, however, now much higher because of the greater number of
listed companies.

Mr.~Hoogewerf said people who probably should have been listed, but about whose wealth not enough is
known, included Liu Chuanzhi, the chairman of Lenovo, one of the world's largest makers of PCs, and
Chen Feng, the founder of Hainan Airlines.

\section{In Recession, China Solidifies Its Lead in Global Trade}

\lettrine{W}{ith}\mycalendar{Oct.'09}{14} the global recession making consumers and businesses more
price-conscious, China is grabbing market share from its export competitors, solidifying a dominance
in world trade that many economists say could last long after any economic recovery.

China's exports this year have already vaulted it past Germany to become the world's biggest
exporter. Now, those market share gains are threatening to increase trade frictions with the United
States and Europe. The European Commission proposed on Tuesday to extend antidumping duties on
Chinese, as well as Vietnamese, shoe imports.

China is winning a larger piece of a shrinking pie. Although world trade declined this year because
of the recession, consumers are demanding lower-priced goods and Beijing, determined to keep its
export machine humming, is finding a way to deliver.

The country's factories are aggressively reducing prices -- allowing China to gain ground in old
markets and make inroads in new ones.

The most striking gains have come in the United States, where China has displaced Canada this year
as the largest supplier of imports.

In the first seven months of 2008, just under 15 percent of American imports came from China. Over
the same period this year, 19 percent did. Meanwhile, Canada's share of American imports fell to
14.5 percent, from nearly 17 percent.

Besides increasing its share of many American markets, China is increasing the value of exports in
absolute terms in some categories. In knit apparel, for instance, American imports from China jumped
10 percent through July of this year -- while America's imports from Mexico, Honduras, Guatemala and
El Salvador plunged 19 to 24 percent in each country, according to Global Trade Information
Services.

Indeed, China's said Wednesday that its export slump eased in September, down just 15 percent, a
strong improvement over August's decline. Economists said it was a sign of an improving global
economy and renewed strength for Chinese exporters.

A similar tale is told around the world, from Japan to Italy.

One reason is the ability of Chinese manufacturers to quickly slash prices by reducing wages and
other costs in production zones that often rely on migrant workers. Factory managers here say
American buyers are demanding they do just that.

``The buyers are getting more and more tough in bargaining for lower prices, especially American
buyers,'' says Liao Yuan, the head of international trade at the Changrun Garment Company, which is
based in southern China and exports jeans to Europe and the United States. ``They offer \$2.85 per
pair of jeans for a package of a dozen, when the reasonable price is \$7.''

Because China produces a diversified portfolio of low-priced and essential items, analysts say the
country's exports can hold up relatively well in a recession. Few other countries can match what has
come to be called the ``China Price.''

``China has a huge advantage,'' says Nicholas R.~Lardy, an economist at the Peterson Institute for
International Economics in Washington. ``They can adjust to market changes very rapidly. They have
flexibility in their labor markets. And as consumers trade down the quality ladder, China can
benefit.''

The expiration of textile quotas in large parts of the world this year has also allowed China to
increase its market penetration.

But equally important are government policies that support this country's export sector -- from
Beijing keeping its currency weak against the dollar to its determination to subsidize exporters
through tax credits and billions of dollars in low-interest loans from state-run banks.

The results have been impressive. All told, in the first half of 2009, China exported \$521 billion
worth of clothes, toys, electronics, grains and other commodities to the rest of the world.

Though that represented a 22 percent decrease from the first half of 2008, it compares favorably to
other major exporters. German exports have fallen 34 percent over the same period. Japanese exports
were down 37 percent and American exports 24 percent, according to Global Trade Information
Services.

Trading powerhouses like Germany are suffering from weaker demand for heavy equipment, automobiles
and luxury goods. But the value of exports from oil-producing countries, like Russia and Saudi
Arabia, has fallen even more.

One reason is that the price of oil has plummeted from last year's record highs. But since oil is
priced in dollars and the value of the dollar has fallen markedly, so have the value of American
imports from these countries -- over 45 percent in the case of Russia's exports to the world.

Meanwhile, American imports from Saudi Arabia have fallen 65 percent.

China's market share gains are mostly at the expense of countries like Japan, Italy, Canada, Mexico
and Central America -- in industries that China has long sought to dominate.

China's share of furniture imports in the United States has grown to 54 percent, from 50 percent,
over the last year, while furniture exports to the United States from Canada and Italy have plunged
40 percent from a year ago. In Europe, Chinese textiles and apparel have gained market share in
every major country, after the quota expiration in January. Not long ago, Italy's shoe imports were
dominated by Romania; now China has a commanding share.

Japan once relied on electronics shipments to the United States, but every year for the past decade
Japan has lost market share to China. This year is no different. In 1999, electronics goods from
Japan made up 18 percent of America's electronics imports. Today, that figure is down to 7 percent.

China's market share has climbed 10 to 20 percent from a year ago. Together, the gains are helping
China maintain its large trade surplus with the rest of the world, reviving worries about global
trade imbalances -- and once again putting the spotlight on China's currency, the renminbi.
Specialists note, however, that much of China's gains stem from the fact that it is increasingly
assembling products whose components are made elsewhere, and re-exporting them.

After letting its currency rise against the dollar, beginning in July 2005, China is once again
pegging it closely to the dollar. As the dollar has fallen against other major currencies like the
euro -- about 15 percent since a year ago -- Chinese imports have become more and more competitive.

Now, European officials are clamoring for China to reduce its flood of exports and pressing for
antidumping investigations.

The International Monetary Fund is calling on China to rebalance its economy and allow its currency
to appreciate against other major currencies.

The United States -- which for years complained about China's weak currency and soaring trade
imbalances -- has largely been silent in recent months, analysts say, partly because Washington is
trying to improve relations with Beijing at a time when it desperately needs China to purchase
American debt.

``Obama's interest is not to push China to appreciate the currency, but to get them to pay the
bills,'' Dong Tao, an economist at Credit Suisse says, referring to China's purchases of American
debt.

For its part, Beijing worries that raising the value of its currency could be catastrophic, damaging
exports and diluting the effect of the government's aggressive stimulus package.

But the country's leaders are well aware of the need to shift the economy away from heavy dependence
on exports and toward stronger domestic consumption. Indeed, China is eager to move up the value
chain, by selling higher-priced goods like computer chips, aircraft and pharmaceuticals -- all of
which would bring better-paying jobs and healthier economic growth.

Moreover, many economists say that as Chinese consumers become richer, they will buy more of their
own goods. And as the dollar falls, it will make American exports more competitive globally,
including in China. Those trends together could eventually help rebalance global trade -- which
became overly reliant on Americans buying cheap Chinese goods and China buying American debt. Right
now, Beijing worries about growing trade frictions with its biggest trading partners, the European
Union and the United States, and the possibility of some countries initiating protectionist
measures.

Chinese exporters, meanwhile, fear that even as they gain market share the pressure to produce at
low prices will hurt them and the quality of their products.

Ms.~Liao at the Changrun factory says many producers are essentially scavenging to source raw
material.

``Some even go to old factories to collect abandoned fabrics from old stock, so they can save
two-thirds of the cost on raw material,'' she says. ``These fabrics are in very bad shape. They
won't wash, and easily wear out.''

But the discounting period may be here for a while with many economists forecasting a lengthy period
of slow growth in Europe and the United States.

``China is going to get stronger,'' Mr.~Tao at Credit Suisse says. ``Its competitors are getting
weaker in the downturn. And the Chinese state has helped bail out some industries, like the auto
industry; so in the future some new industries may emerge as exporters.''

\section{Exiled From School, H.I.V.-Infected Orphans Learn a Bitter Lesson}

\lettrine{A}{n}\mycalendar{Oct.'09}{14} NHON TAY, Vietnam -- The first day of school was a special
one last month for the 15 children from the Mai Hoa orphanage here. They are infected with H.I.V.,
the virus that causes AIDS, and for the first time they would be allowed to attend the local primary
school.

``The children were so excited,'' said Sister Nguyen Thi Bao, who runs the orphanage and had been
lobbying for three years to enroll them in the government school. ``They had been wishing for this
day to come.''

But when they arrived, they found an uprising by the parents of the other students, who refused to
let their children enter the school together with the infected orphans. Some of the parents hastily
backed away when the orphans walked past.

After a short standoff, the principal, who had agreed to accept the orphans, told Sister Bao that
their papers were not in order and that they could not stay.

The children returned to the orphanage, just a short walk down a country road, where they continue
to study in small classrooms, still exiled from the uninfected world.

``I was so happy to go to the school,'' said a 12-year-old fourth grader for whom Sister Bao
insisted on anonymity to keep her from the spotlight. ``But then I saw that some parents wouldn't
let their children go to school with me because they are scared of my disease.''

The girl said she understood their reaction.

``If I were a normal child, I would be afraid, too, because I wouldn't understand,'' she said. ``I
would feel the same way. But I wouldn't have acted the way they did.''

Sister Bao and officials of the district and the school, the An Nhon Dong Elementary School, have
met with the parents since then, but they remain adamant.

``I don't want my child to be with the AIDS children,'' Nguyen Thi Thuy, 36, said the other day as
she brought her 8-year-old son to school. ``He could be injured, and it's easy to transmit the
disease through blood. And once you're sick, it's difficult to become a normal person again.''

One after another, parents who arrived with their children on small motorbikes raised their voices
in agreement. If the orphans came back, said a man who gave his name only as Tam, he would pull his
son out again.

The story is not surprising, said Eamonn Murphy, Vietnam director for Unaids, the United Nations'
AIDS-fighting agency.

``You go to any rural environment in Asia, and you are going to have similar reactions,'' he said.
``The general lack of understanding leads to this inappropriate reaction and fear.''

Most of the parents here are farmers with little education, but the prejudice seemed to extend to
city folk as well.

``I don't know why we don't isolate people with AIDS,'' said a civil servant in Ho Chi Minh City,
about 20 miles southeast of the village. ``Even with swine flu we isolate people, and this disease
is much more dangerous.''

There is no truth to these fears, Mr.~Murphy said.

``H.I.V. is not contagious from community contact, even if you are sharing cups and saucers and
eating from a communal plate. You can't get H.I.V. from that.''

In recent years, Vietnam's prevention and treatment programs have been improving, Mr.~Murphy said,
although so far only 30 percent of people who need life-saving antiretroviral drugs receive them.

About 290,000 people in Vietnam, a country of 86 million, carry H.I.V. today, and Mr.~Murphy said
that although the rate of increase was slowing, the infection was spreading outside high-risk
groups.

Among those infected, the government estimates that 5,100 are children. Although the law requires
equal treatment, almost none of them have been accepted in schools because of the fears of other
children's parents, Nguyen Vinh Hien, the deputy minister of education, said last month.

He said the ministry would try to enroll at least half of these children in government schools by
next year, but the experience of the Mai Hoa orphans suggests that this will not be easy.

Frightened and angry on that first day, some of the parents seemed heartless, Sister Bao said.

``They were saying the children were going to die anyway, so there's no need for them to study,''
she said. `` 'If they are going to study, let them do it in the orphanage, and not put our children
in danger.' ''

The Mai Hoa AIDS Center, with its green and quiet grounds, was founded by a Roman Catholic order in
2003 as a hospice for patients in the final stages of the disease. It added the orphanage to care
for children of people who died there.

The children are infected as well, Sister Bao said, but are receiving antiretroviral medication.

The buildings behind the classrooms are still a hospice, where a dozen emaciated patients lie on
cots. Altogether, 250 people have died, Sister Bao said, including 90 whose unclaimed ashes are
stored behind the hospice buildings.

Some of those are the remains of the children's parents.

So the orphans of Mai Hoa live suspended between the death that fills the space behind their
classrooms and the life of a world, just down the road, that still will not accept them.

``The children say they want to go to the other school because they want to have friends,'' Sister
Bao said.

But the 12-year-old fourth-grader seemed to have changed her mind.

``I don't want to go to that school,'' she said. ``I already have enough friends here.''

\section{Inside the Islamic Emirate}

\lettrine{A}{}\mycalendar{Oct.'09}{19} young Taliban driver with shoulder-length hair got behind the
wheel of the car. Glancing at me suspiciously in the rearview mirror, he started the engine and
began driving down the left-hand side of the road.

It was some sort of prank\cite{prank}, I hoped, some jihadi version of chicken -- the game where two
drivers speed toward each other in the same lane until one loses his nerve.

Which lane he drove down showed what country we were in. If he continued driving on the left, we had
crossed into Pakistan. If he drove on the right, we were still in Afghanistan.

A mile down the road, traffic signs appeared in Urdu.

We're in Pakistan, I thought to myself. We're dead.

Eight days earlier, a Taliban faction had kidnapped me along with an Afghan journalist, Tahir
Luddin, and our driver, Asad Mangal, during a reporting trip just outside Kabul. The faction's
commander, a man who called himself Atiqullah, had lied to us. He had said we were being moved to
southern Afghanistan and would be freed.

Instead, on Nov. 18, we arrived in Pakistan's tribal areas, an isolated belt of Taliban-controlled
territory. We were now in ``the Islamic emirate'' -- the fundamentalist state that existed in
Afghanistan before the 2001 American-led invasion. The loss of thousands of Afghan, Pakistani and
American lives and billions in American aid had merely moved it a few miles east, not eliminated it.

Through seven years of reporting in the region, I had pitied captives imprisoned here. It was
arguably the worst place on earth to be an American hostage. The United States government had
virtually no influence and was utterly despised.

Since 2004, dozens of missiles fired by American drones had killed hundreds of militants and
civilians. The Taliban had held Afghan, Pakistani and foreign hostages in the area for years,
trading lives for ransom and executions for publicity.

``We're in Pakistan,'' I said out loud in the car, venting my anger.

Atiqullah laughed, and the driver appeared surprised.

``How does he know it's Pakistan?'' the driver asked.

``Because you're driving down the left-hand side of the road,'' I answered.

``How do you know that?'' he asked. ``When were you in Pakistan before?''

Atiqullah smiled and appeared amused by the conversation. He knew I had been to Pakistan many times
on reporting trips.

I was one of dozens of journalists who had written articles detailing how Al Qaeda and the Taliban
had turned the tribal areas into their new stronghold after being driven from Afghanistan in 2001. I
had watched the Pakistani government, then led by President Pervez Musharraf, largely stand by as
the Taliban murdered tribal elders and seized control of the area.

Now, an abstract foreign policy issue was deeply personal. When my wife and family learned that I
was in the tribal areas, their distress would increase exponentially. They would expect that I would
never return.

We arrived in a large town, and I spotted a sign that said ``Wana'' in English. Wana is the capital
of South Waziristan, the most radical area of the seven administrative districts that make up the
tribal areas. We stopped in the main bazaar, and I was left alone in the car with the young driver.

Desperate rationalizations swirled through my mind. Our captors wanted a ransom and prisoners.
Killing us got them nothing. The three of us would survive. They were all delusions, of course.
Simply getting us this far was an enormous victory for them. We would be held here for months or
killed.

Outside the car, dozens of Pakistani tribesmen and Afghan and foreign militants milled around. Each
carried a Kalashnikov assault rifle on his shoulder and had a long, thick beard.

A man with a large turban stopped, peered at me in the back seat and asked the driver a question in
Pashto. The driver looked at me and said a sentence that I thought included the word for martyr. I
told myself the driver had said I was on my way to heaven.

Atiqullah got back into the car, and I felt relief. He had kidnapped us, but more and more I
desperately viewed Atiqullah as my protector, the man who would continue to treat us well as other
militants called for our heads.

Our first Pakistani home was in Miram Shah, the capital of North Waziristan. Two large sleeping
rooms looked out on a small courtyard. One even had a small washroom, separate from the toilet, for
showering.

On the first day there, I went to the bathroom and returned to find Tahir with a fresh cut on his
calf. It looked as if someone had drawn a line across his leg in red ink. A local Waziri militant
had taken out his knife and tried to cut off a chunk of Tahir's calf, saying he wanted to eat the
flesh of an Afghan who worked with Westerners. One of Atiqullah's guards had stopped him.

All day, a parade of random Pakistani militants stopped by the house to stare at us. I felt like an
animal in a zoo. Among them was a local Taliban commander who introduced himself as Badruddin. He
was the brother of Sirajuddin Haqqani, who led the Haqqani network, one of the most powerful Taliban
factions in the region. Miram Shah was its stronghold.

Their father was Jalaluddin Haqqani, an Afghan mujahedeen leader whom the United States and Pakistan
backed in the 1980s when he battled the Soviets. In the 1990s, the United States ended its
relationship with the Haqqanis and many other hard-line Afghan fighters. With Pakistan's help, the
Taliban movement emerged and the Haqqanis joined them.

Badruddin, a tall, talkative man who appeared to be in his early 30s, said he was preparing to make
a video of us to release to the media. He smiled as he showed me a video on his camera of a French
aid worker, Dany Egreteau, who had been kidnapped a week before us as he walked to his office in
Kabul. He was in chains and appeared to have welts on his face. He implored his family and friends
to save him.

``It's a nightmare,'' he said. ``I really beg you to pay.''

I asked if Tahir and I could speak alone with Atiqullah, and I told him we should not make the
video. The American and Afghan governments were more likely to agree to a secret prisoner exchange,
I said, than a public one.

Trying to reduce their expectations, I told him it would be far easier to get prisoners from the
main Afghan-run prison outside Kabul, known as Pul-i-Charkhi. If the Taliban demanded prisoners from
the American-run detention centers at Guant\'anamo Bay, Cuba, and Bagram, Afghanistan, they would
never succeed.

I was not worth that much, I told him, and he should compromise. I did not say it, but I also wanted
to spare my family the pain of seeing me in a video. To my surprise, Atiqullah agreed.

``I am one of those kinds of people,'' he said at one point. ``I am one of those people who like to
meet in the middle.''

Tahir, Asad and I would be allowed to call our families that night to prove we were alive, he said.
Atiqullah told me to emphasize during the call that he wanted to reach a deal quickly. He continued
to cover his face with a scarf. To me, that meant he did not want to be identified because he
planned to release us.

I spent the rest of the day nervously scribbling a list of things I wanted to say to my wife,
Kristen, whom I had married just two months earlier. I added items and then crossed them out. I
wanted to ease my family's fears that I was being tortured, but I also wanted to do everything
possible to free the three of us. I wasn't sure I would have another chance to speak with her.

Late that night, Atiqullah and Badruddin drove us out of town. Atiqullah stopped the car in a dry
riverbed and turned off the engine. He left the headlights on, and we used them to see the number
pad on a small satellite phone. Atiqullah and Badruddin ordered me to tell Kristen that we were
being held in terrible conditions in the mountains of Afghanistan. I dialed my wife's number.

``Hello?'' she said.

``Kristen?'' I said. ``Kristen?''

``David,'' she said, ``it's Kristen. I love you.''

She sounded calm.

``Kristen?'' I asked.

``Yes?'' she said.

``I love you, too,'' I said. ``Write these things down, O.K.?''

``O.K.,'' she said.

She sounded remarkably composed.

``I'm, we are being treated well,'' I said.

``Being treated well,'' Kristen repeated.

``No.~1,'' I said.

``Uh-huh, No.~1,'' Kristen said.

``No.~2,'' I said. ``Deal for all three of us, all three of us, not just me. The driver and the
translator also; it has to be a deal for all three of us.''

``Deal for all three of us,'' she repeated. ``The driver and the translator as well. O.K.''

``Do not use force to try to get us,'' I said.

``Do not use force,'' Kristen repeated.

``Four,'' I said.

``Yes,'' Kristen said.

``Make a deal now or they will make it public,'' I said. ``They want to put a video out to the
media.''

Kristen repeated my words back to me.

``It will make it a big political problem,'' I said.

Atiqullah told me to tell her that this was my last call.

``They said I can't call you again,'' I said. ``They want a deal now and I can't call you again.''

``You cannot call me again,'' she repeated. ``I love you. I love you, honey.''

``I love you, too,'' I said. ``Tell my family I'm sorry.''

``Your family is here, Lee's here with me,'' she said, referring to my older brother.

``I'm sorry,'' I said, ``I'm sorry.''

``It's going to be all right,'' Kristen said calmly. ``I love you. I am praying for you every day.''

Kristen said she wanted to make sure she understood what the Taliban wanted.

``What is the deal?'' she asked.

Atiqullah told me to tell her that he would call The New York Times's Kabul bureau with demands.
``We are very concerned about you,'' Kristen said. ``And we love you, and we're praying for you.''

The satellite phone beeped and abruptly went dead. Kristen was gone.

Standing in the remote darkness of Waziristan at the mercy of Taliban militants, I felt at peace. I
had spoken to my wife for the first time in nine days. I had expected panic or tears, but she
sounded collected and confident. Her words ``It's going to be all right'' would linger in my mind
for months. Her composure would sustain me.

Atiqullah and Badruddin then told me to call The Times's bureau in Kabul. But instead of ordering me
to make specific demands, they instructed all three of us to exaggerate our suffering.

``We are in terrible conditions, Tahir is very sick,'' I told Chris Chivers, a close friend and
Times reporter, who answered the phone. I was ordered to tell Chris that Atiqullah was not with us
-- even though he was, in fact, standing beside me.

Tahir then spoke to Chris and asked him to tell his family he was alive and in good health.

``They keep telling me that if things go wrong they will repeat the story of Helmand,'' Tahir said,
``so I am just afraid they are going to kill me.''

Tahir was referring to the 2007 kidnapping of an Italian journalist in Helmand Province that ended
in the beheadings of an Afghan journalist and a driver working with him.

Asad then spoke with an Afghan reporter in the bureau.

``I am fine, I am O.K.,'' he said. ``Tell my family that we are in the mountains but we are O.K.''

The conversation dragged on, with Atiqullah continuing to direct me about what to say. When he
ordered me to tell Chris that they would kill the driver and translator first, I refused.

``Kill me first,'' I told Atiqullah, ``Kill me first.''

Chris overheard me and interrupted. ``Nobody needs that, David,'' he said. ``Nobody needs to die.''

``They are threatening to kill the driver and the translator,'' I explained to Chris. ``I have to
tell you, I have to tell you. I don't want to tell you.''

``We understand that they are making those threats,'' Chris said, ``but that will not make our job
easier.''

Chris said that if the Taliban killed anyone it would make government officials angry and make any
deal even more difficult.

``Please don't let them kill the driver and the translator,'' I said. ``Please don't let them kill
the driver and translator.''

``I am sorry about this,'' I added. ``I apologize to everyone.''

``David, this is not your fault,'' Chris said. He urged me to tell Atiqullah to keep calling.

``O.K., all three of us, Chris,'' I said as Badruddin and Atiqullah ordered me to end the call.
``It's gotta be all three of us. I gotta go.''

As Atiqullah drove us back into Miram Shah, I felt relief. Kristen had sounded calm. Chris had said
The Times was doing all it could. I felt I had fought for Tahir and Asad.

We arrived at a new house, and I was again surprised by the good conditions. It had regular
electricity, and we could wash ourselves with buckets of warm water. I received a new set of
clothes, a toothbrush, toothpaste and shampoo. Guards allowed us to walk in a yard, and the weather
was surprisingly warm. We received pomegranates and other fresh food and Nestl\'e Pure Life water
bottled in Pakistan.

The tribal areas were more developed and the Taliban more sophisticated than I expected. They
browsed the Internet and listened to hourly news updates on Azadi Radio, a station run by the
American government. But then they dismissed whatever information did not meet their preconceptions.

Atiqullah said he needed to return to Afghanistan, but two of his men stayed behind to guard us. ``I
will return in 7 to 10 days,'' he promised, then disappeared.

That week, to help us pass the time, we received a shortwave radio and a board game called checkah,
a Pakistani variation of Parcheesi. To my amazement, the guards even brought me English-language
Pakistani newspapers. Delivered to a shop in Miram Shah, the newspapers were only a day or two old.
Instead of beating us as I expected, our captors were at least trying to meet some of our needs.

But as in so much of our seven months in captivity, reasons for optimism would be overtaken by harsh
realities.

For the next several nights, a stream of Haqqani commanders overflowing with hatred for the United
States and Israel visited us, unleashing blistering critiques that would continue throughout our
captivity.

Some of their comments were factual. They said large numbers of civilians had been killed in
Afghanistan, Iraq and the Palestinian territories in aerial bombings. Muslim prisoners had been
physically abused and sexually humiliated in Iraq. Scores of men had been detained in Cuba and
Afghanistan for up to seven years without charges.

To Americans, these episodes were aberrations. To my captors, they were proof that the United States
was a hypocritical and duplicitous power that flouted international law.

When I told them I was an innocent civilian who should be released, they responded that the United
States had held and tortured Muslims in secret detention centers for years. Commanders said they
themselves had been imprisoned, their families ignorant of their fate. Why, they asked, should they
treat me differently?

Other accusations were paranoid and delusional. Seven years after 9/11, they continued to insist
that the attacks were hatched by American and Israeli intelligence agencies to create a pretext for
the United States to enslave the Muslim world. They said the United States was forcibly converting
vast numbers of Muslims to Christianity. American and NATO soldiers, they believed, were making
Afghan women work as prostitutes on military bases.

Their hatred for the United States seemed boundless.

Ten days passed, but Atiqullah did not return as promised. Badruddin now seemed to be in charge.

He moved us to a far smaller, dirtier house. The space we were allowed to walk in was the width of a
city sidewalk and ringed by high walls. The food was unclean and made me sick.

Our first night there, the Taliban commander who owned the house promised to update us every three
days on negotiations for our release. But we would not see him again for months. The guards stopped
taking Tahir to a local doctor for digestive and skin ailments.

And it was increasingly clear that Tahir and Asad would be separated from their families for Id
al-Adha -- a major Muslim holiday that marks Abraham's willingness to sacrifice his son to show his
devotion to God.

Alarmed by the worsening treatment, Tahir and I began a hunger strike in early December. At first,
the guards panicked and begged us to eat. We refused.

After two days, the guards said Atiqullah had called and told them that a deal for our release was
nearly complete. He said he was waiting for approval from President Hamid Karzai of Afghanistan, who
was on a foreign trip. The French aid worker in the video I was shown had been released, they said.
We would be next.

Fearing that our continued defiance would anger our captors and scuttle the deal, we began eating
again.

Instead of releasing us, Badruddin moved us to yet another house. It was larger than the previous
one but felt more like a prison. Twenty-foot-high concrete walls surrounded a small courtyard where
I spent my days walking in circles. For the first time in my life, I began praying several times a
day, and I found that it centered me.

We started preparing our own meals. The food was cleaner and fresher, but cooking for ourselves gave
a worrying sense of permanence to our imprisonment.

Badruddin visited us several days later and promised that negotiations were continuing. But he was
increasingly casual. Any sense of urgency about our release seemed to be fading. Before leaving, he
told me the Taliban would not kill me.

``You are the golden hen,'' he said, clearly expecting me to lay a golden egg.

I asked him to promise not to kill Tahir and Asad. Speaking directly to me in broken English, he
said the Taliban had decided to kill Asad if their demands were not met in a week. After he left,
Tahir and I decided not to tell Asad.

I panicked over the next two days, frantically trying to think of ways to save our young driver.
Since the three of us had arrived in Pakistan on Nov. 18, I had spent hours each day talking
politics, religion and survival with Tahir, but I could barely communicate with Asad.

I spoke little Pashto, he spoke little English. I came up with a routine when a newspaper arrived. I
showed Asad photos and tried to explain what they were about. He laughed, but I felt like a monster.
Asad was an impoverished, hard-working father of two -- and I was going to get him killed.

On the third day after Badruddin's visit, I told one of our guards that I was willing to make a
video -- or do anything they wanted -- to save Asad. The guard said he would check with Badruddin.
The following day, the guard announced that it had all been a misunderstanding. There was no
deadline to kill Asad. I didn't know what the truth was but felt enormous relief.

Several days later, Badruddin arrived to make the video. He promised us that it would go only to our
families, but what he instructed us to say made me think it would be released publicly. As guards
pointed assault rifles at our heads, I called for President Bush and President-elect Obama to meet
the Taliban's demands.

``If you don't meet their demands,'' I said, ``they will kill all of us.''

Tahir and Asad then made similar statements. Badruddin departed, and I told myself that our families
would at least know we were alive.

As December dragged on, tensions in the house steadily grew. Qari, the guard who had nearly shot
Tahir, tore the checkah board to shreds after he repeatedly lost. Then, Tahir and Asad ripped up two
other checkah boards out of frustration as well. Qari began spending hours alone reciting the Koran
and seemed increasingly distant and unstable. I worried that the situation was slowly spinning out
of control.

Several days before Christmas, Atiqullah finally returned. He announced that he had spectacular
news. ``We are here to free you,'' he said, wearing no scarf over his face for the first time. ``We
have come here to release you.''

At first, I was euphoric. My confidence in Atiqullah had not been misplaced. Here was a more
moderate and reasonable Taliban leader who would persevere and release us.

Then, later that night, the conversation turned menacing.

The American military had mounted an operation to arrest Abu Tayyeb on the morning that we were to
interview him, Atiqullah said, referring to the Taliban leader we had been traveling to meet when we
were kidnapped.

Shocked, I told Atiqullah I knew nothing about a military operation.

I had sent text messages from my cellphone to Saudi Arabia before the interview, Atiqullah claimed,
to tip off the American military about Abu Tayyeb's location. Again, I told him I had no idea what
he was talking about.

Finally, he announced that I was a spy, along with other employees of The Times in Afghanistan. His
men had prepared a suicide attack on the paper's Kabul bureau, he said, which he could set off with
a single phone call. His men had nearly kidnapped Carlotta Gall, our bureau chief, but she had left
an interview just before they arrived.

``She was probably given information,'' he said, seemingly convinced that all journalists were
intelligence operatives.

Our imprisonment, I thought, had reached a low point. My colleagues in Kabul were now in danger.
Atiqullah's talk of our imminent release seemed farcical\cite{farcical}.

The following morning, Atiqullah insisted that there was, in fact, a deal. At one point, he said we
would be exchanged within ``days.'' He toyed with me, asking which flights I would take back to the
United States and how many television cameras would be at the airport. He asked me what I would say
to my wife when I saw her.

By this point, I began to doubt everything he said. Then I learned that he had lied to us from the
beginning.

In conversations when our guards left the room, Tahir and Asad each separately whispered to me that
Atiqullah was, in fact, Abu Tayyeb. They had known since the day we were kidnapped, they said, but
dared not tell me. They asked me to stay silent as well. Abu Tayyeb had vowed to behead them if they
revealed his true identity.

Abu Tayyeb had invited us to an interview, betrayed us and then pretended that he was a commander
named Atiqullah.

I was despondent and left with only one certainty: We had no savior among the Taliban.

\section{At Book Fair, a Subplot About Chinese Rights}

\lettrine{A}{s}\mycalendar{Oct.'09}{19} China extends its economic reach, it has also increased
efforts to promote its culture, or ``soft power,'' to counter Western influence and improve its
image in the wider world.

Yet if Chinese goods are accepted everywhere, its arts and literature, embattled at home after
decades of censorship and state control, are proving harder for the government to export.

After years of delicate preparations, China was the ``honored guest'' this past week at the
Frankfurt Book Fair, the largest and most influential book trade event, based on the number of
publishers represented. But what Beijing hoped would be a celebration of its cultural achievements
turned into a tug of war between control and free speech, as much a showcase for Chinese dissidents
as the state's approved writers.

Mao Zedong said that power flowed from the ``wielders of the pen,'' not only from the gun. But the
chairman would not be amused to find books like ``Mao: The Unknown Story,'' an indictment of his
rule that is banned in China, displayed alongside the official Chinese exhibit at this year's fair,
which ended Sunday.

When the German organizers and diplomats urged the Chinese to allow a prominent storyteller and
musician, Liao Yiwu, to come to Frankfurt, the authorities refused to lift his overseas travel ban,
and told him to stop talking about it.

A symposium\cite{symposium} preceding the book fair titled ``China and the World -- Perceptions and
Realities,'' became a major confrontation. Fair organizers withdrew invitations to two dissident
writers the Chinese wanted to exclude, Dai Qing and Bei Ling, but welcomed them at the last minute
after criticism by journalists and politicians. When the writers made statements, the Chinese
delegation walked out, only to return after an abject apology by the fair's director, Jürgen Boos.

``We did not come to be instructed about democracy,'' declared Mei Zhaorong, China's former
ambassador to Germany.

Unlike the exquisitely choreographed ceremonies during the Beijing Olympics, the fair presented a
messier and more ambiguous portrait of China on the rise -- a country still deeply uncomfortable
with its own discordant voices, yet eager to become more competitive with the West in the realm of
ideas.

China controlled its own massive display of books, artwork and authors at the fair, including even
books from Taiwan, to underline its assertion of ``One China.'' But it could not censor the 2,500
books about China displayed by others. And while Beijing had many consultations with the German
government and arguments with the fair organizers, it ultimately did not push to prevent dissidents
and critics -- even representatives of the Dalai Lama -- from attending the event.

The book fair is not the Beijing Olympics and ``cannot be controlled,'' said Mr.~Boos. He apologized
for mishandling the symposium, but said: ``It is the beginning of a cultural dialogue. And dialogue
is not easy.''

Still, Chinese officials did not attend dissident events, ``which were full of people who already
agreed with the dissidents,'' said the German novelist Tanja Kinkel. ``They were preaching to the
choir,'' she said.

The Chinese themselves were annoyed. With SpiegelOnline headlining its coverage ``China, the
Unwelcome Guest,'' several official Chinese delegates told colleagues that Europe's politicians and
news media were strongly biased.

Li Pengyi, a delegation member and vice president of China Publishing Group Corporation, said
happily that China had sold nearly 900 copyrights here. But he complained about the coverage.

``We don't feel we've been hospitably treated,'' he said. ``China sent more than 2,000 people to
Frankfurt. And now this barrage of criticism.''

Zhao Haiyun, spokesman for China's General Administration of Press and Publication, said that
instead of focusing on literature, the media had focused on human rights and censorship. ``The
German media are very biased,'' he said.

Even so, the Chinese did not pull out. The Beijing leadership sent Xi Jinping, China's vice
president and heir apparent to President Hu Jintao, a measure of the political weight they attached
to the event.

Michael Naumann, a former German culture minister and now publisher and editor of Die Zeit, a
prominent weekly newspaper, said German organizers misjudged the complications of honoring China in
a year laden with controversy, including the 20th anniversary of the fall of the Berlin Wall, the
20th anniversary of the crushed Tiananmen Square democracy movement and the 60th anniversary of
Chinese Communist Party rule.

``I think the people who run the book fair were kind of na\"ive when they invited the Chinese,'' he
said. ``But opening this enormous window of the book fair to Chinese writers, whether they are
censored or not, will give them a way to sniff out the open forum of intellectual debate.''

Since 2004, China has pursued what it calls its ``going out'' policy on the cultural front, trying
to square its economic influence and new status as a global power, while trying to defuse criticism
on issues like Tibet, Taiwan and human rights.

There have been yearlong cultural exchanges with many countries; the opening of hundreds of language
teaching centers known as Confucius Institutes; new foreign-language services from official media
like Xinhua and CCTV; and new interest in foreign platforms like the Kennedy Center and the
Europalia festival in Brussels.

There have been other furors. When China was featured at the 2004 Paris Book Fair, officials
initially persuaded the French not to invite the Nobel literature laureate Gao Xingjian, a French
citizen whose books are banned in China.

But Frankfurt, with its 7,300 publishers and 300,000 visitors, was a much riskier venture.

Jing Bartz has been the fair's chief representative in Beijing since 2003 and negotiated strenuously
with Chinese publication officials. ``China has really wanted to use this platform to promote
Chinese culture,'' she said. ``On the other side, they are worried because they can't use Chinese
rules to do it.''

What helped persuade China was the cultural trade gap. At the 2005 Beijing book fair, the Chinese
were shocked that German publishers sold 600 copyrighted works to China while the Chinese sold just
one to Germany, Mrs.~Bartz said.

Chinese officials worried particularly that the Dalai Lama might attend, or that books would be
displayed from adversaries like the banned movement Falun Gong.

The breakthrough came in 2006, said Mrs.~Bartz, when Shi Zongyuan, then head of the General
Administration of Press and Publication, told organizers: ``We just have to make it very clear what
is our guest of honor program, and what are the other events.''

China invested \$15 million and managed nearly every detail of its exhibition. There was much
argument over what translations to finance. The 20 new German-published volumes China financed
include works by major writers, like Jiang Rong's ``Wolf Totem,'' Yu Hua's ``Brothers,'' and Xu
Zechen's ``Running Through Zhongguancun.''

Mr.~Xu's hit, about a migrant hawking pirated DVDs and fake IDs in the capital, was unexpected. But
of some 100 newly translated titles that China promoted, most are banal introductions to China from
state publishers.

``The government has not put on such a concentrated, large-scale event before to promote Chinese
literature, so I think it's a good opportunity,'' said Mr.~Xu, 31. ``Because of the government's
involvement, there are inevitably going to be these ideological problems. But we just have to be
responsible to ourselves.''

Since the uproar over the symposium last month, said Mr.~Boos and Mrs.~Bartz, China has appeared
more relaxed. Officials eventually gave up protesting the attendance of those like the Uighur
independence advocate Rebiya Kadeer; the Dalai Lama's envoy, Kelsang Gyaltsen; Ms.~Dai, Mr.~Bei or
Mr.~Gao.

``They tried to learn,'' Mrs.~Bartz said. But she confirmed that while the Chinese were ``very
satisfied with the business results'' of the fair, ``they don't really feel they were welcomed as
guests here.'' The word went down from the top, she said, not to react to demonstrations or
provocations from protesters or journalists.

Back in China, however, the fair has not brought any noticeable easing of restrictions.

Mr.~Liao, the writer and musician, was imprisoned from 1990 to 1994 after he wrote a poem about the
Tiananmen massacre. Despite an invitation here -- he hoped to promote his book about China's
downtrodden, known in English as ``The Corpse Walker'' -- the police would not lift a ban on his
going overseas.

In a telephone interview, Mr.~Liao said it was not a complete loss for him or other underground
writers, given the publicity. ``Only by going through these incidents, it seems, can we become known
to the outside world,'' he said.

\section{Party Elder Still Jousts\cite{joust} With China's Censors}

\lettrine{F}{or}\mycalendar{Oct.'09}{19} nearly two decades, the Communist Party strove to wipe out
the national memory of Zhao Ziyang, the reform-minded party secretary who opposed the use of force
against pro-democracy protesters in 1989.

So when a former aide of Mr.~Zhao's, Du Daozheng, disclosed in May that he had helped secretly
record Mr.~Zhao's memoir for posthumous publication, Mr.~Du's daughter refused to let him walk
outside alone for fear of possible repercussions.

She need not have worried. On June 25, a top official in charge of propaganda showed up at Mr.~Du's
western Beijing apartment with a reassuring message from Zhongnanhai, the headquarters of the
Communist Party and the government. Mr.~Du said he was told that, as an old friend of Mr.~Zhao's,
``Zhongnanhai and party central can understand why you did this.''

Mr.~Du used to be among those who delivered such judgments. Until he was ousted in 1989 with
Mr.~Zhao, he served as head of the government's press and publications administration, an agency
that helps enforce censors' orders.

Now he spends his days jousting with such officials, trying to foist unmentionable topics like
Mr.~Zhao's career into the public domain. Helping with Mr.~Zhao's memoir -- a rare look at the
party's inner conflicts that was published this May outside China -- was a particularly daring
thrust.

But strategic ventures into forbidden territory are characteristic of his monthly scholarly journal,
Yanhuang Chunqiu. In 2005, he published articles on Hu Yaobang, the former party leader whose death
helped set off the Tiananmen protests. Infuriated authorities threatened to reduce copies of the
magazine to pulp, according to Mr.~Du's daughter, Du Mingming.

After a string of journal articles last year touched on Mr.~Zhao's accomplishments, party
authorities issued an internal regulation so precisely focused that it could have been named after
Mr.~Du. The order forbids retired government or party officials to serve as publication directors.

Party sources say Jiang Zemin, the now-retired leader who replaced Mr.~Zhao, was irritated by the
articles and instigated\cite{instigate} the pressure on Mr.~Du to step down. Sitting in the
magazine's musty offices, Mr.~Du said he dealt with the order by reshuffling titles.

``I just ignore it,'' he said. ``I am old enough and tough enough that if there is any pressure from
the government, I can hold on here.''

MR. Du survives such skirmishes because he is 86, wily and quietly supported by certain party
luminaries. He says as many as 100 former party officials back his magazine's attempts to draw
lessons from the party's buried past and nudge it toward democratic reforms. Some current officials
also sympathize with the effort, he suggests. ``Nobody dares close it,'' he said, lest that provoke
a reaction from ``old cadres.'' Last year supporters promised him, ``If the magazine closes, we will
take to the streets,'' he said.

They said: ``We are old. We are in our 80s. We have heart problems. We will probably die in the
streets.''

``So the conservatives don't take any action,'' Mr.~Du said, ``because they are afraid of that
responsibility.''

Others suggest the party can afford to be tolerant. Russell Leigh Moses, a Beijing-based analyst of
China politics, said that Mr.~Du and other liberal-minded party ``elders'' posed no particular
threat to today's Communist Party, so slaps on the wrists sufficed.

``I admire the courage and the conviction, but the conservatives really won this battle some time
ago,'' he said. ``I really see him as a tragic figure, still holding the flag after most of the
armies have left the field.

``He is fighting a struggle against the political tenor of the times, as well as against time
itself,'' Mr.~Moses said.

Mr.~Du is not, however, fighting with himself. He sees his modest magazine, printed on
newsprint-quality paper and distributed to some 100,000 subscribers for about a dollar a copy, as
``the best thing he has done in his life,'' his daughter said.

The struggle between truth and propaganda has been a constant theme in Mr.~Du's life. He was an
early Communist Party loyalist, dropping out of middle school at age 14 to join the battle against
Japanese invaders. After the Communists rose to power in 1949, he dutifully -- and falsely --
reported the party's claims of record harvests and free food as a reporter for Xinhua, the state-run
news agency.

But by April 1959, he could no longer reconcile the discrepancy with reality. In a 4,900-word letter
to a superior, he documented widespread famine and disease in the countryside.

Within two months, his letter was turned against him during the campaign against antiparty
``rightists.'' He was publicly condemned 17 times -- once before an auditorium filled with 6,000
people -- and dismissed from his job and party post.

He escaped even worse punishment, as many did, by betraying others to his government tormentors. The
four people he named ``suffered greatly,'' he said. One refused until his dying day to forgive him.

Mr.~Du was persecuted anew during the Cultural Revolution. ``It was as though I was sent to hell and
back,'' he wrote in his magazine in January.

Once rehabilitated, he rose quickly in the ranks of the state-run media. He held a vice ministerial
position during the 1989 pro-democracy protests.

Mr.~Du tried to marshal support for Mr.~Zhao's position that the crisis should be resolved
peacefully. In a letter signed by some 30 other government officials, Mr.~Du urged party leaders not
to use force against the demonstrators, according to his daughter.

When the troops opened fire on the night of June 3, he, his wife and some friends ``cried from our
hearts'' in his living room, Mr.~Du said.

``We all shared one feeling,'' he said. ``The Communist Party is over.''

IT was three years before he could discuss the tragedy with Mr.~Zhao, who was removed as party
secretary and placed under house arrest. Mr.~Du, who was also ousted, said he urged Mr.~Zhao to
record his version of events for history's sake.

In the ensuing years, Mr.~Du said, he filled two notebooks with Mr.~Zhao's words, then switched to a
tape recorder. Four other former officials also pitched in. Mr.~Du hid his copies of the tapes in
his daughter's underwear drawer and later she transported the copies to Hong Kong. After Mr.~Zhao,
still under house arrest, died in 2005, Bao Pu, the son of Mr.~Zhao's top aide, began transcribing
and translating other copies of the recordings.

Ever the strategist, Mr.~Du recommended that the memoir be published only after the 20th anniversary
of the crackdown in June. But with the support of Mr.~Zhao's family, in May Mr.~Bao arranged
publication of a Chinese version in Hong Kong and an English version, titled ``Prisoner of the
State: The Secret Journal of Premier Zhao Ziyang.''

Hong Kong bookshops have reportedly sold 100,000 copies. ``I have not seen such excitement about a
book in years,'' said Lam Mingkei, owner of Causeway Bay Bookstore, a prominent bookseller in Hong
Kong.

Mr.~Du said he believed that the democratic ideals expressed in Mr.~Zhao's book and in the pages of
his beloved journal would eventually take hold, though not, he predicted, under the current leaders.
``If the Communist Party refuses to take political reform, then there must be some other force that
rises up to carry it out,'' he said.

In the meantime, he says, he will defend his journal's role as a liberal windsock. Said his
daughter: ``My father knows how to fight.''

\section{Beijing's Air Is Cleaner, but Far From Clean}

\lettrine{T}{his}\mycalendar{Oct.'09}{19} city's network of arterials\cite{arterial}, five ring
roads that are fed by nine more freeways, is barely two decades old, but it is already
sclerotic\cite{sclerotic}. Roughly four million vehicles clog Beijing roads, seven times the number
about 15 years ago. On any given day, another 1,500 new vehicles join the crush.

So it is no surprise that Beijing has some of the worst air pollution of any big city on earth.

No, the surprise is this: Beijing's air is actually getting cleaner.

China may have a hard-earned reputation for long-neglected and fearsome environmental problems, from
poisoned rivers to chemical-belching smelters. But the nation's capital, Beijing, is trying hard to
clean up its dirty air.

The results show up not only in recent pitch-perfect October days, but in the data that for years
have presented a bleak picture of pollution here.

Through September, the government counted 221 days in which the 0-to-500 pollution index -- the
lower the number, the better -- was below 101. It was the greatest number of ``blue-sky days,'' as
the city calls them, since daily measurements were first published in 1998.

At the same time, the city has recorded only 2 days with dangerously high air pollution. That is the
lowest number in a decade, and fully 17 days fewer than were logged in the same period in 2000.

Outside experts caution that the city's measurements are not just imprecise -- they do not measure
the tiniest particulates that are most damaging to lungs, for instance -- but potentially
misleading.

The government index is a 24-hour average of readings from monitors citywide. In contrast, a single
monitor operated by the United States Embassy in downtown Beijing, where traffic is heavy,
consistently reports less healthy air conditions than the official index.

And Beijing's air remains far from pristine\cite{pristine} by any measure. The average concentration
of particulates in city air during 2008, for example, was six times the ideal standard recommended
by the World Health Organization. Indeed, Beijing has yet to meet the W.H.O.'s interim air standards
for developing countries -- or even the less stringent standards posted by China's national
government.

Yet if the air remains bad, experts say, the trend toward clearer skies is both steady and
undeniable.

``For those of us who have been monitoring air pollutants for about 10 years, we see a clear
reduction in pollution,'' Zhu Tong, a professor and air pollution scientist at Peking University's
College of Environmental Sciences and Engineering, said in an interview.

In the past decade, in fact, authorities have moved against air pollution problems with a tenacity
that some environmentalists in developed nations, pitted against industry lobbyists and balky
political machinery, can only envy\cite{envy}.

Consider: since 1999, the city has ratcheted up its new-car emissions standards from nonexistent to
the level of Euro IV, the same clean-tailpipe requirements that are now enforced across the European
Union nations. The rest of China will move to Euro IV next year.

With the switch to Euro IV, Beijing environmental officials last year ordered a 90 percent reduction
in the sulfur content of gasoline and diesel fuel, as well as significant new cuts in polluting
compounds like benzene and aromatic hydrocarbons in gasoline.

Nor is that all. More than 4,100 of the 20,000 city buses run on clean-burning compressed or
liquefied natural gas -- the largest such fleet in the world. The Beijing subway system, currently
about 125 miles of tunnels and overhead tracks, is undergoing a breakneck expansion that will nearly
triple its length in the next five years.

Nor is that all, either. Before last year's Olympics, the government imposed a regimen based on even
and odd license-plate numbers that effectively banned half of all private automobiles from the road
on weekdays. Officials -- and the public -- liked it so much that a modified version, banning one in
five cars, is now a permanent rule.

The city also bans heavy trucks from entering the city during the day; after dark, the 18-wheelers
can be seen lining up on suburban freeways, waiting for permission to enter. The oldest, dirtiest
automobiles -- called ``yellow-label'' cars, after the sticker glued to their windshields -- are
banned from the center city altogether.

And in July, 11 national government ministries took a leaf from the Obama administration's Cash for
Clunkers program and began offering rebates of \$440 to \$880 to people who trade in old
high-pollution cars and trucks for new ones. Officials estimate that the program will sweep about
2.7 million dirty autos from the roads nationwide.

Impressive as all that is, it may not be enough. While Beijing's air is cleaner now than in the
recent past, the improvements have been largely, though not totally, offset by a sevenfold increase
in the number of vehicles since 1993.

Many of the easy cleanup measures, like improving fuel quality and banning the dirtiest vehicles,
have already been made. Future improvements will be harder to attain -- and yet the flood of cars
onto the streets is only growing.

The government insists that it is winning the race. Because older autos emit as much as 30 times the
pollution of new ones, officials say, scrapping a relatively few clunkers makes room for many new
replacements. Moreover, improvements in gasoline and diesel fuel quality have lowered pollution from
all vehicles, including the older ones.

But the city's cleanup drive is not stopping at its highways.

In the past five years, Beijing has converted 60,000 boilers and commercial heaters to run on clean
natural gas instead of coal. The city's four coal-fired power plants -- all of them major polluters
-- have installed state-of-the-art pollution scrubbers.

Nearly 2,900 gas stations and petroleum storage tanks have been equipped with recycling controls.
Hundreds of heavily polluting factories have been moved from central Beijing, including a coking
coal plant and a steel mill that is scheduled to depart by the end of 2010.

The mill is a prime example of environmental action, China-style. A unit of the Shougang Group,
China's fourth largest steelmaker, it has been one of Beijing's biggest employers. When its 140-mile
move to neighboring Hebei Province is completed next year, 65,000 Beijing workers will lose their
jobs.

On the other hand, thousands more are getting a better life. Consider Wang Shuqin, a 64-year-old
resident of one of Beijing's century-old hutongs, or alleys, who recently got a new electric space
heater for her two-bedroom apartment.

Like most all hutong denizens, Mrs.~Wang and her husband, Liu Xingye, have weathered decades of
Beijing winters by feeding round bricks of coal into a squat pot-bellied stove. Every day they
emptied the stove's ashes, turning their bedroom floor black. Every night they worried about
carbon-monoxide poisoning.

Every year, they burned 1,200 one-kilogram coal bricks -- one and one-third tons of coal -- to stay
warm. Until now: this month, Beijing's city government gave the couple a two-thirds discount on the
electric heater, and a laughably low nighttime rate for electric power, 3 cents a kilowatt-hour.

Since 2004, Beijing has replaced 94,000 pot-bellied coal stoves with efficient electric heaters,
eliminating the filth that came from chimneys burning roughly 100,000 tons of coal a year.

Because of that and other changes, Beijing's sulfur emissions have been cut almost in half compared
with 2000, and industrial dust has decreased by two-thirds.

``It's cleaner. It's much safer,'' Mrs.~Wang said, flashing a thumbs up. ``Now we no longer worry
about getting poisoned.''

Increasingly, neither does the rest of Beijing.

\section{Software Pirates in China Beat Microsoft to the Punch}

\lettrine{A}{t}\mycalendar{Oct.'09}{19} shops in the bustling Xinyang market in Shanghai, fake Apple
iPhones and Bose speakers were displayed alongside bootleg copies of Microsoft's Windows 7 operating
system, a week before it officially was to go on sale.

``Which version do you want? Ultimate? Normal? English or Chinese?'' one shopkeeper asked, proudly
pointing out her ample supply of discs packed in unmarked white boxes.

People in mainland China have been able to buy pirated copies of the newest version of Microsoft's
Windows franchise this month for just 20 yuan, or \$2.93, each -- a fraction of list prices, which
are as high as \$320.

Windows 7's ``early release'' in China underscores the challenge major software makers face in
trying to make money in China, the world's second-largest PC market, after the United States.

The U.S.~research firm IDC estimated that about 80 percent of software sold in China last year was
pirated. While that figure is falling, it is still double the global average and about four times
the average in developed markets like the United States and Japan.

``The big issue that is driving piracy in China today is price,'' said Matthew Cheung, an analyst at
the research firm Gartner. ``If you're trying to sell a program that costs 2,000 yuan to a student
living on 400 yuan a month, that's simply not going to work out for most consumers.''

In a nod to such pressures, Microsoft cut the price of its Office 2007 Home and Student Edition to
199 yuan last year from 699 yuan.

And Microsoft will sell its low-end Windows 7 Home Basic version for 399 yuan, a modest price by
Western standards but still 15 times as much as is charged for pirated copies.

Violation of intellectual property rights has been a sore spot in China's relations with its major
trading partners, even as it has cracked down on rampant piracy of everything from Gucci bags to
software.

A Chinese court jailed four people in August for spreading a bootleg version of Microsoft's Windows
XP, in what Xinhua, the official news agency, called the nation's biggest software piracy bust.

``A lot of people are used to getting away with it for a long time,'' said Steve Vickers, president
of FTI-International Risk. ``There are signs that law enforcement is picking up, and that should
help things improve.''

Business Software Alliance, a trade association created by the software industry, said the sector
had lost more than \$6.6 billion in China last year to piracy.

Most experts agree that piracy in China is a long-term issue, but many say that conditions should
improve as software makers cut prices, users become more educated and living standards rise.

``Piracy in China is reducing year by year because the government is placing more attention on it
and prices between the real and fake have narrowed,'' said Qian Liyong, director of the E.U.-China
Project on the protection of intellectual property rights, based in Beijing.

Gartner estimated that software piracy rates in mainland China would fall as low as 50 percent by
2012, putting it almost on a par with rates in developed Asian markets like Hong Kong.

Customer education is also improving in mainland China as many realize the dangers of installing
pirated software, which sometimes comes with viruses and spyware.

``This is a long-term issue -- 10-20 years -- it's not just going to go away in an instant,'' said
Edward Yu, chief executive of Analysys International, a research firm.

In a bid to tackle the problem, Microsoft began an unconventional campaign in China last year that
caused a black screen to be displayed every hour for users of pirated versions of Windows XP.

But that just caused thousands of irate users to migrate to free software from domestic companies
like Kingsoft.

Some say that free Web-based software, supported by advertising, may ultimately help to reduce
piracy in China by letting third parties pay for development costs.

``Because of the Internet, we are seeing a trend that software is by and large becoming free for
consumers from point to point,'' Mr.~Yu said.

\section{In Face of Sanctions, China Premier Warms to Iran}

\lettrine{T}{he}\mycalendar{Oct.'09}{19} Obama administration's effort to marshal\cite{marshal}
global support for new sanctions against Iran hit another snag on Thursday, when China said it was
seeking to increase cooperation and high-level exchanges with Iran, suggesting that it would not
support additional punitive measures over Tehran's nuclear program.

The comments by Prime Minister Wen Jiabao came a day after Russia's prime minister, Vladimir
V.~Putin, said that sanctions were ``premature'' and two days after its foreign minister, Sergey
V.~Lavrov, said that threatening Iran while talks were under way would be ``counterproductive.''

China and Russia are crucial to a muscular sanctions campaign, analysts said, because of their
geopolitical and commercial links to Iran and because they have historically resisted sanctions in
the United Nations Security Council, where each holds a veto.

But China has a deepening reliance on Iran for energy, and in a meeting with Iran's first vice
president, Mohammed Reza Rahimi, in the Great Hall of the People, Mr.~Wen emphasized cooperation
between the countries in energy and trade and greater coordination in international affairs,
according to the official Xinhua news agency.

Obama administration officials said they had long viewed China as being even more reluctant than
Russia to act against Iran. China and Russia often act in unison, the officials said, with China
typically following Russia's lead in Security Council deliberations.

Still, the administration pointed to a recent statement signed by the two countries, as well as the
United States, Britain, France and Germany, that threatened Iran with ``consequences'' if it did not
negotiate over its nuclear program. Those countries presented the same message to Iran at talks in
Geneva on Oct.~1.

Speaking of the Iranians, Ben Rhodes, a spokesman for the National Security Council, said, ``They've
heard directly in those meetings in Geneva that all the parties involved expect them to take
action.''

The administration was also buoyed\cite{buoy} by what it viewed as a shift in Russia's position when
President Dmitri A.~Medvedev said after a meeting with President Obama at the United Nations last
month that ``sanctions rarely lead to positive results, but sometimes, sanctions are inevitable.''
He reiterated his views to Secretary of State Hillary Rodham Clinton in Moscow on Tuesday, American
officials said.

Yet some analysts said the administration should pay more attention to the first part of
Mr.~Medvedev's statement than to the second. Russia has always had little faith in the value of
sanctions, and his statement reaffirmed that skepticism.

``The administration overreacted to Medvedev's comments,'' said David J.~Kramer, a trans-Atlantic
fellow at the German Marshall Fund. ``Putin's comments are the ones to focus on. And I would bet
that Lavrov made sure his comments were in sync with Putin's.''

The latest statements from Russia and China appear to expose a rift\cite{rift} with the United
States about how to press Iran in the negotiations. The United States wants to back up diplomacy
with a credible threat of sanctions; Russia and China believe such an approach would backfire by
pushing Iran into a corner.

China will continue to play a ``constructive role'' in finding a peaceful solution, according to
Xinhua. Russia says the focus should be solely on diplomacy.

On Monday, Iranian officials will meet in Vienna with officials of several countries to discuss
details of an accord reached in Geneva to ship most of Iran's publicly declared stockpile of lightly
enriched uranium to Russia, where it would be further enriched. It would then be returned to Iran,
where it would fuel a research reactor in Tehran.

Iranian officials were in Beijing to attend a meeting of a regional security group that includes
China and Russia. China has used the group to strengthen its ties.

This week, Mr.~Wen and Mr.~Putin attended the signing of a framework agreement by a state-run
Chinese energy company and Gazprom of Russia to potentially build two pipelines carrying natural gas
from Siberia to China.

China's rapid growth has forced it to look at many foreign nations for energy. It sees Iran as a
large provider of oil, and has ignored Western calls for sanctions or withdrawing business from
other resource-rich nations.

It imports half its oil from the Middle East, and those shipping routes are secured by the United
States. It will have to balance its desire to lock up energy supplies with its interest in ensuring
those supplies can be shipped to China.

``The question is whether they can do this skillfully, because they haven't been playing this global
game for decades,'' said Jon B.~Alterman, the Middle East program director at the Center for
Strategic and International Studies.

\section{Students' Advocate Was Once Convicts' Last Hope}

\lettrine{A}{tlanta}\mycalendar{Oct.'09}{19} -- ``Pick your head up, buddy,'' Tom Dunn said to
Darius Nash, who had fallen asleep during the morning's reading drills. ``Sabrieon, sit down,
buddy,'' he called to a wandering boy. ``Focus.''

Mr.~Dunn's classroom is less than three miles from his old law office, where he struggled to keep
death row prisoners from the executioner's needle. This summer, after serving hundreds of death row
clients for 20 grinding, stressful years, he traded the courthouse for Martin Luther King Jr.~Middle
School.

The turmoil of middle school turns many teachers away, said the school's principal, Danielle
S.~Battle. Students' bodies and minds are changing, and disparities in learning abilities are
playing out.

``A lot of people will say, 'I'll do anything but middle school,' '' she said.

But this is precisely where Mr.~Dunn chose to be, having seen too many people at the end of lives
gone wrong, and wanting to keep these students from ending up like his former clients. He quotes
Frederick Douglass: ``It is easier to build strong children than to repair broken men.''

The school has institutional architecture that brings prisons to mind, but Ms.~Battle has warmed it
with colorful paint and brighter light. Ninety-three percent of students are black and 5 percent
Hispanic; some 97 percent qualify for free or reduced lunch.

``I just walked in here and fell in love with the place,'' Mr.~Dunn recalled. His day begins at 8
a.m., when he stands by the school's buzzing metal detector, checking bags, as nearly 600 students
file through in a half-hour. It is not a popular job, but he uses the time like a politician working
the plant gate at shift changes. Saying hello with a smile, he taps the bags, peeks inside, sends
the students along. But he is also no-nonsense, with a ``hey-hey-hey!'' to pull back the ones who
try to slip around the detector. Ms.~Battle drops by to greet the students.

As a lawyer, Mr.~Dunn said, he saw his job as ``telling stories,'' to help judges see each client as
a human being who may or may not have done terrible things, but who suffered wrongs at trial or
earlier in life -- and who deserved fairness under the rule of law, perhaps even mercy.

He told clients' stories while defense counsel in the Army Trial Defense Service, in Florida, in New
York State and most recently at the Georgia Resource Center, the nonprofit law firm he led. Though
the center does not keep a scorecard, the strategy has resulted in delayed executions, commuted
death sentences and even overturned convictions, said Brian Kammer, who took over as executive
director.

``If you're just talking about the legal niceties of the case, you're boring the heck out of the
audience,'' Mr.~Kammer said. ``You're squandering the moral force of your argument.''

After decades of accumulating such stories, Mr.~Dunn said, he recognized a common thread: the lack
of a supportive authority figure like a teacher, of a helping hand that might have meant ``the
difference between a good life and a ruined life.''

Illness forced his decision to leave the law. In 2006, he ignored a sore throat and worked through
two months of grueling hearings in four cases back to back. Bacteria entered his bloodstream,
causing toxic shock; the infection caused deterioration in his spine and led to congestive heart
failure. He recovered, but not fully; this year, Mr.~Kammer recalled, Mr.~Dunn met with the staff
and said: ``I have the heart of a 70-year-old man. If I continue to do this work at the level I want
to do it, I'm going to die.''

The same day that he left the center, he showed up at the Atlanta training program of Teach for
America. During his training, he focused on special education, recalling that he saw learning
disabilities ``in nearly every case'' on death row. He now works mainly in classrooms that blend
special education students with the general population.

When he interviewed with the administrators at King, though, he encountered skepticism. ``I was just
baffled by why he'd want to come here,'' said Barbara Shea, an assistant principal, standing with
Mr.~Dunn as the hallways cleared and an afternoon class began. ``I tried to warn him -- I wanted him
to understand it was not an easy job.''

Across the hall a classroom door opened, and the teacher pulled a tall, angry student out by his arm
and asked Mr.~Dunn for help. ``He's having a Taylor moment,'' Ms.~Shea said, referring to a girl in
the class. ``He wants to sock her in the mouth.''

The boy walked toward the exit, but Mr.~Dunn argued him back to a chair where he sat, stormy and
silent. Mr.~Dunn talked softly to him, helping him to settle down. Walked to his office, Mr.~Dunn
chuckled at the thought that Ms.~Shea might have seen him as a dreamer. ``You can't be a starry-eyed
idealist and do defense work in capital cases for 20 years,'' he said.

Propped against the wall in the office -- actually, a converted teacher workroom that his colleagues
pass through to get to a restroom -- is a clipboard with a paraphrase of a quotation from former
Justice Harry A.~Blackmun of the United States Supreme Court: ``From this day forward, I no longer
shall tinker with the machinery of death.'' It was from a dissent in a 1994 case in which the
justice argued that ``the death penalty experiment has failed.''

Mr.~Dunn was in the office in August when he got word that the Supreme Court had ordered a lower
court to reconsider the case of one of his best-known former clients, Troy Davis. At that moment,
being out of the game was ``really hard,'' Mr.~Dunn said.

But in the classroom, there is no hint of regret. In the afternoon, a student, Shamon Nations,
abruptly asked, ``What made you come to school and give up your other job?'' He replied, ``Because I
love you guys.'' Somehow, it does not sound saccharine.

Shamon was not satisfied. ``Yeah, but what about the money?''

``I made a lot more money last year,'' Mr.~Dunn acknowledged. ``But it's not about money.''

Between classes, as he walked down the halls and the class bell rang, he stopped stragglers who
might have been tempted to keep wandering. ``Hey, buddy,'' he asked one. ``Where are you going?'' He
slipped an arm around the boy's shoulder, and used it for leverage to give a gentle shove in the
right direction.

\section{Calling Story of Boy and Balloon a Hoax, Sheriff Seeks Felony Charges}

\lettrine{A}{}\mycalendar{Oct.'09}{19} story that transfixed millions of television viewers last
week -- about a child aboard a spaceshiplike balloon as it floated uncontrollably across the
Colorado landscape -- was concocted by the boy's parents to gain publicity, law enforcement
officials said Sunday.

``They put on a very good show for us, and we bought it,'' Sheriff Jim Alderden of Larimer County
said Sunday morning at a news conference.

Calling the episode an elaborate hoax, Sheriff Alderden said he expected to file felony\cite{felony}
charges against the parents -- Richard Heene and his wife, Mayumi -- who claimed Thursday that their
6-year-old son, Falcon, had climbed aboard the homemade helium balloon before it soared thousands of
feet into the skies above Fort Collins, Colo.

The report prompted a large-scale aerial and ground pursuit, but the authorities learned hours later
that Falcon had never left the ground. The motivation, the sheriff said, was a shot at reality
television stardom.

The allegation of a hoax capped days of furious speculation about the Heene family, and in
particular about Mr.~Heene, who fashioned himself a scientist and television star but had struggled
to gain attention for his show ideas.

Neighbors and former colleagues of Mr.~Heene described a man eager for fame, whose explosive temper
was paired with a wild imagination. On his MySpace page, under the heading ``Who I'd like to meet,''
Mr.~Heene said he wanted to interview ``real aliens from outer space.''

A production company that had been developing a show about the Heenes -- who had previously starred
in two episodes of ``Wife Swap'' on ABC -- said it was no longer working with the family, but did
not say why.

Mr.~and Ms.~Heene had not been arrested as of Sunday evening, but Sheriff Alderden said that among
the charges being considered were three felonies: conspiracy to commit a crime, contributing to the
delinquency of a minor, and attempt to influence a public servant. The last carries a possible
prison term of six years. The charges could also include a misdemeanor, filing a false report.

The sheriff said child protection authorities had been notified about the case, although the three
Heene children, ages 6, 8 and 10, have not been taken from their parents. The local authorities have
also asked federal officials to look into the case.

The sheriff said his conclusions were based on separate interviews of the Heenes and their three
children, as well as on searches of their computers, e-mail records and documents in their home. He
said the decision to release the balloon and tell the authorities that Falcon was aboard was made
two weeks ago.

``These people wanted to get some national attention, with the ultimate hope that they would end up
with some form of a TV deal,'' said the Larimer County undersheriff, Ernie Hudson.

The Heenes did not comment on the sheriff's accusations on Sunday, but when Mr.~Heene was approached
by a reporter for The Associated Press at a Wal-Mart on Sunday morning, he said with tears edging
his eyes, ``This thing has become so convoluted.'' In earlier interviews, he had vigorously denied
that the episode was a hoax.

Wittingly or not, the Heenes produced an astonishing reality show on Thursday. Their report of a
missing child -- first made in the father's phone call to a local TV station, then followed up by
the mother's seemingly panicked phone call to 911 -- captured the country's collective attention as
news helicopters tracked the balloon's ride across northern Colorado.

Family members told the authorities that Falcon's brother Bradford had claimed he saw the 37-pound
boy in the compartment when the balloon accidentally broke off its moorings. Televised by cable news
for more than an hour, the wayward balloon was found to be empty when it came to a landing,
prompting fears that Falcon had fallen out. Hours later, it was discovered that he had been home the
whole time, hiding in an attic\cite{attic} above the garage.

Sheriff Alderden said he and other law enforcement authorities had genuinely believed they were
trying to rescue a small boy in a balloon hundreds of feet above the ground, and for several hours,
they shared the same anguish millions of television viewers did. ``We were very worried that the
life of a small child, a 6-year-old child, may indeed be in jeopardy\cite{jeopardy},'' he said.

But an interview the Heenes gave to CNN on Thursday night fueled speculation that the ordeal had
been a publicity stunt. On ``Larry King Live,'' Falcon said to his father, ``You guys said that, um,
we did this for the show.''

That remark led the authorities to re-interview the parents separately, Sheriff Alderden said. He
said Colorado law forbade him to comment on whether the Heenes had confessed.

According to former colleagues, the parents met while enrolled at an acting school in Los Angeles.
Mr.~Heene had long held dual interests in scientific tinkering and television: in 2001 he produced
an instructional video about turning cardboard boxes into forts for children.

``He was very ambitious and wanted to get some sort of show going,'' said Vincent LeGrow, who edited
actors' reels for a small company Mr.~Heene founded.

Mr.~LeGrow said Mr.~Heene had once spoken with the actor William Shatner about potentially hosting a
show. Mr.~Heene was the co-host of a show called ``The Science Detectives,'' and sometimes said he
was affiliated with the Science Detective Research Group, though few records exist for the company.
Some of the ``Detectives'' segments now appear on YouTube, alongside Mr.~Heene's theories about
U.F.O.'s and civilization on Mars.

He apparently had other media ideas: he told a Tennessee newspaper in 2006 that he had produced a
tornado-chasing documentary and was seeking a TV distributor, and two years later, according to an
entry in an academic journal, he was producing a different documentary about electromagnetic fields
in storms.

Mr.~Heene's projects never attracted measurable attention, but his propensity for chasing tornadoes
and building flying saucers attracted the producers of ``Wife Swap,'' who cast the family for two
episodes, including one last March. The production company behind ``Wife Swap,'' RDF, worked on a
standalone show about the Heenes at one point, but the company said it was ``no longer in active
development with the family.'' It did not respond to further questions.

Mr.~Heene was apparently pursuing a reality TV job as late as last month, when he signed up for an
account on RealityWanted.com, a Web site that connects reality television casting agents and
aspiring contestants, said Mark Yawitz, a co-founder of the site.

Mr.~Yawitz was stunned\cite{stun} that the balloon flight might have been a publicity stunt.

``There are easier ways to get on a reality TV show,'' he said.

\section{Health Care Poses Stiff Tests for Top Democrats}

\lettrine{A}{s}\mycalendar{Oct.'09}{19} she pulls together a health care bill, Speaker Nancy Pelosi
has been exceedingly direct, lawmakers say, asking them explicitly what it will take to win their
vote.

Across the Rotunda, Senator Harry Reid, the majority leader, is trying a different tack, acting like
what one participant in closed-door sessions described as a coach, urging key chairmen not to get
bogged down in pride of authorship and to keep their eyes on the legislative ball.

While they may have different styles and different sets of Democrats to assemble behind separate
proposals, Ms.~Pelosi and Mr.~Reid have an identical goal: passage this year of a major health care
overhaul.

Achieving that end and enacting a top priority of President Obama will present perhaps the stiffest
test yet of the skills of the two Congressional leaders. How they perform could influence not only
the political fate of their own colleagues, but the standing of the administration as well.

No one expects it to be easy.

``It is not painless to do this,'' said Senator Christopher J.~Dodd, the Connecticut Democrat who
oversaw the legislation generated by the Senate health committee.

Mr.~Reid and Ms.~Pelosi are at the moment working independently to merge sweeping health care bills
produced by two committees in the Senate and three in the House into bills each must squeeze through
the respective chambers. Then, Democrats hope, they will negotiate one final version between them
before sending it to Mr.~Obama's desk.

Ms.~Pelosi, who has been meeting incessantly with fellow Democrats on health care, said she was now
focused on producing the strongest possible legislation that can pass the House to provide her side
of the Capitol with as much leverage as possible in future talks with the Senate.

``I'm not into dealing with the politics of the Senate,'' she said last week. ``They do their policy
and their politics. We meet them at the conference table, and that's where we will deal with that.''

Mr.~Reid, similarly, said that he was focused on devising a bill that could win the 60 votes needed
to overcome a likely Republican filibuster\cite{filibuster} and could not yet turn to any task
beyond that.

The challenge for Ms.~Pelosi is to write a measure with sufficient coverage and benefits to appease
the left wing of her caucus without alienating\cite{alienate} too many of the moderate and
conservative Democrats whose votes she needs.

``She's tough,'' said Representative Steny H.~Hoyer of Maryland, the majority leader. ``She's a
great negotiator. She knows at some point in time, she has got to get to 218.''

Mr.~Reid may have the more difficult job since Ms.~Pelosi, of California, has a larger majority as
well as stricter House rules that limit opportunities for Republicans to slow the process.

Mr.~Reid has to come up with a bill that can win at least procedural backing from virtually all 58
Democratic senators and two independents aligned with his party, which will necessitate navigating
divisions between the majority of his members who are pushing for a public option and centrist
Democrats opposed to the idea. All while holding Republicans at bay.

``It requires the wisdom of Solomon and the patience of Job,'' said Senator Evan Bayh, Democrat of
Indiana and one of the centrists whose support Mr.~Reid is struggling to win.

At the same time, Mr.~Reid is facing what could be a difficult re-election contest next year in
Nevada, where his poll numbers are lagging\cite{lag}. He cannot afford to be seen as the architect
of a bill that creates a public uproar or alienates his allies in organized labor. The competing
demands were evident this weekend as Mr.~Reid left the health debate in the Capitol on Thursday to
fly home for a series of events.

As the health fight enters the next stage, his notoriously\cite{notoriously} independent colleagues
are complicating his political life even more. A significant number of them are stepping up their
push for some form of public health insurance option to be included in the Senate version of the
bill -- a potential deal breaker for a few other Senate Democrats.

``There are 52 solid Democrats for the public option,'' said Senator Tom Harkin, the Iowa Democrat
who is chairman of the health committee. ``Only about five Democrats oppose it. Should the 52 give
in to the five? Or should the five go along with the vast majority of the Democratic caucus?''

Mr.~Reid is working with the chairmen of the health and finance committees as well as his leadership
team, in regular consultation with the White House, in trying to meld the two bills.

``We're going to work very hard with the two chairmen and the White House to see what, if anything,
we're going to do,'' Mr.~Reid told reporters.

He could choose to leave more contentious\cite{contentious} elements out of the measure he brings to
the floor and allow backers of the public option, for instance, to try to add it through an
amendment. If they can assemble the necessary votes, it would be a triumph\cite{triumph}; if they
lose, they would have had their chance.

House Democratic leaders are determined that their initial bill will include some type of government
insurance plan, leaving Ms.~Pelosi facing difficult choices in deciding what this public plan should
look like.

Liberals say it should be modeled on Medicare and should use Medicare fees as a basis for paying
doctors and hospitals, as in the bill approved by the House Ways and Means Committee.

But centrist and conservative Democrats say the public plan should negotiate rates with doctors and
hospitals, as private insurers do.

Supporters of a government plan say it would hold down costs by putting competitive pressure on
private insurers. Medicare rates are normally much lower than private insurance rates, so using them
could save money for the government and allow Congress to direct more money to help low- and
middle-income families buy insurance.

In the end, Ms.~Pelosi and Mr.~Reid will not have to go it alone. Senate and House leaders say
Mr.~Obama must weigh in at some point and put his heft behind the eventual legislation, a move the
leaders hope will eliminate any last doubts among both liberal Democrats who fear the emerging plan
does not go far enough and conservative Democrats who worry about costs and other potential
complications.

But at the moment, it is Ms.~Pelosi and Mr.~Reid who are on the spot.

\section{Obama Aides Go on TV to Criticize Wall Street}

\lettrine{S}{triking}\mycalendar{Oct.'09}{19} a populist tone, several of President Obama's leading
advisers on Sunday issued stern warnings to Wall Street. They said big banks must not resist greater
government oversight now that they have regained their financial footing through taxpayer financed
bailouts.

Rahm Emanuel, the White House chief of staff, said Americans ``have a right to be frustrated and
angry'' at reports that a year after the government used \$700 billion to save major lending
institutions from collapse, Wall Street appears poised to hand out another round of hefty bonuses.

Mr.~Emanuel also scolded banks for turning around as soon as they became healthy and fighting
against the regulations the president and his supporters in Congress deem necessary tohead off the
same kind of financial crisis in the future.

``The risks that they took, took the economy to a place, it was near a depression,'' Mr.~Emanuel
said in an appearance on ``State of the Union'' on CNN, referring to the financial institutions.

``They have a responsibility to be part of the solution, not being the obstacle.''

David Axelrod, one of President Obama's closest advisers, called the bonuses and the lobbying
``offensive.''

``The most offensive thing is we haven't seen the kind of increase in lending that we should,'' he
said on ``This Week'' on ABC.

``There are a lot of small businesses, creditworthy businesses around this country, who still can't
get the capital they need to grow, which is important for our economy, and you've seen these same
institutions spend tens of millions of dollars lobbying the Congress to try and stop financial
regulatory reform.''

Wall Street's bonuses are not new fodder for outrage. Resentment, however, has grown in recent
months as relief for some business sectors has not stopped unemployment from rising.

Valerie Jarrett, another close adviser to President Obama, said in an appearance on ``Meet the
Press'' that while the bailout and stimulus packages of the last year had stopped the economy from
sliding into a depression, ``the unemployment rate is still much too high.''

``The president will not be satisfied until every single American who wants to work has a job,'' she
said. When asked what the administration was considering doing to stimulate job creation,
Ms.~Jarrett said all ideas were on the table. But when asked specifically whether the administration
would support a second stimulus package, she said the full effects of the first recovery package had
yet to be felt.

``We've only spent a little less than half of the money,'' she said, referring to the first stimulus
package.

Describing unemployment as a ``lagging indicator,'' she added: ``Let's wait and see. Let's let the
recovery bill do its job, and then we'll see.

Senator Christopher J.~Dodd, Democrat of Connecticut and chairman of the Senate Banking Committee,
issued his own warnings to Wall Street on ``Meet the Press'' on NBC.

``Now, obviously, when you see these bonuses being paid out, it's a source of outrage in the
country, and it should be. What are these people thinking about at these companies?'' he asked.

``These firms on Wall Street need to understand that what they're doing by providing these bonuses,
particularly when they received so much federal money, is an outrage,'' he said.

\section{The Battle Between the White House and Fox News}

\lettrine{T}{he}\mycalendar{Oct.'09}{19} Obama administration, which would seem to have its hands
full with a two-front war in Iraq and Afghanistan, opened up a third front last week, this time with
Fox News.

Until this point, the conflict had been mostly a one-sided affair, with Fox News hosts promoting tax
day ``tea parties'' that focused protest on the new president, and more recently bringing down the
presidential adviser Van Jones through rugged coverage that caught the administration, and other
news organizations, off guard. During the health care debate, Fox News has put a megaphone to
opponents, some of whom have advanced far-fetched theories about the impact of reform. And even
farther out on the edge, the network's most visible star of the moment, Glenn Beck, has said the
president has ``a deep-seated hatred for white people.''

Administration officials seemed to have decided that they had had enough.

``We're going to treat them the way we would treat an opponent,'' Anita Dunn, the White House
communications director, said in an interview with The New York Times. ``As they are undertaking a
war against Barack Obama and the White House, we don't need to pretend that this is the way that
legitimate news organizations behave.''

Ah, but pretending has traditionally been a valuable part of the presidential playbook. Smiling and
wearing beige even under the most withering news media assault is not only good manners, but also
has generally been good politics. While there is undoubtedly a visceral thrill in finally setting
out after your antagonists, the history of administrations that have successfully taken on the media
and won is shorter than this sentence.

Not that they haven't tried. In his second Inaugural Address, Ulysses S.~Grant said he had ``been
the subject of abuse and slander scarcely ever equaled in political history.'' President William
McKinley labeled a gathering of the press a ``congress of inventors,'' and President Franklin
D.~Roosevelt assigned less favored press members to his ``Dunce Club.'' Sometimes the strategy
worked -- or caused no lasting damage. McKinley, like Grant, was elected to a second term. Roosevelt
also won a third and fourth.

As Americans turned to TV for news, enmity from presidents soon followed. Vice President Spiro
T.~Agnew said ``self-appointed analysts'' at the Big Three networks exhibited undisguised
``hostility'' toward President Richard M.~Nixon, subjecting his speeches to ``instant analysis and
querulous criticism.'' Later, in the dispute with The Times over the Pentagon Papers, Mr.~Nixon's
national security adviser, Henry Kissinger, accused the newspaper of treason.

Neither of the Bush presidents had a particularly cozy relationship with the press. George H.W. Bush
finished the campaign in 1992 with a bumper sticker that suggested, ``Annoy the Media. Vote Bush.''
And George W.~Bush, in the words of ABC's Mark Halperin, viewed ``the media as a special interest
rather than as guardians of the public interest.'' Bill Clinton, too, distrusted the press, as did
others in his administration. When Vincent Foster, Mr.~Clinton's deputy White House counsel,
committed suicide in 1993, he left behind a note accusing the Wall Street Journal's editorial page
of lying.

Even though almost all the critiques contained a kernel of truth, in each instance the folks who had
the barrels of ink, and now pixels, seemed to come out ahead. So far, the only winner in this latest
dispute seems to be Fox News. Ratings are up 20 percent this year, and the network basked for a week
in the antagonism of a sitting president

It could all be written off as a sideshow, but it may present a genuine problem for Mr.~Obama, who
took great pains during the campaign to depict himself as being above the fray of over-heated
partisan squabbling. In his victory speech he promised, ``I will listen to you, especially when we
disagree.''

Or not. Under the direction of Ms.~Dunn, the administration has begun to punch back. On Sept.~20,
the president visited all the Sunday talk shows save Fox News', with Ms.~Dunn explaining that Fox
was not a legitimate news organization, but a ``wing of the Republican Party.''

The one weapon all administrations can wield is access, and the White House, making it clear that it
will use that leverage going forward, informed Fox News not to expect to bump knees with the
president until 2010. But Fox News, as many have pointed out, is not in the access business. They
are in the agitation business. And the administration, by deploying official resources against a
troublesome media organization, seems to have brought a knife to a gunfight.

Tactics aside, something more fundamental is at risk. Even the president's most avid critics admit
he exudes a certain cool confidence. The public impression of him is that if anyone were to, say,
talk trash on the basketball court with Mr.~Obama, he would not find much space for rent in
Mr.~Obama's head.

Mr.~Obama has also shown a consistent ability to disarm or at least engage his critics. When he
eventually sat for an interview with the Fox News personality Bill O'Reilly two months before the
election, it made for great television. But for the time being, d\'etente seems very far away and
the gap seems to be widening.

On the official White House Web site, a blog called Reality Check provides a running tally of
transgressions by Fox News. It ends with this: ``For even more Fox lies, check out the latest
'Truth-O-Meter' feature from Politifact that debunks a false claim about a White House staffer that
continues to be repeated by Glenn Beck and others on the network.''

People who work in political communications have pointed out that it is a principle of power
dynamics to ``punch up `` -- that is, to take on bigger foes, not smaller ones. A blog on the White
House Web site that uses a ``truth-o-meter'' against a particular cable news network would not seem
to qualify. As it is, Reality Check sounds a bit like the blog of some unemployed guy living in his
parents' basement, not an official communiqu\'e from Pennsylvania Avenue.

The American presidency was conceived as a corrective to the royals, but trading punches with cable
shouters seems a bit too common. Perhaps it's time to restore a little imperiousness to the
relationship.

\section{Foreclosures Force Ex-Homeowners to Turn to Shelters}

\lettrine{T}{he}\mycalendar{Oct.'09}{19} first night after she surrendered her house to foreclosure,
Sheri West endured the darkness in her Hyundai sedan. She parked in her old driveway, with her
flower-print dresses and hats piled in boxes on the back seat, and three cherished houseplants on
the floor. She used her backyard as a restroom.

The second night, she stayed with a friend, and so it continued for more than a year: Ms.~West --
mother of three grown children, grandmother to six and great-grandmother to one -- passed months on
the couches of friends and relatives, and in the front seat of her car.

But this fall, she exhausted all options. She had once owned and overseen a group home for homeless
people. Now, she succumbed\cite{succumb} to that status herself, checking in to a shelter.

``No one could have told me that in a million years: I'd wake up in a homeless shelter,'' she said.
``I had a house for homeless people. Now, I'm homeless.''

Growing numbers of Americans who have lost houses to foreclosure are landing in homeless shelters,
according to social service groups and a recent report by a coalition of housing advocates.

Only three years ago, foreclosure was rarely a factor in how people became homeless. But among the
homeless people that social service agencies have helped over the last year, an average of 10
percent lost homes to foreclosure, according to ``Foreclosure to Homelessness 2009,'' a survey
produced by the National Coalition for the Homeless and six other advocacy groups.

In the Midwest, foreclosure played a role for 15 percent of newly homeless people, according to the
survey, reflecting soaring rates of unemployment -- Ohio's reached 10.8 percent in August -- and
aggressive lending to people with damaged credit.

At a shelter for women and children run by the West Side Catholic Center in Cleveland, where
Ms.~West now lives, foreclosure accounted for zero arrivals in 2007, the center's executive
director, Gerald Skoch, said. Last year, two cases emerged. This year, the number has already
reached four.

Similar increases have been reported at shelters in California, Michigan and Florida, where a
combination of joblessness and the real estate bust have generated unusually severe rates of
foreclosure.

Most people who become homeless because of foreclosure had been low-income renters whose landlords
stopped making their mortgage payments, leaving them scrambling for new housing with little notice
and scant savings, according to the survey and interviews with shelters.

But in recent months, there has been a visible increase in the number of former homeowners showing
up in shelters. Like Ms.~West, most have landed there after months trying to stave off that fate.

``These families never needed help before,'' said Larry Haynes, executive director of Mercy House in
Santa Ana, Calif. ``They haven't a clue about where to go, and they have all sorts of humiliation
issues. They don't even know what to say, what to ask for.''

Many start off camping out in cars, particularly in warmer places.

``We've seen a rise in people sleeping in their cars,'' said Rick Cole, city manager in Ventura,
Calif., which recently allowed car-camping in designated areas. ``Some are foreclosed former
homeowners, and some couldn't afford their rent. People will give up their house before they give up
their car.''

Those with means try to rent homes or apartments, though tainted credit often makes that impossible.
Growing numbers are landing in motels that rent by the week, cramming whole families into single
rooms and using hot plates as kitchens. But as unemployment expands, many are losing the wherewithal
to remain.

Many take refuge with families and friends, occupying extra bedrooms, basements and attics. But such
hospitality rarely lasts.

So, as lean times endure and paychecks disappear, homeless shelters are absorbing those who have run
out of alternatives.

For Ms.~West, whose youthful appearance belies her age, in her mid-50s, the nights spent on couches
in other people's homes were uncomfortably familiar. She grew up an only child in a housing project
in Neptune, N.J., where her mother slept in the lone bedroom, and she occupied a pullout sofa in the
living room.

``I've always had this dream of doing better,'' she said. ``I always wanted to own my own house.''

She realized that dream shortly after arriving in Cleveland with her husband and two children in the
early 1990s. At first, they rented. But one fall afternoon, Ms.~West found herself on a block lined
with leafy trees in Mount Pleasant, a neighborhood east of the Cuyahoga River that was a magnet for
middle-class black families like hers. Red brick homes with wooden porches sat on ample lots. Public
schools were a few blocks away.

When she saw an ad in the Sunday paper offering a house on that very block, she bought it for
\$45,000; for the \$9,000 down payment she used the savings her mother had left her when she died.
She and her husband assumed the mortgage from the previous owner, with affordable payments of less
than \$400 a month.

Ms.~West then had a job as a maintenance worker at an apartment complex for about \$9 an hour. Her
husband earned about \$10 an hour as a truck driver. As the years passed, they added shrubbery to
the front yard and photos of children's birthday parties to the walls.

``I thought that was going to be my house,'' she said.

She tapped her inheritance to buy another house on nearby Union Street, paying \$15,000 in cash for
a light-blue, vinyl-sided A-frame. She turned the house into a home for five homeless people. She
did their laundry, reminded them to take their medications and cooked meals, while collecting
payments of up to \$750 a person each month from the agencies that placed them.

Over the years, Ms.~West and her husband spent more than they earned. They used credit cards to
finance restaurant meals. They bought a new S.U.V.

At the group home, Ms.~West's compensation slipped as the state limited benefit payments. Yet every
month brought the same thicket of bills -- water, electricity, gas, plus food for the people under
her charge.

In 2001, Ms.~West and her husband took out a \$67,000 mortgage on the Union Street house -- which
had increased considerably in value -- to refinance high-interest debts, assuming payments of nearly
\$700 a month.

Two years later, her husband left her.

``It just took the life out me,'' she said. ``I was in a very bad state, a very depressed situation.
Things just kind of went downhill. I just didn't care anymore.''

By 2005, she was broke. She sold the brick house to her cousin, disbanded the group home and moved
in. She paid what bills she could through temporary jobs as a signature collector for petition
drives. But as many months passed without work, the bills piled up past due.

By the next year, terse letters were coming from the mortgage company -- notices of delinquency,
then threats of foreclosure. Much of the neighborhood was in a similar state. Broken windows sat
unrepaired at a two-story apartment block across the street, where tattered curtains flapped in the
breeze. The city boarded up abandoned homes to deter vagrants, drug addicts and prostitutes.

Ms.~West wrote to her mortgage company, seeking lower payments. But with tainted credit and no
full-time job, she was not a candidate for a deal. Fliers beckoned with relief as companies offered
to negotiate with her lender for lower payments. But when she called, the companies demanded upfront
payments as high as \$500.

``I told them, 'if I had that money, I wouldn't be going into foreclosure,' '' she said.

In the spring of 2008, Ms.~West accepted an offer from the mortgage company: move out, hand over the
keys and collect \$2,500. She sold what furniture she could and put the rest on the street --
tables, beds, a couch.

Her uncle had said she could stay with him for a while. But when she called him to say she was on
the way, he told her that his girlfriend was uncomfortable with the arrangement. Ms.~West's daughter
was in a cramped rented house with her boyfriend and her two children. Her son was in a rooming
house.

So Ms.~West, a stylish woman with a penchant for shiny lipstick and glittering jewelry, wound up
camping in her car. She listened to the radio to drown out the voices of prostitutes trawling the
street. She meditated. (``Just blank out everything in your mind,'' she said. ``Just go to a place
that's peaceful, like a beach.'') She prayed.

``It was scary,'' she said. ``Here I am, alone, and I don't have nowhere to go.''

The next day, she moved in with a friend, remaining there for about three months. For several more
months, she stayed with the cousin who had bought her old brick house and was living there with her
husband and seven children. Toys lay scattered across the floor. The walls vibrated with music,
television and the sounds of children. She lay awake on the couch, a vagabond in the one place that
had once felt so solid.

``I was losing my mind,'' she said.

She was grateful to be inside -- particularly during the Cleveland winter -- yet never comfortable
or stable enough to plan beyond the next day.

``You know in the back of your mind that people don't really want you there,'' she said.

Sometimes, she lived out of her car, spending days at the public library, where she washed up in the
restroom and used a computer to scan meager\cite{meager} job listings.

Finally, a woman she met on the street took her in and helped her formulate a recovery plan. She
signed up for food stamps. She enrolled at a community college in a three-month, state-financed
training program that would give her a certificate for an entry-level job in biotechnology, putting
her in position to earn as much as \$16 an hour.

In September, she got a bed at the homeless shelter, reluctantly accepting that she needed her own
space to re-establish her life.

``I never wanted to go to the shelter because of the stigma\cite{stigma},'' she said. ``I'm a very
independent person. I felt like I got myself into this situation, and I've got to get myself out.
But I knew I couldn't just keep going back and forth and staying with these people and not moving
forward with my life.''

She sleeps in a twin bed with a flower-print duvet, in a small room painted lavender. Her plants
line the windowsill. She keeps to herself, reading motivational books, as she prepares to start
classes next month.

She is working again, taking care of senior citizens in their homes part time, and saving money.

By December, she will exhaust the shelter's 90-day limit, so she is hurrying to line up a house to
rent while arranging a subsidy through the West Side Catholic Center.

She is still shaken by the past and anxious about the future, but she is again looking ahead.

``I do want to eventually own a house again,'' she said. ``That's the American dream. That's what
everybody wants.''

\section{Face-to-Face Socializing Starts With a Mobile Post}

\lettrine{T}{witter}\mycalendar{Oct.'09}{19} and Facebook ask users to answer the question: What are
you doing right now?

But for many urbanites in their 20s and 30s, two other questions are just as important: Where are
you, and can I come join you?

For them, a fast-growing social networking service called Foursquare is becoming the tool of choice.
A combination of friend-finder, city guide and competitive bar game, Foursquare lets users ``check
in'' with a cellphone at a bar, restaurant or art gallery. That alerts their friends to their
current location so they can drop by and say hello.

``It's planned serendipity,'' said Emily Woolf, 24, a strategic planner living in Brooklyn who
checks in on Foursquare when she wants to grab coffee or a drink with friends. ``At this point, I
don't even bother texting or calling my friends. I just check Foursquare to see if they're nearby
and go meet them.''

Just seven months old with about 60,000 users so far, Foursquare is still getting off the ground --
especially when compared with supersize services like Facebook and Twitter, which have millions of
members. But that underground status is part of Foursquare's appeal, its fans say. It is not yet
cluttered with celebrities, nosy mothers-in-law or annoying co-workers.

``On Twitter, there are more than 3,000 people that follow me, and Facebook is more of a business
community now,'' said Annie Heckenberger, 36, who works at an advertising agency in Philadelphia.
``Foursquare is more of the people that I actually hang out with and want to socialize with.''

It is akin to knowing about a hip new club before everyone else, said Deborah Schultz, an analyst
with the Altimeter Group who specializes in trends in social media.

``There will always be people who love new technology and want to test it out, kick the tires,'' she
said. ``Once those services become too big and the bridge-and-tunnel crowd shows up, they can lose
some of that initial interest.''

One factor that might help Foursquare retain its intimate feel is that most of its members are
picky\cite{picky} about who can see the real-time footprints that they are leaving across the cities
in which they live.

Foursquare emerged from the ashes of an earlier mobile service called Dodgeball, which was
introduced in 2004 by Dennis Crowley and a classmate from New York University and was sold to Google
in 2005.

After Google killed the service, which relied heavily on text messages, Mr.~Crowley revived the
concept. He and his business partner, Naveen Selvadurai, introduced Foursquare in March during South
by Southwest, an annual technology and music conference that draws media types, Internet
entrepreneurs and technology fans.

Other companies, like BrightKite, Loopt and Google Latitude, are also offering services aimed at
helping friends find each other on the go. But Foursquare has attracted more attention than the
others, in part because it incorporates elements of gaming and social competition.

The system awards points and virtual badges to players depending on how often they go out and which
places they visit. Users who frequent a particular place enough times are crowned ``mayor'' of that
particular location.

``People are very territorial about their mayorships,'' Mr.~Crowley said. ``It's almost like
bragging rights.''

Ms.~Heckenberger says she once even leapt out of bed to reclaim the title at her local watering
hole, the Swift Half Pub, where another player had briefly wrested away the honor. ``Foursquare has
become like an obsession for me,'' she said.

So far, Foursquare has no revenue, and the company is still developing its business model.
Mr.~Crowley and Mr.~Selvadurai say they are focusing on trying to build up the infrastructure,
expand the user base and develop a database of locations. They see the service as developing beyond
a nightlife game into a ``service that encourages people to do new things and get rewarded,''
Mr.~Crowley said.

Foursquare is available in 31 cities in the United States, including New York, Chicago, San
Francisco, Boston and Washington. The service is also operating in London, Amsterdam and three
Canadian cities.

Mr.~Crowley said the fierce competition of Foursquare players to claim the mayorship of their
favorite haunts suggested that corporate sponsorships could be a lucrative source of revenue. Some
small businesses, like Sugar Mama's Bake Shop in Austin, Tex., are starting to advertise special
deals on the service, although Mr.~Crowley has yet to begin charging businesses to put deals in the
database.

``For a small business with a limited advertising budget, it's a great way to promote ourselves,''
said Olivia O'Neal, owner of Sugar Mama's. The shop offers Foursquare mayors a free cup of coffee
each time they come in, and regular patrons receive their 10th cupcake free. ``There are about 67
people currently working on those offers, and for a small family-owned business like ours, that's a
really big number,'' Ms.~O'Neal said.

Eventually, Mr.~Crowley said, he would like to work with businesses on sponsored badges. For
example, Starbucks might agree to give a player who visits 10 different Starbucks locations in a
week a special badge and a free coffee drink.

Foursquare's potential has intrigued some notable entrepreneurs and investors. In September, the
company landed a \$1.35 million round of venture financing, led by Union Square Ventures, which
previously backed Twitter and Meetup.

Also contributing to the financing round were Kevin Rose, founder of the social news site Digg; Jack
Dorsey, co-founder of Twitter; and Ron Conway, an angel investor who had backed Google and PayPal.

``Being the hot new thing is a factor, but that's not totally the driving force behind Foursquare,''
said Fred Wilson, general partner at Union Square Ventures. ``There's a part of it that's just fun,
and you can't ignore how important that is in terms of usage.''

\section{In Hawaii's Health System, Lessons for Lawmakers}

\lettrine{I}{mee}\mycalendar{Oct.'09}{19} Gallardo, 24, has been scooping ice cream at a
H\"aagen-Dazs shop at Waikiki Beach for five years, and during that time the shop has done something
its counterparts on the mainland rarely do: it has paid for her health care.

Ms.~Gallardo cannot imagine any other system.

``I wouldn't get coverage on the mainland?'' Ms.~Gallardo asked. ``Even if I worked? Why?''

Since 1974, Hawaii has required all employers to provide relatively generous health care benefits to
any employee who works 20 hours a week or more. If health care legislation passes in Congress, the
rest of the country may barely catch up.

Lawmakers working on a national health care fix have much to learn from the past 35 years in Hawaii,
President Obama's native state.

Among the most important lessons is that even small steps to change the system can have lasting
effects on health. Another is that, once benefits are entrenched\cite{entrench}, taking them away
becomes almost impossible. There have not been any serious efforts in Hawaii to repeal the law,
although cheating by employers may be on the rise.

But perhaps the most intriguing lesson from Hawaii has to do with costs. This is a state where
regular milk sells for \$8 a gallon, gasoline costs \$3.60 a gallon and the median price of a home
in 2008 was \$624,000 -- the second-highest in the nation. Despite this, Hawaii's health insurance
premiums are nearly tied with North Dakota for the lowest in the country, and Medicare costs per
beneficiary are the nation's lowest.

Hawaii residents live longer than people in the rest of the country, recent surveys have shown, and
the state's health care system may be one reason. In one example, Hawaii has the nation's highest
incidence of breast cancer but the lowest death rate from the disease.

Why is Hawaiian care so efficient? No one really knows.

In dozens of interviews, doctors and hospital and insurance executives in Hawaii offered many
theories, including an active population that is culturally disinclined to hospitals, a significant
military presence and a health care market dominated by a few not-for-profit organizations.

But there was another answer: With nearly 90 percent of the populace given relatively generous
benefits, patients stay healthy and health providers have the money and motivation to innovate.

If true, it's a crucial lesson. Health care overhaul efforts at the state and national levels have
so far been largely confined to providing bare-bones insurance coverage to those in need. But
changing the way care is provided has been given short shrift, and medical experts warn that costs
could soar if overhaul legislation passes. After expanding coverage in 2006, Massachusetts is only
now tackling the cost problem as expenses continue their inexorable\cite{inexorable} rise.

But the Hawaii experience suggests that overhauling health insurance before changing the way care is
provided could work, eventually. With more people given access to care, hospital and insurance
executives in Hawaii say they have been able to innovate efficiencies. For instance, the state's top
three medical providers are adopting electronic medical records -- years ahead of most mainland
counterparts.

The Hawaii Medical Service Association, the state's largest insurer and a Blue Cross Blue Shield
member, recently offered the nation's only statewide system whereby anyone for a nominal fee can
talk by phone or e-mail, day or night, to doctors of their choosing.

Kaiser Permanente Hawaii, which covers about 20 percent of the state's population, screens 85
percent of its female members ages 42 to 69 for breast cancer, among the highest screening rates in
the country.

One result of Hawaii's employer mandate and the relatively high number of people with health
insurance is that hospital emergency rooms in the state are islands of relative calm. In 2007, the
state had 264 outpatient visits to emergency rooms per 1,000 people -- 34 percent lower than the
national average of 401.

Dr.~Ray Sebastian splits his time between the emergency room at Kapi'olani Medical Center at Pali
Momi and a hospital in Los Angeles. Nearly all of his poorest patients in Hawaii have routine access
to family doctors who can provide follow-up care, while fewer than half of those in Los Angeles do,
he estimated. So, he said, the emergency room in Hawaii is not clogged with patients suffering minor
problems like medication adjustments and cold symptoms, and patient waiting times are a small
fraction of those in Los Angeles:

``It's like greased lightning here,'' he said.

Other states tried employer-mandated care only to repeal the efforts after employers threatened to
move across state lines. Hawaii's isolation forestalled such threats, and its
paternalistic\cite{paternalistic} plantation history made employer-provided care an easy fit.

In interviews, leaders and employers in Hawaii referred with surprising earnestness to an ``aloha
spirit'' and a sense of familial obligation known as ohana to justify providing care to nearly
everyone.

Hospital executives said they never overbuilt their facilities because the Pacific Ocean meant they
could not delude\cite{delude} themselves into thinking, as their mainland counterparts sometimes do,
that they would be able to attract patients from afar.

Since supply tends to drive demand in health care, this may be one reason Hawaii residents use fewer
health care services -- they get a third fewer magnetic resonance imaging tests and are admitted to
the hospital 26 percent less than the United States average.

There are clear problems with Hawaii's system. Hospitals on the outer islands are small and losing
money. With unemployment rising, so, too, are the ranks of the uninsured -- which is now 10.7
percent of nonelderly adults. Only Massachusetts has a lower share of uninsured adults, and the
national share is 20.4 percent. And there is growing evidence that as the economy has slowed and
premiums have risen, employers have hired more part-time workers who are ineligible for benefits.

Barbara Zacchini, owner of Pizzeria Zacchini on the island of Hawaii, said she makes sure that her
17 part-timers work less than 20 hours a week so she does not have to pay for their care.

``I'm for universal health care,'' Ms.~Zacchini said, ``but it's tough to run a business in this
state and in this economy.''

Some employers are ducking the law altogether. A 61-year-old travel agent in Honolulu said her boss
refused to provide health insurance although he is required by law to do so. She cannot find another
job, so she asked that her name not be used.

She has not been to a primary-care doctor or a gynecologist in years and goes to the emergency room
when she needs care. ``I could have an alien called cancer growing inside me, but who knows?'' she
said. ``It worries me.''

Hawaii law requires employers to offer standardized health plans with low co-pays, no deductibles
and few out-of-pocket costs. Cliff Cisco, a senior vice president at the Hawaii Medical Service
Association, said that having a standardized and popular benefit has helped keep administrative
costs to just 7 percent of revenue, among the lowest in the nation.

Indeed, many in Hawaii are worried that legislation moving through Congress could, if it supersedes
Hawaii law, allow employers to reduce the quality of care provided. So legislators in Hawaii have
pushed for provisions exempting the state.

Chad Buck, owner of Hawaii Food Service Alliance, a grocery distribution company with 140 employees,
said he feared national health care legislation might allow his competitors to provide low-cost,
high-deductible plans in place of the generous benefits now required by the state.

``I don't want to compete against low-quality health care, and I don't want my employees to have a
cheap second,'' Mr.~Buck said.

Richard Caldarazzo, 25, a manager at Lulu's Waikiki Surf Club, said he had worked at restaurants in
Ann Arbor, Mich., and Chicago and had never gotten health insurance.

After he moved to Hawaii and got a job at Lulu's, ``I was really surprised when they told me I'd get
insurance,'' Mr.~Caldarazzo said. ``My parents couldn't believe it.''

\section{Doctor Wants Canada's Athletes in Front, for Flu Shots}

\lettrine{C}{anadians}\mycalendar{Oct.'09}{19} will not be vaccinated against the H1N1 flu until
November, and the federal government has released a priority list for who should receive the shot
first. Another suggestion for the list -- to move Team Canada near the front of the line in
preparation for the Vancouver Olympics in February -- has generated swift and strong reaction from
the public.

The proposal by Robert G.~McCormack, the team's chief medical officer, became public on Tuesday,
when his comments were published by the CBC. The federal government has so far resisted the idea,
but McCormack remains hopeful.

``It has the potential to be emotionally charged because it's presented as elite athlete getting
special treatment,'' he said from Vancouver, where he is an associate professor of orthopedic
surgery at the University of British Columbia. ``But our position is that they need to be
protected.''

Because waiting times for treatments are a perpetual\cite{perpetual} political issue for Canada's
public health-care system, any suggestion that a particular group might receive priority can be
explosive.

Although several readers posted comments on the CBC's Web site backing the idea, they were a clear
minority among the several hundred people who offered comments.

``What makes them more worthy than the rest of us?'' one reader wrote. Another added, ``I suppose
that next they will want all N.H.L. hockey players to get priority!''

The issue is about timing rather than availability. The Canadian government has ordered 50.4 million
doses of vaccine for a country with a population of about 33.5 million. McCormack estimates that
about 200 athletes from Canada will compete.

The Canadian government's priority list is similar to those of many other countries and includes
pregnant women, people under 65 with chronic health conditions, children ages 6 months to 5 years,
as well as essential health-care workers.

McCormack said he was still talking with Health Canada, the federal department that is buying the
vaccines. But the Public Health Agency of Canada, the branch of Health Canada that set the priority
list, appears unwilling to change its plan to accommodate the athletes.

Nadia Mostafa, a spokeswoman for the agency, said in an e-mail message, ``Other groups, such as
Olympics workers, volunteers and athletes, were considered for prioritization but not selected,
given expected vaccine availability prior to the Olympics and the absence of epidemiological
evidence to warrant early prioritization of this group.''

Dr.~Patricia Daly is the chief medical health officer for the Vancouver Coastal Health authority,
and her duties include preparing vaccination recommendations for athletes attending the Vancouver
Games. Her suggestions, which are being relayed to national teams by the Vancouver organizing
committee, do not share McCormack's sense of urgency. She wants athletes to receive injections for
swine flu and seasonal flu at least 14 days before the Games begin.

``We said about two weeks, and I think that's fine,'' she said from Vancouver. ``Two weeks is the
time needed for full protection.''

McCormack has a two-part argument for giving athletes priority for flu vaccines. Although Olympic
athletes are obviously fit, he said, they are also a high-risk group for the swine flu. Training,
travel and team accommodations put them in constant contact with others, increasing the risk of
infection. He also pointed to research showing that intense exercise weakens the immune system.

Second, he said, many athletes will travel overseas for training and competitions just as Canada's
vaccination program begins and may not return until afterward. It is important to protect them
early, he said, to prevent the swine flu from disrupting their training and competition schedules.

He said he was not suggesting that athletes move ahead of pregnant women and high-risk groups. But
his request suggested that, at the very least, they would be in line at the same time.

``If you lose a month of training, it's very serious,'' McCormack said. ``These are people who train
10, 12 years for one day of competition. We live vicariously through our athletes.''

\section{In Chinatown, Sound of the Future Is Mandarin}

\lettrine{H}{e}\mycalendar{Oct.'09}{22} grew up playing in the narrow, crowded streets of
Manhattan's Chinatown. He has lived and worked there for all his 61 years. But as Wee Wong walks the
neighborhood these days, he cannot understand half the Chinese conversations he hears.

Cantonese, a dialect from southern China that has dominated the Chinatowns of North America for
decades, is being rapidly swept aside by Mandarin, the national language of China and the lingua
franca of most of the latest Chinese immigrants.

The change can be heard in the neighborhood's lively restaurants and solemn church services, in
parks, street markets and language schools. It has been accelerated by Chinese-American parents,
including many who speak Cantonese at home, as they press their children to learn Mandarin for the
advantages it could bring as China's influence grows in the world.

But the eclipse of Cantonese -- in New York, China and around the world -- has become a challenge
for older people who speak only that dialect and face increasing isolation unless they learn
Mandarin or English. Though Cantonese and Mandarin share nearly all the same written characters, the
pronunciations are vastly different; when spoken, Mandarin may be incomprehensible to a Cantonese
speaker, and vice versa.

Mr.~Wong, a retired sign maker who speaks English, can still get by with his Cantonese, which
remains the preferred language in his circle of friends and in Chinatown's historic core. A bit
defiantly, he said that if he enters a shop and finds the staff does not speak his dialect, ``I go
to another store.''

Like many others, however, he is resigned to the likelihood that Cantonese -- and the people who
speak it -- will soon become just another facet of a polyglot neighborhood. ``In 10 years,''
Mr.~Wong said, ``it will be totally different.''

With Mandarin's ascent has come a realignment of power in Chinese-American communities, where the
recent immigrants are gaining economic and political clout, said Peter Kwong, a professor of
Asian-American studies at Hunter College.

``The fact of the matter is that you have a whole generation switch, with very few people speaking
only Cantonese,'' he said. The Cantonese-speaking populace, he added, ``is not the player anymore.''

The switch mirrors a sea change under way in China, where Mandarin, as the official language, is
becoming the default tongue everywhere.

In North America, its rise also reflects a major shift in immigration. For much of the last century,
most Chinese living in the United States and Canada traced their ancestry to a region in the Pearl
River Delta that included the district of Taishan. They spoke the Taishanese dialect, which is
derived from and somewhat similar to Cantonese.

Immigration reform in 1965 opened the door to a huge influx of Cantonese speakers from Hong Kong,
and Cantonese became the dominant tongue. But since the 1990s, the vast majority of new Chinese
immigrants have come from mainland China, especially Fujian Province, and tend to speak Mandarin
along with their regional dialects.

In New York, many Mandarin speakers have flocked to Sunset Park, Brooklyn, and Flushing, Queens,
which now rivals Chinatown as a center of Chinese-American business and political might, as well as
culture and cuisine\cite{cuisine}. In Chinatown, most of the newer immigrants have settled outside
the historic core west of the Bowery, clustering instead around East Broadway.

``I can't even order food on East Broadway,'' said Jan Lee, 44, a furniture designer who has lived
all his life in Chinatown and speaks Cantonese. ``They don't speak English; I don't speak Mandarin.
I'm just as lost as everyone else.''

Now Mandarin is pushing into Chinatown's heart.

For most of the 100 years that the New York Chinese School, on Mott Street, has offered language
classes, nearly all have taught Cantonese. Last year, the numbers of Cantonese and Mandarin classes
were roughly equal. And this year, Mandarin classes outnumber Cantonese three to one, even though
most students are from homes where Cantonese is spoken, said the principal, Kin S.~Wong.

Some Cantonese-speaking parents are deciding it is more important to point their children toward the
future than the past -- their family's native dialect -- even if that leaves them unable to
communicate well with relatives in China.

``I figure if they have to acquire a language, I wanted them to have Mandarin because it makes it
easier when they go into the workplace,'' said Jennifer Ng, whose 5-year-old daughter studies
Mandarin at the language school of the Church of the Transfiguration, a Roman Catholic parish on
Mott Street where nearly half the classes are devoted to Mandarin. Her 8-year-old son takes
Cantonese, but only because there is no English-speaking Mandarin teacher for his age group.

``Can I tell you the truth?'' she said. ``They hate it! But it's important for the future.'' Until
recently, Sunday Masses at Transfiguration were said in Cantonese. The church now offers two in
Mandarin and only one in Cantonese. And as the arrivals from mainland China become old-timers, ``we
are beginning to have Mandarin funerals,'' said the Rev. Raymond Nobiletti, the Cantonese-speaking
pastor.

At the Chinese Consolidated Benevolent Association, which has been the unofficial government of
Chinatown for generations and conducts its business in Cantonese, the president, Justin Yu, said he
is the first whose mother tongue is Mandarin to lead the 126-year-old organization. Though he has
been taking Cantonese lessons in order to keep up at association meetings, his pronunciation is
sometimes a source of hilarity\cite{hilarity} for his colleagues, he said.

``No matter what,'' he added, laughing, ``you have to admire my courage.''

But even his association is being surpassed in influence by Fujianese organizations, said Professor
Kwong of Hunter College.

Longtime residents seem less threatened than wistful. Though he is known around Chinatown for what
he calls his ``legendarily bad'' Cantonese, Paul Lee, 59, said it pained him that the dialect was
disappearing from the place where his family has lived for more than a century.

``It may be a dying language,'' he acknowledged. ``I just hate to say that.''

But he pointed out that the changes were a natural part of an evolving immigrant neighborhood: Just
as Cantonese sidelined Taishanese, so, too, is Mandarin replacing Cantonese.

Mr.~Wong, the principal of the New York Chinese School, said he had tried to adjust to the subtle
shifts during his 40 years in Chinatown. When he arrived in 1969, he walked into a coffee shop and
placed his order in Cantonese. Other patrons looked at him oddly.

``They said, 'Where you from?' `` he recalled. `` 'Why you speak Cantonese?' '' They were from
Taishan, he said, so he switched to Taishanese and everyone was happy.

``And now I speak Mandarin better than Cantonese,'' he added with a chuckle. ``So, Chinatown -- it's
always changing.''

\section{China Warns Pirates Who Commandeered Ship}

\lettrine{C}{hina}\mycalendar{Oct.'09}{22} has vowed to rescue the crew of a coal-laden cargo ship
seized by pirates far off the coast of Somalia, but a negotiated settlement is more likely, maritime
experts said Wednesday.

European naval officials said that the 25 Chinese crew members of the seized ship, the De Xin Hai,
were uninjured but that the captors were moving the vessel toward the Somali coast, where the
hostages were expected to join 120 other hostages from other countries who are also awaiting their
freedom.

The seizure of the ship, which was attacked on Monday 700 nautical miles off the coast of Somalia,
suggests that pirates are traveling farther offshore in their effort to evade the international
flotilla that has been sent to the Gulf of Aden to maintain security in one of the world's busiest
shipping lanes. The ship had been en route to India from South Africa, filled with about 76,000 tons
of coal.

It is not the first time a Chinese vessel has been attacked by pirates off Somalia; a fishing
vessel, the Tianyu No.~8, was hijacked by pirates last November and a second vessel, the Zhenhua 4,
repulsed a pirate attack last December. But maritime experts said the seizure of the De Xin Hai
represented the first time that a vessel had been raided so far out at sea, well beyond the
corridor, which is 300 nautical miles off the coast, in which ships can expect some protection.

``We've effectively clamped down and prevented many acts of piracy, which has pushed them further
out into the ocean,'' said Cmdr. John Harbor, a spokesman for the European Union Naval Force, whose
main task is to protect ships carrying food aid to Somalia.

Despite a brief lull during the monsoon season that just ended, pirate attacks off the Somali coast
have increased this year.

In a report released Wednesday, the London-based International Maritime Bureau documented 47 attacks
during the first nine months of 2009, up from 12 during the same period last year. Despite the spike
in attempted raids, only one in nine was actually successful, compared with one in six last year,
according to Cyrus Mody, the bureau's manager.

``It's primarily because of increased cooperation between naval vessels and merchant ships,''
Mr.~Mody said in a telephone interview. ``The masters of the ships have also learned how to prevent
pirates from boarding until help can arrive.''

On Tuesday, China issued strong words, suggesting that it might attempt to rescue those captured.

``We will continue to follow developments closely and make all-out efforts to rescue the hijacked
ship and personnel,'' said Ma Zhaoxu, a Foreign Ministry spokesman. Last year China sent three
warships to the region, although they, like most patrol vessels, stay within the Gulf of Aden.

Shipping experts said the pirates now in command of the De Xin Hai would probably continue heading
toward the Somali coast. Once close to shore, the pirates were expected to drop anchor and begin
negotiations with the ship's Chinese owners.

Then the haggling presumably would begin.

``The normal trend has been a negotiated resolution based on the idea that the safety of the crew is
of paramount importance,'' said Commander Harbor of the European Union.

If recent settlements are any guide, the results will be costly. In August, the owner of a German
ship paid \$2.7 million to free a crew of two dozen that included Filipinos, Russians and
Ukrainians.

\section{Leader of China's Uighur Minority Builds a Stage Across the Globe}

\lettrine{I}{n}\mycalendar{Oct.'09}{22} what has become a familiar vocal pas de deux, Rebiya Kadeer,
the exiled Uighur leader, stepped off a plane in Tokyo on Tuesday and immediately began accusing the
Chinese government of secretly executing members of the Uighur minority and illegally detaining
hundreds of others.

``I wish the killing would stop,'' she said, her braided gray hair topped by a distinctive square
hat. Her words, spoken in the Uighur language, were instantly picked up by international news
agencies and broadcast by the Japanese media.

The Chinese Foreign Ministry immediately fired back, condemning Japan for granting Ms.~Kadeer, 62, a
visa for her weeklong visit, much of which will be devoted to giving speeches on what she says is
China's suppression of the country's Uighurs, who make up the largest ethnic group in the
northwestern region of Xinjiang.

To China, she is a terrorist and the unseen hand behind rioting in Xinjiang last July between
Uighurs and Han Chinese that killed 197 people -- most of them Han -- and injured 1,600.

``Some forces in Japan have planned for Kadeer to go to Japan to engage in separatist activities
aimed at China,'' Ma Zhaoxu, a Foreign Ministry spokesman, said during a regularly scheduled news
conference on Tuesday. ``The Japanese side has ignored China's staunch objections and allowed Rebiya
to enter, and China expresses its strong dissatisfaction.''

Since the unrest last summer, Ms.~Kadeer has become the internationally recognized face of the
Uighur people, a Muslim, Turkic-speaking minority who have long had a contentious relationship with
China's Han majority. Wherever she goes -- from Germany to New Zealand -- she handily draws
attention.

A year ago, however, Ms.~Kadeer was hardly noticed and her cause -- greater autonomy for China's
Uighurs -- largely unknown beyond a small, lonely band of rights advocates. ``Until this year, I
think a lot of Chinese would have had trouble identifying Rebiya Kadeer,'' said Michael Davis, a
professor at the Chinese University of Hong Kong who studies China's relationship with its national
minorities.

Ms.~Kadeer has China to thank for her notoriety, which seems to increase each time she leaves
Washington, her home since being freed from a Chinese prison in 2005 on the condition she go into
exile.

In August, Chinese officials unleashed a fierce publicity campaign against her after the Melbourne
International Film Festival invited her to attend a screening of a documentary about her life.
Diplomats and Chinese citizens tried, and failed, to persuade festival organizers to rescind their
invitation and pull the film. Last week, she barnstormed New Zealand for a few days and then flew to
Germany, where she spoke at the Frankfurt Book Fair, infuriating China, which was the guest of
honor.

``It is just not right to welcome a country where executions are a daily occurrence and human rights
are treated with disrespect,'' she said in a speech there.

For the most part, China's efforts to convince the world that she is a terrorist have failed. To
bolster its contentions, the government released transcripts of a phone conversation in which she
warned her brother to stay off the streets of Urumqi, the regional capital, on the day of an illegal
rally by Uighur students. The demonstration, aimed at drawing attention to the death of two Uighurs
10 days earlier, gave way to the worst rioting in years.

Ms.~Kadeer, who leads the World Uighur Congress, has denied any hand in the unrest, saying Chinese
security forces provoked the murderous rampage by Uighur mobs and the retaliatory violence by Han in
the days that followed.

So far, Chinese courts have handed down 12 death sentences to those convicted of taking part in the
riots. She says that others have been quietly executed without trial.

China's eight million Uighurs have long yearned for a standard-bearer like the Dalai Lama, whose
beatific mien is associated with the aspirations of the Tibetan people. Although she does not have
his stature, his fluency in English or his aura of holiness, Ms.~Kadeer, a wealthy businesswoman
before her imprisonment, has mastered public relations and the well-honed sound bite.

``The Chinese government took a grandmother with 11 children and gave her an international image,''
said J.~Bruce Jacobs, who teaches Asian Studies at Monash University in Australia. ``I think their
plan totally backfired.''

Here in China, few question the official description of Ms.~Kadeer as a terrorist and find the
positive attention she garners abroad confounding and hurtful. Su Hao, director of the Center for
Strategic and Conflict Management at China Foreign Affairs University, blamed the foreign media,
which he said had purposely played down evidence of Ms.~Kadeer's role in organizing the riots. ``The
West uses figures like Rebiya Kadeer to hold China back,'' he said.

Dong Guanpeng, director of the Global Journalism Institute at Tsinghua University in Beijing,
teaches Chinese officials how to grapple with local crises and says that China has to do a better
job getting its message out to foreign audiences. ``We've already proven domestically that she is
the enemy of the Chinese people,'' he said from Chongqing, where he was conducting a training
session for local officials. ``We just have to find a way to make the truth more acceptable to
everyone else,'' he said.

\section{\$13 an Hour? 500 Sign Up, 1 Wins a Job}

\lettrine{A}{s}\mycalendar{Oct.'09}{22} soon as the job opening was posted on the afternoon of
Friday, July 10, the deluge\cite{deluge} began.

C.R. England, a nationwide trucking company, needed an administrative assistant for its bustling
driver training school here. Responsibilities included data entry, assembling paperwork and making
copies.

It was a bona-fide opening at a decent wage, making it the rarest of commodities here in northwest
Indiana, where steel industry layoffs have helped drive unemployment to about 10 percent.

When Stacey Ross, C.~R. England's head of corporate recruiting, arrived at her desk at the company's
Salt Lake City headquarters the next Monday, she found about 300 applications in the company's
e-mail inbox. And the fax machine had spit out an inch-and-a-half thick stack of r\'esum\'es before
running out of paper. By the time she pulled the posting off Careerbuilder.com later in the day, she
guessed nearly 500 people had applied for the \$13-an-hour job. ``It was just shocking,'' she said.
``I had never seen anything so big.''

Ms.~Ross had only a limited amount of time to sort through the r\'esum\'es. While C.~R. England has
not been immune to the downturn, it has added significantly to its stable of drivers and continued
to hire office staff members to support them. Ms.~Ross was also trying to fill more than two dozen
other positions.

The 34-year-old recruiter decided the fairest approach was simply to start at the beginning,
reviewing r\'esum\'es in the order in which they came in. When she found a desirable candidate, she
called to ask a few preliminary questions, before forwarding the name along to Chris Kelsey, the
school's director. When he had a big enough pool to evaluate, she would stop. Anyone she did not get
to was simply out of luck.

She dropped significantly overqualified candidates right away, reasoning that they would leave when
the economy improved. Among them was a former I.B.M. business analyst with 18 years experience; a
former director of human resources; and someone with a master's degree and 12 years at Deloitte \&
Touche, the accounting firm.

Over the course of four days, Ms.~Ross forwarded 61 r\'esum\'es to Mr.~Kelsey, while rejecting 210
others. The remainder never even got a look. Many were, in fact, never uploaded to the company's
internal system because there were too many.

Just before the advertisement was removed, a standard one-page r\'esum\'e arrived from Tiffany
Block, 28, who lived in nearby Portage and had lost her job four months earlier as an accounts
receivable manager at a building company when it closed its Indiana office.

Someone she knew had applied for the job and had said so on Facebook. Ms.~Block went to the
company's Web site and filed an application online, which many others had not. By doing do, her
application went directly into the company's system. She was hardly optimistic, since she had not
had an interview in months.

Ms.~Ross, however, passed it on the next day to Mr.~Kelsey.

Attendance at Mr.~Kelsey's school has surged during the recession. Mr.~Kelsey, 33, had just promoted
one of his three administrative assistants, who handle the paperwork needed for drivers to hit the
road. He needed a replacement quickly.

The overwhelming response astonished him. He asked Cheree Seawood, one of his current assistants, to
go through the r\'esum\'es and help pick out several to interview. To make the task easier, he
decided they should be even more rigorous in ruling out anyone who appeared even slightly
overqualified. Mr.~Kelsey, an ardent New England Patriots fan, compared his personnel strategy to
the team's everyman approach.

``We like to get the fair and middling talent that will work for the wages and groom them from
within,'' he said.

In other words, he said, he did not want the former bank branch manager Ms.~Ross had sent, or the
woman who had once owned a trucking company, or even the former legal secretary.

He also realized that in this climate he could afford to be extra picky and require trucking
industry experience.

The company eventually settled on eight people to interview, inviting in the first two just five
days after the job was posted.

In the past, Mr.~Kelsey had mostly ad-libbed interviews, but this time he asked his company's human
resources department for help. They sent him a list of 13 questions, as well as an eight-page packet
with 128 questions grouped under 50 ``competencies.'' He decided he would ask them all.

At the end of each hourlong interview, he and Ms.~Seawood each jotted down a rating for each
applicant and then compared them.

Invariably, the candidates' job search travails came up. One woman who lost her job had started
working as a waitress and confessed she had come directly from her job on the overnight shift.

But Mr.~Kelsey resolved to keep his personal sympathies at bay. ``If you start judging applicants on
want or need, eventually that want, or need, will go away when they get the job and their financial
situation stabilizes,'' he said. ``Then you're left with whatever skills they have.''

Before Ms.~Seawood called Ms.~Block to schedule an interview, she had been getting increasingly
depressed.

``I felt like, I'm 28 years old, and I don't have a job,'' she said. ``What am I doing with
myself?''

But Mr.~Kelsey was immediately impressed when she came in on the second day of interviews. Dressed
in a conservative business suit, Ms.~Block patiently answered all of the 100-plus questions.
Mr.~Kelsey liked that she remained consistent in her answers and showed independence.

Afterward, Mr.~Kelsey gave Ms.~Block a 9; Ms.~Seawood rated her at a point lower.

The next week, however, Ms.~Seawood gravitated to a different candidate. The woman had just had nose
surgery and came in wearing a protective mask. Besides her qualifications, the fact she had not
tried to postpone impressed Ms.~Seawood.

But when Mr.~Kelsey invited the woman back, the interview was a disaster. She grew visibly irritated
amid his battery of questions.

Mr.~Kelsey immediately called Ms.~Block to ask if she could come in for a second interview.

Was an hour from now too soon?

Momentarily panicked, Ms.~Block quickly assented.

Mr.~Kelsey marched through many of his questions again. Then, trying to gauge her ability to be
assertive among truck drivers, he added a new hypothetical: if she were in the stands at a baseball
game and a foul ball came her way, would she stand up to try to catch it, or wait in her seat and
hope it fell her way?

The other finalist had said she would wait. But Ms.~Block said immediately that she would jump up to
grab it.

Mr.~Kelsey decided he had found his hire.

\section{In Atlanta, String of Black Mayors May Be Broken}

\lettrine{S}{ince}\mycalendar{Oct.'09}{22} 1973, this Southern capital has elected a succession of
black mayors, sometimes to the consternation\cite{consternation} of residents in the largely white,
prosperous neighborhood of Buckhead in the north.

But the current race to succeed Mayor Shirley Franklin in the Nov. 3 election has upended
expectations here in what Chris Rock, in his new documentary, ``Good Hair,'' calls ``the city where
all major black decisions are made.'' The front-runner, Mary Norwood, is one of Buckhead's own, a
white Junior Leaguer running as a populist outsider.

The three top candidates in the six-way race have each maintained that Atlanta has moved beyond
using race as a qualification for public office. But the ascendancy of Ms.~Norwood may also reflect
the decline of the city's black majority and the recession's sour effect on the mood of the voters.

The city has changed significantly since Mayor Franklin squeaked\cite{squeak} to victory without a
runoff in 2001. It has grown by more than 100,000 people since 2000, according to census estimates,
and the influx of many whites and Hispanics has narrowed the black majority to 56 percent from 61
percent.

Atlanta is still a draw for black professionals, and the percentage of blacks in the metropolitan
area has grown slightly, but in the city itself the pool of likely black voters is estimated at just
barely a majority. Many of the city's public housing projects, where black votes once could be
marshaled in a bloc, have been demolished.

Ms.~Norwood, who has held an at-large City Council position for eight years, and tirelessly attended
neighborhood meetings across the city, has galvanized white voters and managed to attract
significant support among blacks. Though she has often voted in Republican primaries in this heavily
Democratic city, some polls show her with more black support than either of her two top opponents,
who are both black: Lisa Borders, the president of the City Council, and Kasim Reed, a lawyer and
former state lawmaker who resigned his office to run for mayor.

The election will probably result in a runoff between Ms.~Norwood and Mr.~Reed or Ms.~Borders,
although Ms.~Norwood is so far ahead in polls that there is talk she could win outright.

``It would be a major game change in this town if a Buckhead Betty became mayor,'' said Tom Houck, a
former newspaper columnist here, who is white, using a mocking term for the well-heeled women of the
north side. ``Atlanta is a symbol for black Americans, more than Los Angeles, more than Chicago,
more than Baltimore.''

Mr.~Houck spoke recently to a primarily black audience at a forum about race in the campaign, where
some of those present were intent on electing a black mayor and others asked what good, exactly,
black leadership had done the city.

Ms.~Norwood has addressed the race question only obliquely, though her campaign photographs and
videos emphasize her interactions with black voters.

``Dr.~King said we should be evaluated on who we are, not what we look like,'' she said in an
interview on Tuesday. ``I'm focused on public safety, city service delivery, quality of life issues
and growing the city. That's what the citizens of Atlanta are interested in.''

Ms.~Norwood has set the tone by relentlessly attacking the Franklin administration's record on crime
and city finances, forcing the other candidates to distance themselves from the mayor.

``When you attack City Hall, you're also implicitly attacking, to a degree, black politics,'' said
Michael Leo Owens, a political science professor at Emory University. ``And this is a message that
in some ways plays well with the white electorate.''

At a time of high anxiety over taxes and crime, it also resonates with voters of all races. The
candidates have spent the bulk of their time in debates and forums arguing over who has the
experience necessary to fix the city's money problems and who has the best public-safety plan.

Some voters, particularly younger ones, seem to agree that race should not be a factor in their
choice at the polls.

``This is a majority white country and Barack Obama's president,'' Tyronia Morrison, a 30-year-old
lawyer who is black, said after a candidate forum in southwest Atlanta. ``We need to rise above that
and get back to the issues.''

Asked if she thought black voters cared about the race of the mayor, Ms.~Borders said: ``Folks that
do not have what they need, who've been marginalized in some way, they want someone that they think,
'Hey, I can relate to them.' If you feel like you've been treated fairly, then your sensitivity is
not that high.''

Ms.~Borders and Mr.~Reed condemned a memo from an ad hoc group, the Black Leadership Forum, that
surfaced in August suggesting that blacks unite behind Ms.~Borders, whom the memo described as the
most electable black candidate.

The controversy over the memo obscured the fact that, as Steve Suitts, an Emory lecturer, wrote in
an op-ed article Wednesday in The Atlanta Journal-Constitution, ``White voters, not black voters,
end up voting most often for a candidate of their own race in the South.''

This mayoral election is the first for an open seat since the death of Maynard Jackson, the city's
first black mayor and a political kingmaker, in 2003.

``Maynard basically was able, because of the esteem in which he was held, to pick his successors,''
said Bob Holmes, the author of a new biography of Mr.~Jackson. ``If he endorsed them, then the black
community united behind them. And he's no longer here.''

Presumably, the mantle of the Jackson machine would have fallen to Mr.~Reed, who ran Ms.~Franklin's
campaigns and who has been endorsed by Andrew Young, who succeeded Mr.~Jackson as mayor. Mr.~Reed
has raised the most money, according to the latest filings, but has lagged behind in the polls.

Both Mr.~Reed and Ms.~Borders have battled voter fatigue after the presidential election and
struggled to differentiate themselves with recession-size campaign treasuries. ``None of them is
considered a charismatic leader,'' Mr.~Holmes said of the three candidates.

That bodes well for Ms.~Norwood, who has positioned herself as the candidate of change.

``In this instance,'' Mr.~Owens said, ``difference equals white.''

\section{Democrats Lose Big Test Vote on Health Legislation}

\lettrine{D}{emocrats}\mycalendar{Oct.'09}{22} lost a big test vote on health care legislation on
Wednesday as the Senate blocked action on a bill to increase Medicare payments to doctors at a cost
of \$247 billion over 10 years.

The Senate majority leader, Harry Reid, Democrat of Nevada, needed 60 votes to proceed. He won only
47. And he could not blame Republicans. A dozen Democrats and one independent crossed party lines
and voted with Republicans on the 53 to 47 roll call.

The Medicare bill has become a proxy for larger issues in the debate over legislation to overhaul
the health care system.

Mr.~Reid said the bill, by averting big cuts in physician fees, guaranteed that doctors would
continue accepting Medicare patients. But since none of the costs were offset or paid for,
Republicans said it was fiscally irresponsible, and some Democrats said they shared that concern.

By addressing doctors' fees in a separate bill, Senate Democrats could hold down the cost of the
broader health legislation, keeping it within the limits set by President Obama. House Democrats are
considering a similar tactic. Republicans said it was a transparent ploy to hide the cost of a
health care overhaul.

Democrats had hoped that by passing the Medicare bill they could appease doctors and secure their
support for the broader legislation.

Senate Democratic leaders said the bill to protect doctors' fees had strong support from the White
House, the American Medical Association and AARP.

Among the Democrats who voted against the party leadership were Senators Evan Bayh of Indiana, Kent
Conrad of North Dakota, Russ Feingold of Wisconsin, Claire McCaskill of Missouri, Bill Nelson of
Florida and Ron Wyden of Oregon.

``I will vote for the doctor fix -- when it is funded,'' Mr.~Nelson said.

Mr.~Wyden said, ``On the eve of a historic debate on health care, it's essential to show a
commitment to real reform,'' which includes fiscal responsibility.

Mr.~Reid said the Republicans had suddenly ``gotten religion'' and were being very
frugal\cite{frugal}. In the past, he said, they did not worry about paying for tax cuts or for the
drug benefit added to Medicare in 2003.

Under current law, doctors face a 21.5 percent cut in Medicare fees in 2010 and then annual 5
percent cuts for several years. Since 2003, Congress has stepped in to postpone such cuts, but it
has usually found ways to offset the cost to the government.

The bill this year, by Senator Debbie Stabenow, Democrat of Michigan, had no offset and would have
repealed the current cost-cutting formula, known as the sustainable growth rate.

The Senate Republican leader, Mitch McConnell, said he felt vindicated by the vote.

``In the Senate's first vote on health care spending this year,'' Mr.~McConnell said, ``a bipartisan
majority rejected the Democrat leadership's attempt to add another quarter-trillion dollars to the
national credit card without any plan to pay for it. With a record deficit and a ballooning national
debt, the American people are saying enough is enough.''

Despite the Senate vote, proponents of sweeping health care legislation moved ahead on other fronts.

At a meeting of the House Democratic Caucus, Speaker Nancy Pelosi indicated that she would push for
a ``robust'' liberal version of a government-run health insurance plan, to compete with private
insurers, if she could get the 218 votes needed to win approval in the full House.

An aide to the House Democratic leadership said Ms.~Pelosi had told the caucus that she had 200
votes, ``or a little over 200,'' for this option, which would use Medicare rates as a basis for
paying hospitals and doctors. Under another option, the government plan would negotiate rates with
providers, as private insurers do.

Ms.~Pelosi said the first alternative saved more money and would give the House leverage in
negotiations with the Senate. But Representative Earl Pomeroy, Democrat of North Dakota, said he
could not vote for this proposal because it would be damaging to his district and was likely to be
dropped in negotiations with the Senate.

In North Dakota, as in many other rural areas, Mr.~Pomeroy said, Medicare rates are far below those
paid by private insurers.

Democrats in both chambers agreed Wednesday on one point: they want to revoke\cite{revoke} the
exemption\cite{exemption} from federal antitrust law that health insurance companies have long
enjoyed.

By a vote of 20 to 9, the House Judiciary Committee approved the antitrust change, which is likely
to be included in the broader health care legislation. It would outlaw price-fixing, bid rigging and
``market allocations'' by companies that sell health insurance or medical malpractice insurance.

Mr.~Reid and the chairman of the Judiciary Committee, Senator Patrick J.~Leahy, Democrat of Vermont,
said they would offer a similar plan as an amendment\cite{amendment} to health legislation in the
Senate.

``Criminal conduct that would land people in jail in other industries is legal when insurers do
it,'' Mr.~Leahy said.

The insurance business has been largely exempt from federal antitrust law since 1945. The Supreme
Court ruled in 1944 that insurance was interstate commerce subject to federal antitrust law. But the
insurance industry won a reprieve nine months later, when Congress passed the McCarran-Ferguson Act.

States regulate health insurance and have sued insurers for anticompetitive conduct. But state
officials said they had received little help from the federal government.

Richard Blumenthal, the attorney general of Connecticut, said: ``I feel strongly that the exemption
should be repealed. It is a legal and historical anomaly\cite{anomaly} that there is such a sweeping
and all-encompassing exemption for an entire industry. It makes no sense as a matter of logic or
law.''

\section{Fate of White House Counsel Is in Doubt}

\lettrine{E}{very}\mycalendar{Oct.'09}{22} morning, Gregory B.~Craig, the White House counsel, sits
at the conference table of the Roosevelt Room with the rest of the president's senior staff members.

His colleagues greet him, friendly as always. He updates the room on his issues and listens as they
update him on theirs.

The one issue that does not come up? Mr.~Craig himself.

As President Obama's top lawyer, Mr.~Craig has been at the center of thorny decisions on closing the
prison at Guant\'anamo Bay, Cuba, and revising interrogation and detention policies, problems that
have bedeviled\cite{bedevil} the new administration and generated fierce battles inside and outside
the White House. And for months now, he has endured speculation in print and around the White House
about whether he is on the way out.

Mr.~Craig said he had no plans to leave, and White House officials said the president still had
faith in him. But colleagues and Democrats close to the White House said they expected him to move
on around the end of the year, and they have been talking about possible replacements. Whether
Mr.~Craig leaves or not, the discussion of his fate has grown so persistent and the situation so
awkward that one supporter calls it ``the Greg Craig watch.''

It is a classic and not particularly savory\cite{savory} Washington story. When an administration
stumbles, whispers begin and fingers point in search of someone to blame. At a certain point,
assumptions can become self-fulfilling, and an official in the cross hairs finds it harder to do the
job. In Mr.~Craig's case, friends said he was unfairly being made a scapegoat for decisions
supported across the administration.

``Everybody who's here in Washington has seen this play before, when internal second-guessing and
rivalries burst out into public,'' said Elisa Massimino, executive director of Human Rights First,
an advocacy group that has worked with Mr.~Craig to reverse President George W.~Bush's policies.

``What he did was jump in with both feet,'' Ms.~Massimino added. ``Had he not done that, we would
not be nearly so close as we are to finally turning this problem around.''

Mr.~Craig, a former adviser to Senator Edward M.~Kennedy and President Bill Clinton, is a savvy
capital veteran. He declined to comment on his plans, but told the National Law Journal that he was
staying. Attorney General Eric H.~Holder Jr., who called Mr.~Craig ``extremely effective,'' said he
had maintained equanimity despite the anonymous sniping.

``Having been in Washington as long as he has been, there's a familiarity to this, an inoculation
against it,'' Mr.~Holder said in an interview. ``So that helps, I'm sure, in dealing with this
situation.''

Mr.~Craig was handed one of the most difficult portfolios at the White House. He drafted executive
orders banning torture and ordering the Guant\'anamo prison closed within a year. Over the
objections of the Central Intelligence Agency, he recommended the release of Justice Department
memorandums describing harsh interrogations. And he was at the center of the White House decision to
reverse itself and withhold photographs of abuse of detainees.

Mr.~Craig took much criticism for those decisions, for not doing more to build consensus within the
administration or prepare the political ground in Congress. And the prospect of closing Guant\'anamo
by the one-year deadline he helped set has proved more difficult than anticipated and now appears
unlikely to be met.

At moments, it has looked as if Mr.~Craig's authority has been trimmed back. Rahm Emanuel, the White
House chief of staff, assigned Pete Rouse, a senior adviser with deep ties to Capitol Hill, to
oversee Guant\'anamo issues.

Similarly, after Mr.~Craig started the search that produced the Supreme Court nomination of Justice
Sonia Sotomayor, Mr.~Emanuel assigned the confirmation fight to Ronald A.~Klain and Cynthia Hogan,
aides to Vice President Joseph R.~Biden Jr.~with long experience handling judicial appointments.

In both instances, White House officials said that Mr.~Craig remained involved but that it made
sense to tap people with political backgrounds to manage political issues, particularly since
Mr.~Craig had so many other duties, like scrutinizing legislation, vetting appointees and selecting
judges. But articles in The Wall Street Journal, The Washington Post and Politico questioned
Mr.~Craig's job security, citing anonymous critics.

In an interview, Mr.~Emanuel said Mr.~Craig should not be faulted, especially on Guant\'anamo.

``The president believes he has done a very good job and continues to do a very good job,''
Mr.~Emanuel said. ``The notion that you're going to blame him is ridiculous. He didn't create
Guant\'anamo. He is trying to work within the system to meet the president's goal.''

Many in the White House are angry at the whispering campaign. Mr.~Craig, 64, is well liked, a
veteran Washington lawyer with an easy smile, a friendly manner and thick wavy hair that makes him
look like a member of the Kennedy family he once served.

He studied law at Yale with Bill and Hillary Rodham Clinton and joined the Clinton White House in
1998 to fight impeachment. Longtime aides resented\cite{resent} the newcomer. When the announcement
of his appointment described Mr.~Craig as the ``quarterback'' of the impeachment\cite{impeachment}
defense, some Clinton aides, including Mr.~Emanuel, derisively referred to him as ``QB.'' (All these
years later, Mr.~Emanuel said he liked and respected Mr.~Craig.)

Mr.~Craig later broke with the Clintons to become one of Mr.~Obama's earliest backers for president.
He told Newsweek he was ``blown away'' by Mr.~Obama's books; he was such an enthusiast that
Mr.~Obama referred to him as one of the ``Kool-Aid kids.'' After the election, Mr.~Obama made him
White House counsel, overruling advisers who thought he was not the best fit or who favored other
candidates.

Mr.~Craig was hardly the only adviser advocating a one-year deadline to close the Guant\'anamo
prison; Defense Secretary Robert M.~Gates, among others, also argued for it. ``That's how you force
the action and get things done, and the truth is we're making good progress on it,'' said David
Axelrod, Mr.~Obama's senior adviser. ``But the thing is, it was a legal and policy thicket.''

Tom Malinowski, Washington director of Human Rights Watch, said that while the White House had made
mistakes in the closing of Guant\'anamo, fueling an uproar in Congress, it would be unfair to assign
those mistakes to Mr.~Craig.

``To make Greg the fall guy, if that indeed is what they're doing, is profoundly disingenuous,''
Mr.~Malinowski said.

Whatever the motivation, the talk does its damage.

``In the White House, in particular, power is the perception of power,'' said Bradford A.~Berenson,
a former associate White House counsel under Mr.~Bush. ``And if people perceive you to be under
attack and on the way out, it limits your effectiveness.''

\section{Eyeing a Democratic Win, the President Joins Corzine in New Jersey}

\lettrine{H}{oping}\mycalendar{Oct.'09}{22} to spare his party an embarrassing
defeat in one of the year's most closely watched elections, President Obama campaigned on Wednesday
evening alongside New Jersey Gov. Jon S.~Corzine, saying that a Democratic victory in the state
would help the president carry out the political change he was elected to bring to Washington.

Mr.~Corzine is the first incumbent governor to seek re-election since the economic crisis hit last
fall and is locked in a tight race with the Republican nominee, Christopher J.~Christie, who has
blasted Mr.~Corzine for New Jersey's rising taxes and unemployment.

But Mr.~Obama told a crowd of more than 3,000 at Fairleigh Dickinson University that it is unfair to
blame Mr.~Corzine for a global recession that has battered every state in the nation. And he praised
Mr.~Corzine for helping mitigate\cite{mitigate} the pain the downturn has inflicted on middle-class
families by preserving funding for education and health care and making New Jersey one of the first
states to start its own economic stimulus plan.

``For the last four years, you've had a leader who's put the interests of hard-working New Jersey
families ahead of the special interests,'' he said, calling Mr.~Corzine one of the ``best partners I
have in the White House.''

Mr.~Obama also accused Republicans of ``selective memory'' by ignoring the fact that the economy
collapsed during the final month of President George W.~Bush's term, and warned that Mr.~Christie
would reimpose the same ``trickle down'' economic policies that led to the downturn.

``They made the mess, and now they're complaining about how fast we're cleaning it up,'' Mr.~Obama
said, bringing a roar of laughter and applause from the crowd.

Mr.~Corzine was clearly buoyed by the support from Mr.~Obama, who won New Jersey handily last fall
and remains highly popular. Speaking before the president, Mr.~Corzine promised to be brief
``because I know who you're here to see,'' and used the joint appearance to align himself closely
with Mr.~Obama's policies.

Mr.~Corzine said that, like Mr.~Obama, he had fought for expanded health care programs, and
increased funding for early education and alternative energy.

``I'm proud to be President Obama's partner in building a better America and certainly a better New
Jersey,'' he said.

The appearance was part of a political blitz for Mr.~Obama, who appeared in New York on Tuesday to
raise more than \$3 million for two Democrats and was scheduled to stump in Massachusetts and
Virginia next week.

Mr.~Christie was careful not to criticize Mr.~Obama, but said that New Jersey voters were
savvy\cite{savvy} enough to judge Mr.~Corzine on his own record.

``Air Force One is leaving tonight,'' he said in a television appearance on Wednesday. ``If they
vote for Jon Corzine, President Obama isn't moving in to run the state.''

Mr.~Corzine's political fortunes have improved considerably since Mr.~Obama's first campaign
appearance on his behalf in June. For much of the spring and summer, Mr.~Christie held a
double-digit lead in the polls, and some Democrats openly wondered whether the governor should be
replaced at the top of the state party's ticket.

But Mr.~Corzine has used his financial advantage to close the gap in recent weeks, and one key
factor in the Nov. 3 election will be his ability to motivate the 600,000 first-time voters who cast
ballots for Mr.~Obama last year.

With the Democrats' prospects having grown dimmer in Virginia -- the only other governor's race this
fall -- a procession of Obama aides and administration officials have come to New Jersey to bolster
Mr.~Corzine in recent weeks -- including Vice President Joseph R.~Biden Jr.~and Valerie Jarrett, a
senior adviser to Mr.~Obama.

Mr.~Obama urged audience members to rekindle the verve they felt during his own campaign last year
and rally behind Mr.~Corzine, and he tried to make a virtue of the governor's image as an unpolished
politician. Conceding that Mr.~Corzine's ``frizzy'' hair and ``scraggly beard'' were the antithesis
of the slick politician, Mr.~Obama said that the governor been an honorable public servant with the
courage to level with the public about the problems ahead.

``I'm going to ask you to cast aside the cynics and the skeptics and show that leaders who do what's
right and what's hard are rewarded and not rejected,'' Mr.~Obama said.

Earlier in the day, Mr.~Obama arrived at Newark Liberty International Airport without incident. But
on Tuesday, Secret Service officials said that they arrested a private security guard and charged
him with making threats to harm the president.

\section{2 Deals Hint at Revenue for Twitter}

\lettrine{T}{witter}\mycalendar{Oct.'09}{22} gets 55 million monthly visitors, it has raised \$155
million in venture capital, and it has generated intense interest from Hollywood to Iran. But it
hasn't earned much revenue and certainly no profit.

Back-to-back deals on Wednesday to make the company's steady stream of posts available to Microsoft
and Google's search engines may point to a potential new source of cash. How large, however, is not
known. The terms of the deals were not disclosed and Evan Williams, Twitter's chief executive, said
in an interview that revenue was ``not the focus of the deals.''

Microsoft said it did not plan to put ads on its Twitter search service for now, and Google said ads
might appear at a later date.

The deals represent the latest evidence of the intense interest in what is known as the real-time
Web -- the constant stream of posts and updates on Twitter, Facebook and similar services. Unlike
traditional Web pages and blogs, that real-time information has not been easily integrated by search
engines.

Microsoft has already included Twitter data in a service of its search engine, Bing. It demonstrated
the service at Web 2.0, a technology conference in San Francisco. Google said that it would offer a
similar feature soon.

The deals are not exclusive and fit into Twitter's approach to doing business, Mr.~Williams said. He
raised the possibility of reaching similar agreements with other companies.

``A core of our philosophy has always been that Twitter is a distributed network and there's
multiple in points and out points that serve different users and different uses,'' Mr.~Williams
said.

Unlike most Web companies, which seek to drive users to their Web sites, Twitter has said it does
not care whether users see it through third-party applications on phones or computers or through
sites like Bing.

``It's `let a thousand flowers bloom,' '' Mr.~Williams said. ``Other folks will attack the problem
differently and present it in a different context that makes sense for different use cases.''

On Wednesday, Microsoft also said that it had reached a separate agreement to soon include status
updates from Facebook in Bing, but it gave few details of what that service would look like.

Sheryl Sandberg, Facebook's chief operating officer, said that only status updates that were public
would be shared with Microsoft.

Several search startups, including Twitter itself, have search services tailored to uncover Twitter
posts and other ``real-time'' data. But none has yet figured out how to mine the most relevant
results in the way companies like Google have for Web search, Mr.~Williams said.

Twitter has experimented with one way to do it, showing users which topics people are most actively
talking about on the site.

Bing will experiment with different approaches to show users the most relevant results by filtering
out duplicates and trying to rank posts in order of importance based on who their authors are and
other factors.

``We are bringing you the best of real time right into the Web results,'' said Yusuf Mehdi, senior
vice president for Microsoft's online audience business group, who demonstrated the service at the
conference.

For now, Bing users interested in searching Twitter will use a separate service, at
bing.com/twitter. Eventually, the company plans to further integrate Twitter search into Bing.

Google plans to offer a dedicated service to search Twitter and to blend Twitter posts into its main
search results.

Marissa Mayer, Google's vice president for search and user experience, said that access to Twitter's
trove of real-time data would ``improve our relevance, our comprehensiveness and our quality.''

\section{The Billion Designers of Windows 7}

\lettrine{T}{here}\mycalendar{Oct.'09}{22} is a scene in the pilot episode of ``Mad Men'' when the
creative genius Don Draper suggests in a client meeting that a campaign for Lucky Strike use the
slogan ``It's toasted'' to play up how the cigarettes are made. A junior Lucky Strike executive
responds, ``But everybody else's tobacco is toasted'' -- an argument brushed aside by the senior
executive in the room.

And a pioneer of modern advertising, Claude C.~Hopkins, created ads for Schlitz beer that proclaimed
its bottles were ``washed with live steam.'' That step was a standard part of the brewing process,
but it became the focus of what was a hugely successful campaign.

Likewise, the Microsoft Corporation regularly asks PC users for feedback about its products. But
after the debacle with Vista, the operating system nobody liked, the company and its advertising
agency, Crispin Porter \& Bogusky, realized that the concept of consumers as an
intrinsic\cite{intrinsic} part of the development process could be an effective selling point for
the Vista replacement, Windows 7.

And so was born a campaign, getting under way on Thursday in six countries, carrying the theme ``I'm
a PC and Windows 7 was my idea.'' It is to run at least through the end of the year, to take
advantage of the computer buying that usually takes place during the holiday shopping season.

The campaign is extensive, like the ``I'm a PC'' campaign, also by Crispin Porter, part of MDC
Partners, which began in September 2008, and the campaign to introduce the Bing search engine, by
JWT, a unit of WPP, that started in June.

•

There will be television commercials for Windows 7, along with print ads, banner ads online, outdoor
posters and branded-entertainment elements with the Fox Broadcasting division of the News
Corporation. Those include a 30-minute variety special on Nov. 8, produced by the creator of
``Family Guy,'' Seth MacFarlane, and tailored segments during pre- and postgame football
programming.

For instance, Windows 7 will sponsor a seven-word recap of a game after it ends. (Take that,
Twitter.)

There are other seven-centric aspects to the campaign, among them video clips on YouTube and
windows.com, each seven seconds long, demonstrating Windows 7 features, and an invitation to
computer users to create similar seven-second videos about features they like.

That invitation, of course, also reinforces the idea that Windows 7 was, as they say in Ziploc ads,
``designed with you in mind.'' Another way that is expressed is by featuring actual Microsoft
customers -- as well as employees -- in the ads.

And many ads use a slogan, ``1 billion = 7,'' suggesting that the billion people who use PCs helped
bring forth the new operating system.

``Our customers co-create the product with us,'' said David Webster, general manager for brand and
marketing strategy at Microsoft in Redmond, Wash. ``We're using the customers' voice to tell our
story.''

Mr.~Webster summarized the points the ads make: ``You told us you want it simpler, we made it
simpler. You told us you want it to boot faster, it boots faster.''

In one poster ad, these words are superimposed over a photograph of a woman: ``I asked for it to use
less memory. Now it uses less memory. I'm a tech goddess.''

In another poster ad, these words appear over a photo of an older man: ``I suggested they make it
less complicated. Guess what? Now it's less complicated. I so rule.''

In commercials, Microsoft engineers say, ``Bring it on; what do you got?'' PC users fire back with
pithy phrases like ``Less clutter, just less clutter.'' And the engineers reply: ``Loud and clear.
We're all over it.''

A top creative executive at Crispin Porter acknowledged the tried-and-true approach the campaign
takes.

``The whole promise, 'We've heard you,' it's not a brand-new technique,'' said Rob Reilly, partner
and co-executive creative director at Crispin Porter, based in the Boulder, Colo., office. ``But a
big company listening to consumers is the way the world is going.''

``It's about how software for the people is now software for the people and by the people,'' he
added.

The campaign does not address the epic fail of Vista -- or as Mr.~Reilly put it, ``the issues that
came with the previous operating system.''

Still, it provides a subtext for the ads in that ``there's a heightened sense of getting this
right,'' he said.

The Windows 7 campaign could benefit from widespread perceptions that the bing.com ads and the ``I'm
a PC'' campaign are working.

The ``numbers are up significantly'' in research about Microsoft, Mr.~Webster said, as more
consumers say they agree with statements like ``Windows is easy to use,'' ``PCs are a good value''
and ``I feel good about this brand.''

``Having an ongoing communication with our customers has had a large impact on overall
favorability,'' he added.

And early coverage of Windows 7 has been positive, carrying headlines like ``Win 7: Microsoft Gets
It Right (Finally),'' in BusinessWeek, and ``Microsoft Reboots,'' in Fortune.

The reaction at Crispin Porter was ``Wow, this is a product we can really stand behind,'' Mr.~Reilly
said, compared with Vista, which ``wasn't the right product'' to help ``build pride back into PC
users.''

•

When Mr.~Reilly was asked about any similarity between what he and Don Draper do for clients, he
laughed and said, ``I don't want to be called Don Draper; my wife will be reading this.''

Still, the Windows 7 campaign, like the Lucky Strike campaign, is about the effort to ``find the
truth in what's part of the brand DNA,'' Mr.~Reilly said, supported in this instance by ``the
authenticity of being able to take real people from around the world and tell their stories.''

The six countries where the Windows 7 ads are to appear initially are Australia, Britain, Canada,
France, Germany and the United States. Others are planned to follow in several weeks.

\section{A Reversal for eBay as Shoppers Return and Revenue Rises, but Profit Still Slides}

\lettrine{E}{bay}\mycalendar{Oct.'09}{22}, the biggest shopping site on the Web, benefited in the
third quarter as consumers gingerly started spending again, but investors were not convinced that
the company had pulled out of its slump.

EBay reported that revenue grew 6 percent, to \$2.2 billion, in the third quarter from the period a
year ago, ending a yearlong run of declining sales.

However, net income fell for a fourth quarter, declining 29 percent, to \$350 million, or 27 cents a
share, from \$492 million, or 38 cents a share. EBay shares were down 4 percent in after-hours
trading.

John J.~Donahoe, eBay's chief executive, attributed the earnings decline to a one-time tax benefit
in the third quarter of last year, lower interest income on eBay's cash and the strong dollar.
Extracting those factors, earnings were flat, he said in an interview.

``These are strong results for a strong company getting stronger,'' Mr.~Donahoe said in a call with
analysts. ``Our strategies are working. EBay has stabilized.''

EBay, which is based in San Jose, Calif., has been on an 18-month mission to turn around its
struggling marketplaces unit, which includes the eBay auction site and other e-commerce sites like
Shopping.com.

That effort is finally bearing fruit, said Sandeep Aggarwal, an analyst at Collins Stewart. ``We
think their marketplace is turning around, and it is positively impacting the business,'' he said.

The total of all transactions except for cars on the eBay Marketplace, an important gauge of growth
at eBay known as gross merchandise volume, increased 7 percent from the same quarter last year,
after significant declines in previous quarters.

EBay now offers free shipping on more items and lower fees for sellers. It has changed its search
algorithm so that listings from high-quality merchants show up higher on the list of results when a
shopper searches for a product. EBay is also doing more fixed-price sales instead of auctions.

Still, not all e-commerce analysts are confident that eBay's turnaround will be successful.

``EBay just kind of reached its saturation point,'' said Sucharita Mulpuru, a vice president at
Forrester, the research firm. ``The universe of people who are willing to buy in a marketplace
setting like eBay is limited.''

EBay's online payments business, which consists of PayPal and Bill Me Later, continued to drive its
growth. Revenue grew 15 percent, to \$688 million.

On Nov. 3, PayPal will open its platform to software developers who want to build applications using
PayPal's technology. Scott Thompson, PayPal's president, said Tuesday that he hoped this would
unleash innovation in the \$30 trillion payments industry

Applications could be as simple as a way to send money to friends over a social network or as
complex as a payroll system for employees in other countries, Mr.~Thompson said.

Revenue at Skype, the online calling service, increased 29 percent, to \$185 million.

In September, eBay announced that it had reached an agreement to sell Skype to a group of investors
led by Silver Lake, a private equity firm. Though Skype's founders are trying to scuttle the deal
with several lawsuits, Mr.~Donahoe said that the company expected the deal to close in the next
quarter.

\section{Yahoo's Profit Triples Despite Sales Decline}

\lettrine{F}{or}\mycalendar{Oct.'09}{22} much of her tenure as chief executive of Yahoo, Carol Bartz
warned investors to expect a steep decline in the company's business. It has not been as bad as she
predicted.

Aided by cost-cutting, the sales of some assets and better-than-expected revenue from display ads,
Yahoo's net income more than tripled in the third quarter. The results topped analysts forecasts,
and investors sent Yahoo's shares sharply higher in after-hours trading.

But revenue fell 12 percent, a sign that many advertisers remain nervous about the pace of economic
recovery. And the results showed that Yahoo was not benefiting from the budding economic rebound as
much as its rival Google.

Analysts said that the results were encouraging but cautioned that Ms.~Bartz's turnaround plan
remained a work in progress.

``The patient is off life support and back in the recovery ward,'' said Jeffrey Lindsay, an analyst
with Sanford C.~Bernstein \& Company. ``But it is certainly not out playing soccer again.''

Yahoo executives, however, hailed the results as a relative success. ``We had a solid third quarter
that signals our major businesses have stabilized,'' Ms.~Bartz said in a statement.

Yahoo said that revenue fell to \$1.57 billion, from \$1.79 billion a year ago. Net revenue, which
excludes commissions paid to advertising partners, was \$1.13 billion, down from \$1.32 billion last
year, and slightly up from analysts forecasts of \$1.12 billion. Net income was \$186 million, or 13
cents a share, up from \$54 million, or 4 cents a share, a year ago, a 244 percent increase. On
average, Wall Street analysts expected that Yahoo would earn 7 cents a share.

Yahoo said that excluding the impact of currency fluctuations and of products it abandoned, revenue
would have fallen 7 percent. In contrast, Google reported that its third-quarter revenue grew 7
percent.

Ms.~Bartz has bristled in the past at the comparison with Google, saying the two companies were in
different markets. While Google derives a large majority of its revenue from search and other text
ads, which have held up relatively well during the recession, Yahoo's business is built on a mix of
search and display ads. Its results tend to be a better indicator of the overall health of the
online advertising market.

But some analysts said they were particularly concerned that Yahoo's search business remained in the
doldrums. The company's revenue from search ads on its own sites was down 19 percent from a year
earlier, a stark contrast to Google's relatively strong growth.

``Something isn't working right in their search business,'' said Christa Quarles, an analyst with
Thomas Weisel Partners. Ms.~Quarles said that if Yahoo's search business continued to erode, Yahoo's
plans for a turnaround could be derailed.

In July, the company signed a blockbuster search and advertising partnership with Microsoft, which
will take over the costly process of indexing and ranking Web sites. Using Microsoft's search
results, Yahoo said it planned to continue offering its own search service, with different features
and a different user interface from Microsoft's Bing.

The deal, which is under review by antitrust regulators, is not expected to close until early next
year.

Yahoo's display advertising business fell by 8.3 percent compared to last year, a lesser drop than
the 14 percent decline in the second quarter. And the business has grown for two consecutive
quarters.

Ms.~Bartz did not participate in a conference call with investors, as is customary, because she was
ill, said Tim Morse, Yahoo's chief financial officer. In a interview, Mr.~Morse said that
advertisers tightened their budgets significantly at the beginning of the year. ``In the second
quarter we saw a loosening. We were very encouraged that the third quarter continued that trend.''

Analysts noted that Yahoo outperformed expectations in part because efforts to get rid of ``low
quality'' ads cost the company far less than it predicted three months ago.

During the quarter, Yahoo also unveiled a major ad campaign aimed at revitalizing the company's
tarnished brand.

The ``It's You'' tagline for the campaign, which is being backed with more than \$100 million in
television, print, outdoor and online ads, is intended to emphasize that some of Yahoo's most
popular products, like its home page, mail and search services, have become easier to personalize.

Yahoo shares closed at \$17.17, down 5 cents, on Tuesday. In after hours, they jumped more than 5
percent, topping \$18.

\section{A New Electronic Reader, the Nook, Enters the Market}

\lettrine{A}{s}\mycalendar{Oct.'09}{22} widely expected, Barnes \& Noble unveiled its Nook
electronic reading device at a splashy news conference on Tuesday to generally positive views from
the publishing community, and offered some details about its whispered-about lending capabilities.

As much as anything, publishers seemed relieved that Barnes \& Noble, which operates the nation's
largest chain of bookstores, had produced a credible alternative to Amazon's Kindle. The Nook,
priced at \$259, went on sale Tuesday afternoon at nook.com, at a price that matched the latest
edition of the Kindle. The Nook will ship starting in late November.

Amazon currently dominates the market for electronic readers. Estimates vary, but according to the
Codex Group, a consultant to the publishing industry, Amazon has sold about 945,000 units, compared
with 525,000 units of the Sony Reader.

Barnes \& Noble opened an e-bookstore in July, and its editions, which are available in ePub and
Adobe PDF versions, can be read on a variety of devices, including Apple's iPhone, the BlackBerry,
Macs and PCs. Barnes \& Noble will continue to support those devices, as well as forthcoming
e-readers from iRex and Plastic Logic.

But it is clear the company is trying to consolidate sales of e-books onto the Nook, which features
a six-inch gray and white reading screen and a color touch screen control panel. In any of the
chain's 1,300 stores, consumers can download books on the Wi-Fi network. Outside the stores,
consumers will access AT\&T's 3G network to download books.

One of the differentiating factors of the Nook is that customers can ``lend'' books to friends. But
customers may lend out any given title only one time for a total of 14 days and they cannot read it
on their own Nook while it is lent.

In an interview, William Lynch, president of Barnes\&Noble.com, said the company would aggressively
market the Nook within its bricks and mortar stores. The Nook also has software that will detect
when a consumer walks into a store so that it can push out coupons and other promotions like
excerpts from forthcoming books or suggestions for new reading. While in stores, Nook owners will be
able to read any e-book through streaming software.

\section{For Decades, Puzzling People With Mathematics}

\lettrine{F}{or}\mycalendar{Oct.'09}{22} today's mathematical puzzle, assume that in the year 1956
there was a children's magazine in New York named after a giant egg, Humpty Dumpty, who purportedly
served as its chief editor.

Mr.~Dumpty was assisted by a human editor named Martin Gardner, who prepared ``activity features''
and wrote a monthly short story about the adventures of the child egg, Humpty Dumpty Jr.~Another
duty of Mr.~Gardner's was to write a monthly poem of moral advice from Humpty Sr. to Humpty Jr.

At that point, Mr.~Gardner was 42 and had never taken a math course beyond high school. He had
struggled with calculus and considered himself poor at solving basic mathematical puzzles, let alone
creating them. But when the publisher of Scientific American asked him if there might be enough
material for a monthly column on ``recreational mathematics,'' a term that sounded even more
oxymoronic in 1956 than it does today, Mr.~Gardner took a gamble.

He quit his job with Humpty Dumpty.

On Wednesday, Mr.~Gardner will celebrate his 95th birthday with the publication of another book --
his second book of essays and mathematical puzzles to be published just this year. With more than 70
books to his name, he is the world's best-known recreational mathematician, and has probably
introduced more people to the joys of math than anyone in history.

How is this possible?

Actually, there are two separate puzzles here. One is how Mr.~Gardner, who still works every day at
his old typewriter, has managed for so long to confound and entertain his readers. The other is why
so many of us have never been able to resist this kind of puzzle. Why, when we hear about the guy
trying to ferry a wolf and a goat and a head of cabbage across the river in a small boat, do we feel
compelled to solve his transportation problem?

It never occurred to me that math could be fun until the day in grade school that my father gave me
a book of 19th-century puzzles assembled by Mr.~Gardner -- the same puzzles, as it happened, that
Mr.~Gardner's father had used to hook him during his school days. The algebra and geometry were
sugar-coated with elaborate stories and wonderful illustrations of giraffe races, pool-hall
squabbles, burglaries and scheming carnival barkers. (Go to nytimes.com/tierneylab for some
examples.)

The puzzles didn't turn Mr.~Gardner into a professional mathematician -- he majored in philosophy at
the University of Chicago -- but he remained a passionate amateur through his first jobs in public
relations and journalism. After learning of mathematicians' new fascination with folding certain
pieces of paper into different shapes, he sold an article about these ``flexagons'' to Scientific
American, and that led to his monthly ``Mathematical Games'' column, which he wrote for the next
quarter-century.

Mr.~Gardner prepared for the new monthly column by scouring Manhattan's second-hand bookstores for
math puzzles and games. In another line of work, that would constitute plagiarism, but among puzzle
makers it has long been the norm: a good puzzle is forever.

For instance, that puzzle about ferrying the wolf, the goat and the cabbage was included in a puzzle
collection prepared for the emperor Charlemagne 12 centuries ago -- and it was presumably borrowed
by Charlemagne's puzzlist. The row-boat problem has been passed down in cultures around the world in
versions featuring guards and prisoners, jealous spouses, missionaries, cannibals and assorted
carnivores.

``The number of puzzles I've invented you can count on your fingers,'' Mr.~Gardner says. Through his
hundreds of columns and dozens of books, he always credited others for the material and insisted
that he wasn't even a good mathematician.

``I don't think I ever wrote a column that required calculus,'' he says. ``The big secret of my
success as a columnist was that I didn't know much about math.

``I had to struggle to get everything clear before I wrote a column, so that meant I could write it
in a way that people could understand.''

After he gave up the column in 1981, Mr.~Gardner kept turning out essays and books, and his
reputation among mathematicians, puzzlists and magicians just kept growing. Since 1994, they have
been convening in Atlanta every two years to swap puzzles and ideas at an event called the G4G: the
Gathering for Gardner.

``Many have tried to emulate him; no one has succeeded,'' says Ronald Graham, a mathematician at the
University of California, San Diego. ``Martin has turned thousands of children into mathematicians,
and thousands of mathematicians into children.''

Mr.~Gardner says he has been gratified to see more and more teachers incorporating puzzles into the
math curriculum. The pleasure of puzzle-solving, as he sees it, is a happy byproduct of evolution.

``Consider a cow,'' he says. ``A cow doesn't have the problem-solving skill of a chimpanzee, which
has discovered how to get termites out of the ground by putting a stick into a hole.

``Evolution has developed the brain's ability to solve puzzles, and at the same time has produced in
our brain a pleasure of solving problems.''

Mr.~Gardner's favorite puzzles are the ones that require a sudden insight. That aha! moment can come
in any kind of puzzle, but there's a special pleasure when the insight is mathematical -- and
therefore eternal, as Mr.~Gardner sees it. In his new book, ``When You Were a Tadpole and I Was a
Fish,'' he explains why he is an ``unashamed Platonist'' when it comes to mathematics.

``If all sentient beings in the universe disappeared,'' he writes, ``there would remain a sense in
which mathematical objects and theorems would continue to exist even though there would be no one
around to write or talk about them. Huge prime numbers would continue to be prime even if no one had
proved them prime.''

I share his mathematical Platonism, and I think that is ultimately the explanation for the appeal of
the puzzles. They may superficially involve row boats or pool halls or giraffes, but they're really
about transcendent numbers and theorems.

When you figure out the answer, you know you've found something that is indisputably true anywhere,
anytime. For a brief moment, the universe makes perfect sense.

\section{Flu Story: A Pregnant Woman's Ordeal}

\lettrine{T}{he}\mycalendar{Oct.'09}{22} group most threatened by swine flu
and most in need of the new vaccine, world health authorities agree, is that of pregnant women.

For example, Aubrey Opdyke.

On June 27, Ms.~Opdyke, a 27-year-old waitress and former high-school swimmer who weighed 135 pounds
before her pregnancy and had no health risks other than a smoking habit, came down with mild flu
symptoms.

She finally came home from the hospital three weeks ago.

``At first, I didn't think anything of it -- just another flu bug,'' Ms.~Opdyke said recently. ``But
it really wrecked me. I probably shouldn't have made it.''

In the four months she was hospitalized, she spent five weeks in a coma, suffered six collapsed
lungs and a near-fatal seizure. High-pressure ventilation blew her up like a molten balloon until
``she looked like she weighed 400 pounds,'' her husband, Bryan, said, and she has stretch marks from
her neck to her ankles. Her muscles and lungs are still so weak that she uses a walker.

While hospitalized, she missed seeing her 4-year-old daughter, Hope, learn to swim and start
pre-school.

And, most important, she lost her baby. Parker Christine Opdyke, almost 27 weeks in the womb, was
delivered by emergency Caesarean section on July 18, when her fetal heart rate plummeted during
Ms.~Opdyke's third lung collapse. Her airways were too blocked to let a breathing tube in, possibly
a side effect of the drugs saving her mother. She lived seven minutes.

On Oct.~1, the Centers for Disease Control and Prevention said 100 pregnant women had been in
intensive care with swine flu and 28 had died. That is a tiny fraction of what are believed to have
been millions of cases in the country. But it is the best argument, federal officials say, for the
drawn-out, expensive effort to make a swine flu vaccine.

Pregnant women are particularly susceptible because they are in the younger age group most likely to
catch this new virus, while those over 50 who have had more flus rarely catch it. Moreover,
pregnancy suppresses the immune system to protect the fetus, and the growing baby makes it harder
for a mother to clear her lungs.

All but a few of the pregnant women who have died or been near death from swine flu are unknown.
Privacy laws prevent health departments from releasing their names, and few families come forward.

The Opdykes did because ``we wanted to get it out there how dangerous it can be,'' Ms.~Opdyke said.

``We have friends who get flu symptoms and say, `Oh, I'm not going to a doctor,' '' Mr.~Opdyke
added. ``And we say, `Do you not understand what we went through?' I can't imagine why there's so
much nonchalance\cite{nonchalance}.''

That nonchalance strikes close to home.

As they said this, Ms.~Opdyke was doing her daily physical therapy, struggling to lift one-pound
weights. Her therapist interrupted to announce that she opposed flu shots.

``Have you ever read the labels?'' she asked. ``They're so full of toxins.''

Asked if she realized that a shot, had it existed in June, might have saved her client and her baby,
she frowned and went back to her clipboard.

Unlike some other families that came forward, the Opdykes are not threatening to sue anyone.

They do not blame her obstetrician\cite{obstetrician}, even though she suggested acetaminophen the
first time Ms.~Opdyke called her and prescribed an antibiotic the second.

``Swine flu just wasn't on our minds at all,'' Ms.~Opdyke said.

Nor are they angry at the Wellington Regional Medical Center. ``I don't think if I'd taken her
anywhere else, she would have survived,'' Mr.~Opdyke said.

Her flu came on gradually, and she never had a high fever. But after a week of feeling exhausted and
achy, she became delirious. When Mr.~Opdyke drove her to the hospital, with his finger on the
door-lock button for fear she would jump out, she could not tell the triage nurse her name.

Her blood oxygen level was below 70; normal is 95, and below 80 is life-threatening. Both lungs were
full of fluid.

Thrashing, she knocked off oxygen masks and pulled out an intubation tube. Panic made her
hyperventilate. Doctors finally sedated her into a paralytic coma to let the ventilator work.

She survived near-failure of her kidneys, then her lungs, damaged by continuous high-pressure
oxygen, began collapsing. Mr.~Opdyke was warned he might have to choose -- her life or that of the
baby, who was just at the border of survivability outside the womb.

``I said, 'Save Aubrey,' '' he said of the woman he married last year. ``I can make another baby,
but I can't replace her.''

Her third lung collapse forced the issue. Parker had to be delivered, but she did not survive.

Now I Lay Me Down to Sleep, a photographers' charity for families of premature infants, offered to
take black-and-white pictures. Ms.~Opdyke's mother bathed Parker and brushed her hair.

``Some people don't agree,'' Mr.~Opdyke said about having the photographs taken. ``But for me there
was no option. Aubrey wasn't awake, she didn't get to bond. If I didn't do it, there's no way I
could make it up to her.''

Aubrey Opdyke started to recover, then developed a bacterial infection common to ventilator
patients. She nearly succumbed, then rallied and doctors started to wean her off the coma-inducing
drugs. She had weeks of hallucinations about a friendly white dog visiting her but could not ask
about it because of the tube.

Mr.~Opdyke took comfort from what he saw as divine signs. Dog is ``God'' backwards, he noted. And
one day, in a panic because he had forgotten to wear his wife's wedding ring around his neck, he saw
a license plate reading ``FAITH.''

Ms.~Opdyke's last crisis was a still-unexplained seizure during a family visit. Her mother had been
painting her nails and blamed herself, thinking the fumes had done it.

Now that she is home, in the townhouse complex they rent near the West Palm Beach airport,
Ms.~Opdyke is struggling to regain her strength. Muscles atrophy into rubber bands; every day in a
coma means two days' recovery.

She apologizes for the messy yard; she would pick it up but barely has the strength to get upstairs
to bed. Mr.~Opdyke has gone back part-time to his job at UPS.

Mr.~Opdyke worries he gave her the flu. ``I was sick two weeks before,'' he said. ``I touch packages
that have been touched a thousand times. If anybody's going to catch it, it's me.''

He said he had not yet looked at the bills. His insurance covered most, he said, ``but it had to be
over a million, and we owe 20 percent; I'm not ready for that many commas.''

Friends from their jobs and a Girl Scout charity that Ms.~Opdyke once ran with her mother have
raised over \$10,000 by holding benefits, and small checks have come in after articles by local
journalists. ``It's a lot of help,'' Ms.~Opdyke said. ``It really shows the compassion of people.''

And does she want another baby?

``Yes,'' she says firmly. ``At first, I didn't. Now I do. But I've got to get my strength back.''


\vocabulary
\end{document}
